\chapter{The Benchmark Suite}\label{ch:bench-pipeline}

The package \module{benchmarks} measures ArchytasZX and a set of rival tools on a fixed
catalogue of quantum circuits, each paired with a specification it must equal. It generates a
false set of deliberately broken equations, schedules data points in the system size $n$ and the
dimension $d$, runs every point in a fresh worker process under CPU, wall, and memory limits,
logs every run as one JSON line, classifies every outcome by a closed list of codes, and writes
aggregate tables and plots. This chapter specifies the pipeline modules in
\file{benchmarks/} and the subject catalogue in \file{benchmarks/subjects/}. The rival-tool
adapters in \file{benchmarks/tools/} are the subject of \cref{ch:bench-tools}; this chapter
names the adapter interface only where the pipeline calls it.

\section{Architecture}\label{sec:bench-architecture}

\Cref{tab:bench-modules} lists the modules. Data flows in one direction: the configuration
selects subjects; the subjects yield equations and their scripts; the false set breaks each
equation once; the runner schedules points, launches workers, and appends records to the log;
the aggregator and the plotter read only the log, the false set, and the catalogue.

\begin{xltabular}{\textwidth}{@{}>{\ttfamily}l X@{}}
\caption{Modules of the benchmark suite.}\label{tab:bench-modules}\\
\toprule
\normalfont Module & Role\\
\midrule\endfirsthead
\toprule
\normalfont Module & Role\\
\midrule\endhead
\bottomrule\endfoot
\_\_main\_\_.py & Command-line entry point (\cref{sec:bench-cli}).\\
config.py, config.toml & The configuration, its overrides, and its hash (\cref{sec:bench-config}).\\
ir.py & The circuit intermediate representation (\cref{sec:bench-ir}).\\
translate.py & The naive translator from the IR to the engine language (\cref{sec:bench-translate}).\\
subject.py & Subjects, equations, versions, domains, and the engine script (\cref{sec:bench-subjects}).\\
sources.py & Source digests that key runs (\cref{sec:bench-sources}).\\
subjects/ & The subject catalogue and its shared builders (\cref{sec:bench-helpers,sec:bench-catalogue}).\\
breaks.py & The false set $F$ (\cref{sec:bench-false}).\\
differs.py & Whether an $F$ pair differs at a point (\cref{sec:bench-differs}).\\
points.py & The point schedule (\cref{sec:bench-points}).\\
stages.py & The ArchytasZX stages and their timing (\cref{sec:bench-stages}).\\
runner.py & Units, phases, resumption, and worker launch (\cref{sec:bench-runner}).\\
worker.py & One run in a fresh process (\cref{sec:bench-worker}).\\
log.py & The append-only results log (\cref{sec:bench-log}).\\
outcomes.py & Outcomes, stages, and cause codes (\cref{sec:bench-outcomes}).\\
stats.py & Repeat statistics (\cref{sec:bench-stats}).\\
aggregate.py & Aggregate tables (\cref{sec:bench-aggregate}).\\
plots.py & Plots (\cref{sec:bench-plots}).\\
reclassify.py & Reclassification of logged runs (\cref{sec:bench-reclassify}).\\
size.py & B3 element counts (\cref{sec:bench-size}).\\
lean.py & The Lean runner (\cref{sec:bench-lean}).\\
memcap.py & Per-child memory caps (\cref{sec:bench-memcap}).\\
memcap\_step.py & The in-step wrapper of the \code{srun} cap (\cref{sec:bench-memcap}).\\
anc.py & The ancillary-file export (\cref{sec:bench-anc}).\\
slurm.py & The cluster job array (\cref{sec:bench-slurm}).\\
\end{xltabular}

\paragraph{Terms.} A \term{subject} is one catalogue file. It declares one or more
\term{equations}; an equation pairs a circuit with a specification. An \term{instance} is one
equation at one \term{version}: \code{symbolic} (no concrete $n$ and $d$), \code{concrete} (a
point $(n,d)$), or \code{restricted} (a symbolic claim cut by restriction clauses). The set $S$
holds every instance of the selected subjects; the \term{false set} $F$ holds one broken copy of
each $S$ instance where a break applies. A \term{unit} is one (instance, tool) pair; a
\term{run} is one execution of a unit at one point and one repeat index.

\section{Command-line interface}\label{sec:bench-cli}

The suite runs as \code{python -m benchmarks ACTION [TARGET] [options]} from the repository
root. \Cref{tab:bench-actions} lists the actions; \cref{tab:bench-flags} lists the options.
Every action first resolves the configuration by \py!config.load(path, sets, profile)!
(\cref{sec:bench-config}).

\begin{xltabular}{\textwidth}{@{}>{\ttfamily}l X@{}}
\caption{Actions of \texttt{python -m benchmarks}.}\label{tab:bench-actions}\\
\toprule
\normalfont Action & Behavior\\
\midrule\endfirsthead
\bottomrule\endfoot
run & \py!runner.run(cfg, (i, n), only)! with the shard \code{i/n} of \code{--shard}
  (\cref{sec:bench-runner}).\\
false & \py!breaks.generate(cfg, RESULTS/false_set.json, platform.jobs)!; then prints one line
  per record: the equation id, the kind (or \code{None}), the label (or \code{None}), the word
  \code{skipped:}, and the comma-joined kinds tried before it.\\
aggregate & \py!aggregate.main(cfg)! (\cref{sec:bench-aggregate}).\\
plot & \py!plots.main(cfg)! (\cref{sec:bench-plots}).\\
anc & \py!anc.main(cfg, Path(TARGET))!; exits with the message \code{anc needs OUT_DIR}
  when \code{TARGET} is absent (\cref{sec:bench-anc}).\\
list & For every selected subject and every one of its $S$ instances, prints the instance id
  left-justified in 50 columns, a space, the category left-justified in 10 columns, and
  \code{sym=True} or \code{sym=False} (\py!subject.has_symbolic!).\\
slurm & \py!slurm.main(cfg, args)! (\cref{sec:bench-slurm}).\\
\end{xltabular}

\begin{xltabular}{\textwidth}{@{}>{\ttfamily}l X@{}}
\caption{Options of \texttt{python -m benchmarks}.}\label{tab:bench-flags}\\
\toprule
\normalfont Option & Meaning\\
\midrule\endfirsthead
\bottomrule\endfoot
TARGET & Optional positional argument; the output directory of \code{anc}, ignored elsewhere.\\
--config PATH & The TOML file to read instead of \file{benchmarks/config.toml}.\\
--profile laptop|cluster & Sets \code{platform.profile} after every \code{--set}.\\
--set KEY=VALUE & One override (\cref{alg:bench-override}); repeatable, applied in order.\\
--shard I/N & The shard of the units this process runs (default \code{0/1}); \code{run} only.\\
--only GLOB & A unit-id glob (\code{EQUATION|TOOL}); repeatable; \code{run} only.\\
\end{xltabular}

Two further entry points exist: \code{python -m benchmarks.worker JOB_JSON_FILE}
(\cref{sec:bench-worker}) and \code{python -m benchmarks.reclassify}
(\cref{sec:bench-reclassify}).

\section{Configuration}\label{sec:bench-config}

\subsection{The configuration file}

The file \file{benchmarks/config.toml} defines every value that determines what runs.
\Cref{tab:bench-config} lists every key with its default. Paths are relative to the repository
root unless absolute.

{\small
\begin{xltabular}{\textwidth}{@{}>{\ttfamily\raggedright\arraybackslash}p{0.32\textwidth} >{\ttfamily\raggedright\arraybackslash}p{0.27\textwidth} X@{}}
\caption{Keys of \texttt{benchmarks/config.toml}.}\label{tab:bench-config}\\
\toprule
\normalfont Key & \normalfont Default & Meaning\\
\midrule\endfirsthead
\toprule
\normalfont Key & \normalfont Default & Meaning\\
\midrule\endhead
\bottomrule\endfoot
subjects.catalogue & "benchmarks/subjects" & Directory of subject files.\\
subjects.include & ["*"] & Subject ids or globs that make up $S$.\\
subjects.exclude & [] & Subject ids or globs removed from $S$.\\
subjects.categories & ["algorithm", "code", "protocol", "circuit"] & Categories kept.\\
false\_set.enabled & true & Whether $F$ is generated and run.\\
false\_set.kinds & ["*"] & Globs over the break kinds allowed.\\
false\_set.assignment & "round\_robin" & \code{round_robin} or \code{random} (\cref{alg:bench-order}).\\
false\_set.seed & 20261106 & Seed of the random assignment and of counterexample candidates.\\
false\_set.beyond\_N & 32 & The $N$ of the \code{beyond_N} kind.\\
false\_set.tiny\_k & 3 & \code{tiny_phase} adds $1/d^k$ of a turn.\\
false\_set.probe\_seconds & 120 & Wall seconds of one counterexample probe.\\
false\_set.equation\_seconds & 1800 & Wall seconds of all kinds of one equation.\\
points.K & 16 & Successful data points per (equation, tool).\\
points.n\_min & 0 & Least $n$ drawn.\\
points.n\_max & 100\_000\_000 & Greatest $n$ drawn.\\
points.d\_min & 1 & Least $d$ drawn.\\
points.d\_max & 2048 & Greatest $d$ drawn.\\
points.distribution & "log\_uniform" & Read by no module; the draws are log-uniform.\\
points.seed & 20261106 & Seed of the phase 2 draws.\\
points.phase2\_extra\_failures & 32 & Cap on failed phase 2 draws per unit.\\
points.symbolic & true & Whether the symbolic version runs (outside $K$).\\
stages.decide & true & Counted only in the wall limit of a run (\cref{sec:bench-launch}); decide always runs.\\
stages.certificate & true & Run the proof-kernel certificate stage.\\
stages.lean\_export & true & Run the Lean export stage.\\
stages.lean\_check & true & Run the Lean check stage.\\
stages.timeout\_seconds & 3600 & CPU seconds per stage per run.\\
stages.wall\_factor & 4 & Wall limit multiplier (\cref{sec:bench-launch}).\\
stages.memory\_fraction & 0.85 & Share of available memory given to workers.\\
repeats.cluster.mode & "adaptive" & Repeat rule of the cluster profile.\\
repeats.cluster.min\_runs & 5 & Least runs per successful point.\\
repeats.cluster.max\_runs & 30 & Most runs per successful point.\\
repeats.cluster.confidence & 0.95 & Coverage of the median interval.\\
repeats.cluster.precision & 0.05 & Allowed relative half-width of the median interval.\\
repeats.laptop.mode & "fixed" & Repeat rule of the laptop profile.\\
repeats.laptop.runs & 3 & Runs per successful point.\\
trace.enabled & true & Runs carry engine timing spans.\\
trace.overhead\_sample & 0.05 & Share of points also run with tracing off.\\
tools.enabled & (12 names) & \code{archytaszx}, \code{pyzx}, \code{quizx}, \code{qcec},
  \code{feynman}, \code{dizx}, \code{qiskit}, \code{cirq}, \code{autoq_para}, \code{tket},
  \code{quantomatic}, \code{published_proofs}.\\
platform.profile & "laptop" & \code{laptop} or \code{cluster}.\\
platform.jobs & "auto" & Concurrent workers; \code{auto} resolves at load.\\
output.results & "benchmarks/results" & Results directory.\\
output.plots & "benchmarks/plots" & Plot directory.\\
slurm.partition & "scavenger" & Slurm partition; any other value is refused.\\
slurm.array & 32 & Array tasks; each runs one shard.\\
slurm.cpus\_per\_task & 8 & Cores per task and workers per task.\\
slurm.mem\_per\_cpu & "4G" & Slurm memory per core.\\
slurm.time & "7-00:00:00" & Slurm time limit.\\
slurm.account & "" & Slurm account; empty: no account line.\\
slurm.venv & "\textasciitilde/.venvs/personal\_lab" & Python environment on the cluster.\\
slurm.setup & "" & A shell file sourced before the environment; empty: none.\\
slurm.commit & "" & The commit run; empty: the latest green \code{dev} commit.\\
\end{xltabular}
}

The worker also reads \code{stages.prepare_timeout_seconds}, which the file does not define; its
default is \code{stages.timeout_seconds} (\cref{sec:bench-worker}).

\subsection{Loading and overrides}

\begin{signature}
\py!load(path: Path | None = None, sets: list[str] | None = None, profile: str | None = None) -> dict!
\end{signature}
\noindent reads \code{path} (default \file{benchmarks/config.toml}) with a TOML parser, applies
every assignment of \code{sets} in order by \py!override!, then sets \code{platform.profile} to
\code{profile} when given, and finally replaces \code{platform.jobs = "auto"} by
$\max(1, \lfloor c/2\rfloor)$, where $c$ is the CPU count (2 when unknown).

\begin{algorithm}[\texttt{override}]\label{alg:bench-override}
Input: the configuration dictionary and one text \code{KEY=VALUE}.
\begin{enumerate}
\item Split at the first \code{=}; with no \code{=}, raise \py!ValueError("--set needs key=value, got ...")!.
\item Strip \code{KEY} and split it at dots into $p_1,\dots,p_m$. Walk $p_1,\dots,p_{m-1}$; when
  a part is missing or does not name a table, raise \py!KeyError("unknown config section ...")!.
\item When $p_m$ is not a key of the last table, raise \py!KeyError("unknown config key ...")!.
  New keys are never created.
\item The value is the TOML literal \code{v = VALUE} parsed by the TOML parser (integers,
  floats, booleans, strings in quotes, arrays); when that parse fails, the stripped text itself
  as a string. Thus \code{--set points.K=4} stores the integer 4 and \code{--set platform.profile=cluster} stores the string \code{"cluster"}.
\end{enumerate}
\end{algorithm}

\py!repeats(cfg)! returns \code{cfg["repeats"][cfg["platform"]["profile"]]}.
\py!path(cfg, section)! returns \code{cfg["output"][section]} as an absolute path, resolved
against the repository root when relative.

\subsection{The configuration hash}

\begin{definition}[Configuration hash]\label{def:bench-digest}
\py!digest(cfg)! deep-copies the configuration, removes the five keys \code{platform.jobs},
\code{output.results}, \code{output.plots}, \code{tools.enabled}, and \code{slurm.commit}
(the constant \py!UNHASHED!), serializes the rest as JSON with sorted keys, and returns the
first 12 hexadecimal digits of its SHA-256 hash.
\end{definition}

The hash keys every run (\cref{sec:bench-log}); changing any hashed value starts a fresh set
of runs, while changing the worker count, the output paths, the enabled tools, or the commit
reuses logged runs.

\section{The circuit IR}\label{sec:bench-ir}

\module{benchmarks.ir} defines the circuits of the catalogue as qudit gate lists over named
wires and registers. All gate classes are frozen dataclasses (value equality, hashable when
their fields are).

\subsection{Wires, registers, and amounts}

A \term{wire} is a name. A \term{register} is a wire name bound to a \term{count}: an
integer or a symbol expression such as \code{"n"} or \code{"n + 1"}. A gate that names a
register acts on every copy of it, the same gate on each. Registers that share a count lie
together in the boundary order, copy-major. A \term{two-count register} is bound to a pair
\code{(outer, inner)}; its copies $(m,j)$ lie $m$-major, and it occurs only in a
\py!Chain! of the outer count, whose copy $m$ acts on row $m$. A wire with no \py!Init! is a
boundary input; a wire with no \py!Post! or \py!Discard! is a boundary output.

\term{Amounts} (the \code{a} of \py!X!, \py!Z!) and phase rules are integers, fractions, or
SymPy expressions in \code{d}, \code{n}, and the phase index \code{k}; phases are in
\term{turns} (a phase $t$ means $e^{2\pi i t}$). Write $\omega = e^{2\pi i/D}$ for a wire of
dimension $D$.

\subsection{Gates}

\Cref{tab:bench-gates} defines every gate by its fields and its action on the computational
basis $\ket{k}$, $k\in\Zd{D}$.

\begin{xltabular}{\textwidth}{@{}>{\ttfamily}p{0.27\textwidth} X@{}}
\caption{Gates of the circuit IR.}\label{tab:bench-gates}\\
\toprule
\normalfont Gate (fields, defaults) & Action\\
\midrule\endfirsthead
\toprule
\normalfont Gate (fields, defaults) & Action\\
\midrule\endhead
\bottomrule\endfoot
Init(q, v=0) & The basis state $\ket{v}$ on \code{q}.\\
Post(q, v=0) & The basis effect $\bra{v}$; \code{q} ends.\\
Discard(q) & Trace out \code{q}; \code{q} ends.\\
Measure(q) & Computational-basis measurement: \code{q} is dephased and stays open, classical.\\
X(q, a=1) & Shift $\ket{k}\mapsto\ket{k+a}$.\\
Z(q, a=1) & Clock $\ket{k}\mapsto\omega^{ak}\ket{k}$.\\
F(q, p=1) & $F^p$ with $F_{jk} = \omega^{jk}/\sqrt{D}$.\\
ZPhase(q, entries=\{\}, rule=None) & $\ket{k}\mapsto e^{2\pi i\,\varphi_k}\ket{k}$, $\varphi$
  from \code{entries} $\{k:\text{turns}\}$ or from \code{rule}, an expression in \code{k}.\\
XPhase(q, entries, rule) & $F\cdot\mathrm{ZPhase}\cdot F^\dagger$.\\
CX(c, t, s=1) & $\ket{c,t}\mapsto\ket{c, t+sc}$, $s=\pm1$.\\
CZ(c, t, s=1) & $\ket{c,t}\mapsto\omega^{sct}\ket{c,t}$, $s=\pm1$.\\
Swap(a, b) & Exchange two wires.\\
PhaseFn(qs, f, support=None) & $\ket{x}\mapsto e^{2\pi i f(x)}\ket{x}$ on several wires,
  $f$ in turns; \code{support} lists the only $x$ where $f$ may be nonzero (\code{None}: any).\\
MCX(cs, t, values=None, a=1) & $\ket{c,t}\mapsto\ket{c, t + a[c = \text{values}]}$;
  \code{values=None} means all 1.\\
MCZ(cs, t, values=None, a=1) & $\ket{c,t}\mapsto\omega^{a t [c=\text{values}]}\ket{c,t}$.\\
Perm(qs, p) & $\ket{x}\mapsto\ket{p(x)}$ on the joint basis of \code{qs}, first wire most
  significant; $p:\N\to\N$.\\
Oracle(xs, y, f) & $\ket{x,y}\mapsto\ket{x, y+f(x)}$, $y$ one wire.\\
Unitary(qs, matrix) & A dense unitary given as nested lists of \code{[real, imag]} pairs;
  reference circuits only.\\
Prepare(qs, vector) & A dense state on initially $\ket{0}$ wires; reference circuits only.\\
Chain(count, body, carries=(), index=None, reverse=False, through=()) &
  A loop over \code{count} copies (\cref{sec:bench-ir-loops}).\\
Tree(count, body, parent, children, level=None, position=None) & A tree of depth
  \code{count} (\cref{sec:bench-ir-loops}).\\
Pairs(over, body, index="j", pair\_index="m") & For every pair of copies $j<k$ of register
  \code{over}, \code{body} on wires \code{"left"} (copy $j$) and \code{"right"} (copy $k$).\\
LowerMCX(over, controls=(), values=(), reverse=False) & For each copy $k$ of register
  \code{over} (dimension 2), from the last copy down (up with \code{reverse}), the flip of copy
  $k$ controlled on every copy $j<k$ at 1 and on the single wires \code{controls} at
  \code{values}.\\
IndexedMCX(over, cells, t) & For each value $v=0,\dots,2^n-1$ of register \code{over} ($n$
  copies, dimension 2, copy $j$ the bit of value $2^j$), in order, the flip of single wire
  \code{t} controlled on \code{over}${}=v$ and on copy $v$ of register \code{cells}
  ($2^n$ copies) at 1.\\
Neighbours(over, s=1) & $\mathrm{CZ}^s$ between copies $j$ and $j+1$ of register
  \code{over}, for every $j$, in copy order.\\
\end{xltabular}

\subsection{Loops}\label{sec:bench-ir-loops}

\paragraph{Chain.} Copy $m$ of a \py!Chain! reads its \code{index} symbol as $m$ in amounts and
phase rules. Its \code{body} acts on copy $m$ of each register of the chain's count and on the
single wires named in \code{carries}. After each copy, every pair \code{(carry, register)}
exchanges values: the register's copy keeps the carry's value and the carry continues with the
register's. Each single wire of \code{through} passes from copy to copy unchanged. With
\code{reverse} the copies run from $\mathit{count}-1$ down to 0. A carry may itself be a
register whose count is the inner count of its two-count partner; it then carries one row from
copy to copy (copy $j$ to copy $j$).

\paragraph{Tree.} Copy $c(L,P)$ of a \py!Tree! ($L = 1,\dots,\mathit{count}$,
$P < b^{L-1}$, $b$ the number of \code{children}) runs \code{body} on its parent value as wire
\code{parent} and on fresh local wires; the final values of the \code{children} wires feed
copies $c(L+1, bP+i)$, $i<b$; the leaves' children take \code{parent}'s place in the boundary,
leaf-major. Symbols \code{level} and \code{position} read $L$ and $P$.

\subsection{Circuits}

\begin{signature}
\py!Circuit(wires, gates, registers={}, dims={}, scalar="1", outputs=None)!
\end{signature}
\noindent holds the wires in input boundary order, the gates, the registers with their counts,
per-wire dimensions (an integer or an expression; \py!dim(q)! returns \code{"d"} when absent),
a global scalar expression, and the output boundary order. \py!count(q)! returns the register
count (the outer count of a two-count register) or \code{None}; \py!inner(q)! the inner count or
\code{None}. \py!output_order()! returns \code{wires} when \code{outputs} is \code{None};
otherwise \code{outputs} restricted to known wires, each omitted wire inserted right after the
wire that precedes it in \code{wires} (first when it has no predecessor).

The constructor function
\py!circuit(wires, gates, *, registers=None, dims=None, scalar="1", outputs=None)! raises
\py!ValueError("<Gate> names undeclared wire(s) [...]")! when a gate names an undeclared wire
(by \py!gate_wires!), and \py!ValueError("outputs ... is not an order of the wires ...")!
when \code{outputs} is not a permutation of \code{wires}.

\py!gate_wires(g)! returns the wires a gate names, in argument order: $(c,t)$ for \py!CX!,
\py!CZ!; $(a,b)$ for \py!Swap!; \code{qs} for \py!PhaseFn!, \py!Perm!, \py!Unitary!,
\py!Prepare!; $(\mathit{cs},t)$ for \py!MCX!, \py!MCZ!; $(\mathit{xs},y)$ for \py!Oracle!; for
\py!Chain! the body's wires, then every carry pair's wires, then \code{through}, deduplicated in
first-seen order; \code{(over,)} for \py!Pairs! and \py!Neighbours!; \code{(parent,)} for
\py!Tree!; $(\mathit{controls}, \mathit{over})$ for \py!LowerMCX!; $(\mathit{over},
\mathit{cells}, t)$ for \py!IndexedMCX!; \code{(q,)} otherwise.

\py!at_index(g, values)! returns \code{g} with every symbol of \code{values} replaced by its
integer in the amounts of \py!X!, \py!Z!, the entries and rule of \py!ZPhase!, \py!XPhase!,
and recursively in a \py!Chain! body; an amount is replaced only when it mentions one of the
symbols. \py!unrolled(name, n)! returns the wire names \code{name0}, \dots,
\texttt{name\{n-1\}}.

\subsection{Abstract grammar}

\begin{lstlisting}[style=grammar,caption={Abstract grammar of the circuit IR.}]
circuit  ::= Circuit(wires: name*, gates: gate*, registers: {name: count | (count, count)},
                     dims: {name: expr}, scalar: expr, outputs: name* | None)
count    ::= int | expr                      expr: SymPy text in d, n (and index symbols)
amount   ::= int | expr                      phase entries/rules in turns, k the basis index
gate     ::= Init(name, int) | Post(name, int) | Discard(name) | Measure(name)
           | X(name, amount) | Z(name, amount) | F(name, int)
           | ZPhase(name, {int: turns}, rule) | XPhase(name, {int: turns}, rule)
           | CX(name, name, +1|-1) | CZ(name, name, +1|-1) | Swap(name, name)
           | PhaseFn(name+, fn, support) | MCX(name+, name, values, int)
           | MCZ(name+, name, values, int) | Perm(name+, fn) | Oracle(name+, name, fn)
           | Unitary(name+, matrix) | Prepare(name+, vector)
           | Chain(count, gate+, (name, name)*, index, bool, name*)
           | Tree(count, gate+, name, name+, level, position)
           | Pairs(name, gate+, index, index) | LowerMCX(name, name*, int*, bool)
           | IndexedMCX(name, name, name) | Neighbours(name, +1|-1)
\end{lstlisting}

\section{The naive translator}\label{sec:bench-translate}

\module{benchmarks.translate} maps a circuit to an engine-language diagram body (the syntax of
\cref{ch:language}): one generator group per gate, every node named, every wire explicit, and
no simplification. The result is the diagram a user would write by hand from the circuit; the
benchmark measures how the engine reduces it.

\subsection{Contexts and symbol writing}

\begin{signature}
\py!Context(values: Mapping[str, int] = {}, literal: frozenset[str] = frozenset())!
\end{signature}
\noindent binds symbols (\code{d}, \code{n}) to integers; a bound symbol in \code{literal} is
written as its value. The methods are memoized in one LRU cache of 65\,536 entries keyed on
(function, sorted literal bindings, type of the expression, expression); unhashable
expressions bypass the cache.
\begin{itemize}
\item \py!dim_token(dim)! returns \code{literal V} when the dimension is an integer, or becomes
  one after substituting the literal symbols; otherwise the expression as text.
\item \py!value(expr)! substitutes the literal symbols; a numeric result is passed through
  \py!sympy.nsimplify!.
\item \py!concrete(expr)! returns the integer after substitution, or \code{None}.
\end{itemize}

Expressions are written by \py!str! of the SymPy expression with \code{**} replaced by
\texttt{\^{}}. A scalar whose rational coefficient $p/q$ has $\max(|p|, q)$ above 160 bits is written
as prime powers when both $|p|$ and $q$ factor over the primes below 1000:
\texttt{-(2\^{}3*5)/(3\^{}2)*(rest)}, the minus sign only for $p<0$, the denominator only when not 1,
the factor \code{*(rest)} only when the non-rational part is not 1; otherwise it is written as
an ordinary expression. A concrete number of turns is reduced into $[0,1)$; a float is first
converted to the nearest rational by \py!nsimplify(..., rational=True)!.

\subsection{Gate encodings}

\Cref{tab:bench-encodings} gives the generator group of every gate; $D$ is the wire's
dimension token, and the scalar column is the factor multiplied into the diagram's scalar.

\begin{xltabular}{\textwidth}{@{}>{\ttfamily}p{0.15\textwidth} X >{\raggedright\arraybackslash}p{0.15\textwidth}@{}}
\caption{Naive encodings of the IR gates.}\label{tab:bench-encodings}\\
\toprule
\normalfont Gate & Nodes & Scalar\\
\midrule\endfirsthead
\toprule
\normalfont Gate & Nodes & Scalar\\
\midrule\endhead
\bottomrule\endfoot
Init v & \azx!X(0 -> 1, D) phase{k: -v*k/D}! (the phase omitted when it simplifies to 0) & $D^{-1/2}$\\
Post v & \azx!X(1 -> 0, D) phase{k: v*k/D}! & $D^{-1/2}$\\
Discard & \azx!G(1 -> 0, D)! & 1\\
Measure & \azx!Z(1 -> 2, D)!, \code{out1} wired to a \azx!G(1 -> 0, D)!; the wire continues from \code{out0} & 1\\
X a & \azx!X(1 -> 1, D) phase{k: -a*k/D}! & 1\\
Z a & \azx!Z(1 -> 1, D) phase{k: a*k/D}! & 1\\
F p & $p \bmod 4$ nodes \azx!F(D)! in sequence & 1\\
ZPhase & \azx!Z(1 -> 1, D)! with the rule (\texttt{phase\{k: ...\}}) or the entries & 1\\
XPhase & \azx!X(1 -> 1, D)! likewise & 1\\
Swap & No node; the two current ends exchange & 1\\
CX s & \azx!Z(1 -> 2, D)! on $c$, \azx!X(2 -> 1, D)! on $t$, $c$'s \code{out1} to $t$'s
  \code{in1} directly ($s=1$) or through two \azx!F(D)! ($s=-1$) & $D^{1/2}$\\
CZ s & \azx!Z(1 -> 2, D)! on $c$, \azx!Z(2 -> 1, D)! on $t$, joined through one \azx!F(D)!
  ($s=1$) or three ($s=-1$) & $D^{1/2}$\\
PhaseFn & Binders to one port of the product dimension, a $Z$ spider with the phase table,
  splitters back (\cref{alg:bench-phasefn}) & see text\\
MCZ, MCX & \py!PhaseFn! on (controls, target) with
  $f(x) = a\,x_t/D_t$ when the controls equal \code{values}, else 0, with the support
  $\{(\mathit{values}, t) : t < D_t\}$; \py!MCX! adds $F$ before and $F^3$ after on the
  target; the target dimension must be concrete & --\\
Oracle & $F$ on $y$, \py!PhaseFn! $h(x,y) = f(x)\,y/D_y$, $F^3$ on $y$ & --\\
Perm & The fresh-port encoding (\cref{alg:bench-perm}) & $T^{-1/2}$\\
Unitary, Prepare & Refused: \code{a dense reference gate has no ZXW translation} & --\\
\end{xltabular}

A wire never touched by any gate receives one identity \azx!Z(1 -> 1, D)!, which is both its
input and its output port. \py!CX! and \py!CZ! raise \py!ValueError! when $s\notin\{1,-1\}$ and
refuse wires of different dimensions. \py!Init! on a wire that has already started raises
\py!ValueError("Init on wire q after it started")!; using a wire after \py!Post! or \py!Discard!
raises \py!ValueError("wire q is used after it ended")!.

\begin{algorithm}[Phase table of \texttt{PhaseFn}]\label{alg:bench-phasefn}
Input: wires $q_1,\dots,q_m$, a function $f$ in turns, an optional support.
\begin{enumerate}
\item Let $B$ be the box of the wires (\cref{sec:bench-boxes}). When $B$ exists, every
  wire $q_i$ that is not boxed receives a fan \azx!Z(1 -> 2, D)! outside the box: the wire
  continues from \code{out0}, and \code{out1} feeds the table; at the end the table's output
  for $q_i$ goes into a \azx!Z(1 -> 0, D)! inside the box.
\item Every dimension must be concrete after substitution; else refuse with \code{a phase table needs literal concrete dimensions}. Let $d_1,\dots,d_m$ be the dimensions and $T=\prod d_i$.
\item Enumerate the joint indices $x$ in row-major order, first wire most significant (with a
  support, only its points inside the ranges, sorted by joint index). Each value $f(x)$ must be
  a rational number of turns (else refuse with \code{a phase table entry is not a rational number of turns}); keep the nonzero values reduced into $[0,1)$.
\item When the entry $t_0$ at index 0 is nonzero, replace every entry $i\in[1,T)$ by
  $(t_i - t_0) \bmod 1$, drop zeros, and multiply $e^{2\pi i t_0}$ into the scalar (per copy of
  $B$).
\item With $m=1$: one node \azx!Z(1 -> 1, literal T) phase{i: t, ...}!.
\item Otherwise a chain of binders: \azx!B(literal d1, literal d2)! takes $q_1$ on \code{in0} and
  $q_2$ on \code{in1}; each further binder \azx!B(literal A, literal di)! takes the previous
  \code{out0} on \code{in0} and $q_i$ on \code{in1}, $A$ the product so far. Then the spider
  \azx!Z(1 -> 1, literal T)! with the table, then splitters in reverse order
  \azx!S(literal A, literal di)!, each giving $q_i$ its \code{out1} and passing \code{out0} on;
  the last \code{out0} becomes $q_1$.
\end{enumerate}
\end{algorithm}

\begin{algorithm}[Fresh-port encoding of \texttt{Perm}]\label{alg:bench-perm}
Refused on a register (\code{a permutation on a register}) and for symbolic dimensions.
\begin{enumerate}
\item With several wires, bind them by binders as in \cref{alg:bench-phasefn} into one port $x$
  of dimension $T$.
\item Create $t$ = \azx!Z(0 -> 1, literal T)! and a binder \azx!B(literal T, literal T)! with $t$
  on \code{in0} and $x$ on \code{in1}.
\item A spider \azx!Z(1 -> 1, literal T^2)! whose entry at index $t_vT + x_v$ is
  $-t_v\,p(x_v)/T \bmod 1$ for every nonzero value.
\item A splitter \azx!S(literal T, literal T)!; \code{out1} into \azx!Z(1 -> 0, literal T)!,
  \code{out0} into \azx!F(literal T)!; the scalar receives $T^{-1/2}$.
\item Split $F$'s output back into the wires with \azx!S(literal A, literal di)!, $A$ the
  product of the dimensions before $q_i$.
\end{enumerate}
\end{algorithm}

\subsection{Boxes}\label{sec:bench-boxes}

A register's nodes go into one node-scoped \bbox{} per count, named by the count text. A gate
between a register and a single wire places the single wire's spider outside the \bbox{}, where
it receives one port per copy. The builder keeps, for each \bbox{}: its count, its node list,
its name (\code{C0}, \code{C1}, \dots{} in creation order of named \bboxes), its index symbol,
extra lines, a per-copy scalar list (tree \bboxes{} only), its parent \bbox, and its chain links.

\py!box_of(qs)! decides the \bbox{} of a gate:
\begin{enumerate}
\item Inside a chain, the registers of the chain's inner count (the \code{wide} ones) are left
  out; the other registers must share one count (else refuse with \code{gate across registers of different counts}) equal to the chain's (else \code{a gate in a chain on a register of another count}). A gate on a wide register goes into the inner \bbox; a gate on a boxed wire
  into the chain's \bbox; otherwise none.
\item Outside chains, a register already consumed by a chain and not recorded as after it is
  refused (\code{a chain register used outside its chain}). Gates on registers listed as
  before or after a chain go into that chain's \bbox{} (refused with \code{a gate before|after a chain on registers of another box} when the wires disagree).
\item Otherwise the common register count, or \code{None}.
\end{enumerate}

A factor $v$ placed in a \bbox{} with no per-copy list enters the global scalar as
$v^{\mu}$, $\mu$ the product of the \bbox's multiplicity and its ancestors'. Before translating,
the builder scans every \py!Chain! at gate position $i$: each register of the chain (not a row
carry) already used by an earlier gate is recorded as \emph{before} the chain, with a
provisional \bbox{} \code{_pre<i>} of the chain's count whose nodes merge into the chain's
\bbox{} when the chain is translated.

\subsection{Loop encodings}

\paragraph{Chain.} A chain becomes one named \bbox{} over its body's nodes, with its index,
plus the line \code{chain NAME: p -> q} per carry, linking the carry's end after the exchange
($p$) to the first port the carry entered in the copy ($q$); a reversed chain writes
\code{q -> p}. A carry or register with no node in the copy receives an identity
\azx!Z(1 -> 1, D)!. A carry that ends in the copy takes the register's end, and the register's
copy then has no output. Gates on a two-count register, or on a row carry, inside the chain go
into a \bbox{} of the inner count nested in the chain's (named \code{<chain>i}); a row carry's
link runs through that \bbox, and a row carry receives an identity in its own \bbox{} before and
after the chain. A chain inside a chain, a chain over a register of another count, a register
used before its chain (outside the before-scan), a carry that is a register (other than a row
carry), and a through wire or register that ends inside the chain are refused.

\paragraph{Tree.} A tree becomes a named \bbox{} with a per-copy scalar list and the line
\code{tree NAME: HEAD <- c1, ..., cb [level L] [position P]}, where \code{HEAD} is the port the
parent entered first and $c_i$ the final ports of the children. A tree inside a loop or on a
register, a body naming another circuit wire, a child or parent that ends in the tree, and a
local wire that neither ends nor is a child are refused. The per-copy factors appear as
\code{copy scalar NAME: EXPR}, the product simplified by \py!powsimp(force=True)!.

\paragraph{Neighbors.} The register's \bbox{} receives a \azx!Z(2 -> 2, D)! on the wire and one
$F$ ($s=1$) or three ($s=-1$) feeding its \code{in1}; the line \code{chain NAME: z.out1 -> f.in0} links copy $j$'s \code{out1} to copy $j+1$'s first $F$. The head \azx!X(0 -> 1, D)! and
tail \azx!Z(1 -> 0, D)! lie in the parent \bbox{}; the scalar receives $D^{1/2}$ per copy and
$D^{-1/2}$ once.

\paragraph{Pairs.} Allowed only outside chains with exactly one \py!CZ! body between
\code{left} and \code{right}. The register's \bbox{} (which takes \code{index}) receives
\azx!Z(2 -> 1, D)! (its \code{in1} the right port) and then \azx!Z(1 -> 2, D)! (its \code{out1}
the left port); one or three $F$ nodes, outside any \bbox{}, join them. The lines are
\code{pairs P<i> of NAME index PAIR_INDEX: f...}, \code{left P<i>: PORT}, \code{right P<i>: PORT}, and \texttt{copy scalar P<i>: D\^{}(1/2)}.

\paragraph{LowerMCX ($D=2$).} Per copy $k$ in the register's \bbox{} (index \code{j}): a fan
\azx!Z(1 -> 2, literal 2)! (before the target; after it with \code{reverse}), the target
\azx!X(2 -> 1, literal 2)! fed by \azx!Ti! after a $Z$ AND spider \azx!Z(1+c -> 1, literal 2)!
($c$ controls), scalar $\sqrt2$ per copy; one \azx!T! per single control into the AND, through
\azx!X(1 -> 1) phase{1: 1/2}! for value 0, the control's fan outside the \bbox. A pair \bbox{}
(index \code{m}) holds one \azx!T! from copy $j$'s fan port to copy $k$'s AND, so that the AND
reads $T\ket{0}+\prod x\,\ket{1} = T\ket{\prod x}$.

\paragraph{IndexedMCX ($D=2$).} Per value copy $v$ in the cells' \bbox{} (index \code{v}): a fan
on the cell, $\ket{v}$ as \azx!X(0 -> 1, literal 2^n) phase{k: -v*k/2^n}! split by a chained
\azx!S(literal 2, literal 2^j)! in a \bbox{} of $n$ copies (index \code{j}) nested in it and
aligned to the address \bbox{} (\code{align A with B}); copy $j$ feeds \azx!X(2 -> 1)! of bit
$j$ with the address fan port, through \azx!X(pi)! and \azx!T! into a $Z$ AND spider with the
cell's \azx!T!, then \azx!Ti! into one target \azx!X(2 -> 1)! on $t$. The scalar is
$2^{-n/2}\cdot2^{1/2}$ per value copy and $2^{1/2}$ per bit copy. It requires two registers with
$2^n$ cells, a single target, dimension 2, and no enclosing loop.

\subsection{Output}

\begin{signature}
\py!translate(c: Circuit, ctx: Context) -> Translation(body: str, nodes: int, counts: tuple[str, ...])!
\end{signature}

\begin{algorithm}[\texttt{translate}]\label{alg:bench-translate}
\begin{enumerate}
\item Check the layout: in the input order and in the output order, the registers of one count
  are contiguous; else raise \py!ValueError("registers of count K are not contiguous")!.
\item Translate every gate in order. Nodes are named \code{g0}, \code{g1}, \dots{} in creation
  order; a node created in a \bbox{} also joins every ancestor \bbox.
\item Give every untouched wire its identity node.
\item Emit, in order: every declaration \code{gK = GEN}; one line \code{wire a -- b, ...} with
  every wire; per \bbox{} \code{repeat TOKEN[ as NAME][ index I][ in PARENT]: nodes}, then its
  \code{chain} line, its extra lines, and its \code{copy scalar}; the pair lines; \code{inputs ...} in wire order; \code{outputs ...} in \py!output_order!; \code{scalar EXPR} when the
  product of every factor and the circuit's scalar (simplified by \py!powsimp(force=True)!)
  is not 1. The \bbox{} token is \code{literal C} for a digit count, \code{literal V} when every
  free symbol of the count is literal, else the count expression.
\item \code{nodes} is the number of declarations; \code{counts} the sorted non-digit register
  counts.
\end{enumerate}
\end{algorithm}

\begin{example}
The Bell circuit \py!circuit(["a","b"], [Init("a"), F("a"), Init("b"), CX("a","b")])! with
symbolic $d$ translates to
\begin{lstlisting}[style=azx]
g0 = X(0 -> 1, d)
g1 = F(d)
g2 = X(0 -> 1, d)
g3 = Z(1 -> 2, d)
g4 = X(2 -> 1, d)
wire g0.out0 -- g1.in0, g1.out0 -- g3.in0, g2.out0 -- g4.in0, g3.out1 -- g4.in1
outputs g3.out0, g4.out0
scalar 1/sqrt(d)
\end{lstlisting}
The factors $d^{-1/2}\cdot d^{-1/2}\cdot d^{1/2}$ combine into $d^{-1/2}$.
\end{example}

\py!Unsupported! (a subclass of \py!Exception!) signals that a circuit cannot be translated at
this version; its message is one of the refusal texts above, and \cref{tab:bench-unsupported}
maps each to a code.

\section{Subjects, equations, and versions}\label{sec:bench-subjects}

\module{benchmarks.subject} defines the declarations of a subject file and builds the engine
script of one instance at one point.

\subsection{Domains}

\begin{signature}
\py!Domain(n_min=0, n_max=None, n_step=1, d_min=1, d_max=None, d_condition="any")!
\end{signature}
\noindent is the set of points where a claim holds: $n\in[n_{\min}, n_{\max}]$ with
$(n - n_{\min}) \bmod n_{\mathrm{step}} = 0$, and $d\in[d_{\min}, d_{\max}]$ satisfying the
condition. \code{None} bounds are open. \py!has_n(n)!, \py!has_d(d)!, and \code{(n, d) in dom}
test membership. The conditions are those of \py!D_CONDITIONS!: \code{any} ($d\ge1$),
\code{ge2} ($d\ge2$), \code{prime}, \code{odd}, \code{odd_prime}, \code{even}; \py!is_prime(v)!
is trial division (false below 2).

\paragraph{Restriction clauses.} A restricted claim adds clauses in the engine's \code{for}
syntax. Spaces are removed; the first character is the symbol (\code{n} or \code{d}).

\begin{lstlisting}[style=grammar,caption={Restriction clauses.}]
clause ::= SYM "=" INT              n = v  or  d = v      (pins the value)
         | SYM ">=" INT             n >= c or  d >= c     (raises the lower bound)
         | "n" "=" [A ["*"]] "m" ["+" C]                  (n = A*m + C, A defaults to 1)
\end{lstlisting}

\py!restrict_domain(dom, clause)! returns the cut domain: \code{d = v} sets
$d_{\min}=d_{\max}=v$; \code{d >= c} raises $d_{\min}$ to $\max(d_{\min}, c)$; \code{n >= c}
raises $n_{\min}$; \code{n = A*m + C} sets $n_{\mathrm{step}}=A$ and $n_{\min}$ to the least
$C + jA \ge n_{\min}$; \code{n = v} pins $n$. \py!clause_holds(clause, value)! tests one value:
$\ge$ by comparison, the progression by $v\ge C$ and $(v-C)\bmod A = 0$, else equality.

\subsection{Specification contexts}

\begin{signature}
\py!SpecContext(D: str, d: str, n: str, N: str, nv: int | None, dv: int | None)!
\end{signature}
\noindent tells a specification writer how to write symbols: \code{D} the node dimension token
(\code{d} or \code{literal 5}), \code{d} the dimension in phases and scalars, \code{n} the count
in expressions, \code{N} the \bbox{} multiplicity token, and the integers \code{nv}, \code{dv}
(\code{None} when symbolic). \py!spec_context(subject, nv, dv, literal_d)! takes $d$ from a fixed
subject dimension (forcing \code{literal_d}); with \code{literal_d} and a concrete $d$ it gives
\code{D = "literal dv"}, \code{d = "dv"}, otherwise \code{D = d = "d"}; always
\code{n = N = "n"}.

\subsection{Equations and subjects}

\begin{signature}
\py!Equation(key, circuit, specification, form="boxed", reference=None, literal=False, note="")!\\
\py!Subject(id, name, category, source, size_parameter, dimension, domain, equations,!\\
\py!        restricted=(), clifford=False, note="")!
\end{signature}

An equation holds \code{circuit}, a function $(n, d)\mapsto$ \py!Circuit!; \code{specification},
a function \py!SpecContext! $\mapsto$ diagram body text; an optional \code{reference}, a
function $(n,d)\mapsto$ \py!Circuit! realizing the specification for tools that read only
circuits; and \code{literal}, which writes every dimension as a literal. Its \term{form} says how
$n$ enters the circuit:
\begin{itemize}
\item \code{boxed}: $n$ is a register count; the circuit function receives the symbol
  \code{"n"}, so a symbolic version exists.
\item \code{unrolled}: the circuit function receives an integer and creates one wire per qudit.
\item \code{fixed}: the circuit has no $n$.
\end{itemize}

A subject's \code{category} is one of \code{algorithm}, \code{code}, \code{protocol},
\code{circuit} (else \py!ValueError!); \code{source} is a bibliography key of the algorithm's
paper; \code{size_parameter} names what $n$ counts, or \code{None}; \code{dimension} is
\code{"general"} or a fixed integer; \code{restricted} lists clause tuples, each a restricted
claim; \code{clifford} marks a Clifford subject. When \code{size_parameter} is \code{None} and
the domain's $n_{\min}\neq n_{\max}$, the constructor replaces the domain's $n$ range by
$[0,0]$. \py!general_d! is true when \code{dimension == "general"}.

\subsection{Instances}

\begin{signature}
\py!Instance(subject, equation, restriction=(), falsified=None)!
\end{signature}
\noindent is one equation at one version. Its \py!id! is \code{SUBJECT/EQUATION}, with
\code{/r<k>} appended for the $k$-th restriction and \code{!} appended for an $F$ instance
(\code{falsified} an \py!Edit!, \cref{sec:bench-false}). \py!domain! is the subject's domain cut
by each restriction clause, with $d_{\min}$ raised to the edit's parameter \code{d_min} when the
edit has one. \py!truth! is true for an $S$ instance.

\py!instances(subjects)! yields, per subject and equation, the unrestricted instance and then one
per restriction. \py!has_symbolic(inst)! is false for an unrolled equation and otherwise true when
the subject has a general $d$ or a boxed equation with a size parameter. \py!axes(inst)! returns
(has size parameter, general $d$). \py!select(subjects, include, exclude, categories)! keeps the
subjects in the categories whose id matches an include glob and no exclude glob
(\py!fnmatch!). \py!catalogue(directory)! (cached) imports every \file{*.py} of the directory not
starting with \code!_!, in sorted order, as module \code{benchmarks_subject_<stem>}, and
returns their \code{SUBJECT} objects; a subject whose id differs from its file stem raises
\py!ValueError!.

\subsection{The engine script}

\begin{signature}
\py!circuit_at(inst, nv, dv, *, expand=False) -> Circuit!\\
\py!script(inst, nv, dv, *, command="decide") -> Script(text, literal, nodes, point_claim)!
\end{signature}

\py!circuit_at! calls the equation's circuit function with $d$ = \code{dv}, else \code{"d"} for a
general subject, always the fixed dimension otherwise; with $n$ = \code{nv} for an unrolled form
(\code{nv = None} raises \py!Unsupported("the circuit has no general form in n")!), \code{"n"}
for a boxed form with a size parameter, else \code{None}. With \code{expand} and a concrete $n$,
register counts and the scalar are evaluated at $n$. The $F$ edit, if any, is applied next;
with \code{expand} the counts and scalar are then evaluated at $n$ and the fixed $d$.

\begin{algorithm}[\texttt{script}]\label{alg:bench-script}
\begin{enumerate}
\item A fixed subject dimension sets \code{dv}. Build the circuit by \py!circuit_at!.
\item The script is \emph{literal} when the subject's dimension is fixed or the equation sets
  \code{literal}; a literal script needs a concrete $d$ (else \py!Unsupported("the equation needs a concrete dimension")!).
\item When not literal, translate with the bound values and no literal symbols; if the
  translator refuses with a message containing \code{literal} and $d$ is concrete, fall back to
  literal. A literal translation binds \code{d} as literal.
\item Build the specification with \py!spec_context(subject, nv, dv, literal)!; an $F$ edit's
  \py!apply_spec! returns it unchanged.
\item The header lines are \code{dim d} (with \code{= dv} when concrete) for a general,
  non-literal subject, and \code{count n} (with \code{= nv}) for a boxed equation with a size
  parameter.
\item The point clauses are \code{d = dv} for a general, non-literal subject at a concrete $d$,
  and \code{n = nv} for a boxed equation with a size parameter at a concrete $n$. A restriction
  clause on a pinned symbol is dropped when it holds at the pinned value and raises
  \py!Unsupported("n = v is outside the restriction ...")! otherwise; other clauses are kept.
\item The text is \code{@rhs =} with the indented header and specification, a blank line,
  \code{load} with the indented header and circuit body, a blank line, and
  \code{COMMAND against @rhs[ for CLAUSE, ...]} (kept clauses first, then point clauses), and a
  final newline. \code{nodes} is the translation's node count; \code{point_claim} is true when
  point clauses exist.
\end{enumerate}
\end{algorithm}

\begin{example}
\code{bell_register/main} at $(n,d) = (3,5)$:
\begin{lstlisting}[style=azx]
@rhs =
    dim d = 5
    count n = 3
    repeat n: Z(0 -> 2, d)
    scalar d^(-n/2)

load
    dim d = 5
    count n = 3
    g0 = X(0 -> 1, d)
    g1 = F(d)
    g2 = X(0 -> 1, d)
    g3 = Z(1 -> 2, d)
    g4 = X(2 -> 1, d)
    wire g0.out0 -- g1.in0, g1.out0 -- g3.in0, g2.out0 -- g4.in0, g3.out1 -- g4.in1
    repeat n: g0, g1, g2, g3, g4
    outputs g3.out0, g4.out0
    scalar d^(-n/2)

decide against @rhs for d = 5, n = 3
\end{lstlisting}
The symbolic version has \code{dim d}, \code{count n}, and the bare line \code{decide against @rhs}.
\end{example}

\subsection{Source digests}\label{sec:bench-sources}

\module{benchmarks.sources} computes what a run key and an aggregate filter compare.

\begin{definition}[Source digest]
The shared digest (cached per catalogue directory) is the SHA-256 hash, over the files in this
order, of each file's name (relative to the catalogue, or its bare name) followed by its bytes:
the helper modules \file{\_*.py} of the catalogue (sorted), every file in a catalogue
subdirectory outside \file{\_\_pycache\_\_} (sorted), then \file{ir.py}, \file{translate.py},
\file{subject.py}. \py!source_digest(inst, false=None)! continues a SHA-256 hash seeded with the
shared digest by the bytes of the subject's file; for an $F$ instance also by the bytes of
\file{breaks.py} and the JSON (sorted keys) of its $F$ record; it returns the first 10
hexadecimal digits.
\end{definition}

\section{Shared builders of the catalogue}\label{sec:bench-helpers}

The modules \file{\_*.py} of \file{benchmarks/subjects/} are not subjects; they provide gates,
specification writers, and code constructions. \file{\_common.py} only re-exports the IR gates,
\py!circuit!, \py!unrolled!, and the subject classes.

\subsection{Gates (\texttt{\_gates.py})}

Qubit gates: \py!H(q)! is $F$; \py!P(q, t)! is \py!ZPhase(q, {1: t})!, $\mathrm{diag}(1,
e^{2\pi i t})$; \py!S!, \py!Sdg!, \py!T!, \py!Tdg! are $P$ at $1/4$, $-1/4$, $1/8$, $-1/8$;
\py!XP(q, t)! is \py!XPhase(q, {1: t})!; \py!Y(q)! is $[Z, X]$ (the factor $i$ left to the
caller). Decompositions:
\begin{itemize}
\item \py!ccz(a, b, c)!: the 7-$T$ phase polynomial of \citet[Fig.~4.9]{nielsen2010quantum}:
  CX$(b,c)$, $T^\dagger(c)$, CX$(a,c)$, $T(c)$, CX$(b,c)$, $T^\dagger(c)$, CX$(a,c)$, $T(b)$,
  $T(c)$, CX$(a,b)$, $T(a)$, $T^\dagger(b)$, CX$(a,b)$.
\item \py!toffoli(a, b, c)!: $H(c)$, \py!ccz!, $H(c)$. \py!cswap(c, a, b)!: CX$(b,a)$,
  Toffoli$(c,a,b)$, CX$(b,a)$.
\item \py!cphase(a, b, t)!: $P(a, t/2)$, $P(b, t/2)$, CX$(a,b)$, $P(b,-t/2)$, CX$(a,b)$.
\item \py!ry(q, t)!: $e^{i\pi t}R_y(2\pi t)$ as $S^\dagger, H, P(t), H, S$.
  \py!cry(c, q, t)!: CX, \py!ry(-t/2)!, CX, \py!ry(t/2)!. \py!ch(c, q)!: \py!ry(1/8)!, CX,
  \py!ry(-1/8)!.
\item \py!AND_GATES(a, b, t)!: Gidney's 4-$T$ temporary AND~\cite{gidney2018halving}:
  Init$(t)$, $H(t)$, $T(t)$, CX$(a,t)$, CX$(b,t)$, CX$(t,a)$, CX$(t,b)$, $T^\dagger(a)$,
  $T^\dagger(b)$, $T(t)$, CX$(t,a)$, CX$(t,b)$, $H(t)$, $S(t)$.
\item \py!qft(qs, d)!: for each $j$, $F(q_j)$ and for each $k>j$ the \py!PhaseFn! on
  $(q_k, q_j)$ with $f(x) = x_0x_1/d^{k-j+1}$; then \py!Swap!$(q_j, q_{n-1-j})$ for $j<n/2$.
\item \py!inverse_gate(g)!: negated amount, $F^{-p \bmod 4}$, negated rule (\code{-(rule)}) or
  entries, $s\mapsto -s$, \py!Swap! itself, negated \py!PhaseFn!, negated \code{a} of
  \py!MCX!/\py!MCZ!, negated oracle function; \py!Perm! raises \py!TypeError!, any other gate
  \py!ValueError!. \py!inverse(gates)! inverts the reversed list.
\item \py!rename(gates, mapping)! renames the fields \code{q}, \code{c}, \code{t}, \code{y},
  \code{a}/\code{b} of \py!Swap!, and \code{qs}, \code{cs}, \code{xs}.
\item \py!dephase_lines(port, name, D)! returns the three DSL lines of a $Z$ copy with one port
  discarded.
\end{itemize}

\subsection{Specification writers (\texttt{\_spec.py})}

\py!defining(c, wires, gates, *, registers, dims, scalar)! returns the translation of the
defining gates under the specification context (literal $d$ exactly when \code{c.D} starts with
\code{literal}); \py!permutation! and \py!diagonal! wrap one \py!Perm! or \py!PhaseFn!.
\py!basis_state(c, values)! writes $\ket{v_1\cdots v_m}$ as X states with scalar
$d^{-m/2}$.

\py!Dsl! accumulates nodes (\code{v0}, \code{v1}, \dots), wires, inputs, outputs, and scalar
factors; \py!text()! writes declarations, one \code{wire} line, \code{inputs}, \code{outputs},
and \code{scalar (f1) * (f2) ...}. \py!bind(b, n, d)! joins $n$ ports of dimension $d$ into one
of $d^n$ (first most significant; one identity spider when $n=1$); \py!split(b, port, n, d)!
inverts it. \py!register_adder(c, n, d)! writes $\ket{a,b}\mapsto\ket{a, a+b \bmod d^n}$: both
registers bound, $a$ copied by \azx!Z(1 -> 2)!, added into $b$ by \azx!X(2 -> 1)!, split back,
scalar $(d^n)^{1/2}$. \py!pauli_gadget(b, ports, letters, turn, D)! appends
$e^{-2\pi i\,\mathit{turn}\,P}$ for a Pauli string: per wire a $Z$ port in its letter's basis
($Y$ by $S^\dagger$ phases $3/4$ and $1/4$ around $F$), an $X$ hub, the phase
$2\,\mathit{turn}$ at index 1 on a \azx!Z(1 -> 0)!, and scalar
$2^{(m-1)/2}e^{-2\pi i\,\mathit{turn}}$ \cite{cowtan2020phase}.

\subsection{Amy-Maslov-Mosca circuits (\texttt{\_qc.py})}

\file{subjects/amm/} holds ten \file{.qc} circuits of the Amy-Maslov-Mosca
benchmark~\cite{amy2014polynomial}, copied from \tool{Feynman}{https://github.com/meamy/feynman}
under its BSD license (\file{amm/LICENSE.md}). \py!read(name)! parses one file: \code{.v} lists
the wires, \code{.i} the inputs, gate lines lie between \code{BEGIN} and \code{END}. Names that
start with a digit get the prefix \code{w}. Gates: \code{H}, \code{X}, \code{T}, \code{T*},
\code{P} ($S$), \code{P*} ($S^\dagger$) on one wire; \code{tof} with 1, 2, or more wires as
$X$, CX, or MCX; \code{Z} or \code{Zd} with at least two wires as MCZ; anything else raises
\py!ValueError!. \py!bits(x, k)! and \py!number(bits)! convert integers and most-significant-first
bit lists.

\py!amm_subject(name, title, fn, note, clean=(), gates=())! builds subject
\code{amm_<name>} (with \texttt{\^{}} and \code{-} removed from the name) of category
\code{circuit}, $d = 2$, one fixed equation \code{main}: the circuit is the file's gates after
$\ket{0}$ on every ancilla (wires outside \code{.i}, and the \code{clean} inputs). The
specification is either the defining \code{gates} after the same initializations, or, when
\code{fn} maps input bits to every wire's output bit, the XOR-oracle construction
\py!_bitwise!: tabulate $\mathit{fn}$ over all inputs; an ancilla $w$ receives
\py!Oracle(ins, w, fn_w)!; a changed input $w$ receives its image on a fresh wire \code{w_out}
right after it, then $w$ is cleared by the oracle of its preimage value on the unchanged inputs
and fresh wires, and post-selected on $\bra{0}$.

\subsection{Qubit stabilizer codes (\texttt{\_stab.py})}

\py!Pauli(x, z, p)! is $i^pX^xZ^z$ on $n$ qubits, stored as supports; equality and hashing use
$(n, \mathit{xs}, \mathit{zs}, p)$. \py!sign()! is the sign $s$ of the Hermitian operator
$s\bigotimes P_j$: with $y = |\mathit{xs}\cap\mathit{zs}|$ and $e=(p-y)\bmod4$, $+1$ for $e=0$,
$-1$ for $e=2$, \py!ValueError! otherwise. \py!pauli(text)! parses \code{[+-]P1...Pn} with
$p = (2[\text{minus}] + \#Y)\bmod 4$. \py!mul(a, b)! returns $ab$ with
$p = p_a + p_b + 2|\mathit{zs}_a\cap\mathit{xs}_b|$; \py!commute(a, b)! tests
$|\mathit{xs}_a\cap\mathit{zs}_b| + |\mathit{zs}_a\cap\mathit{xs}_b|$ even.

\begin{algorithm}[Gottesman standard form, \texttt{standard\_form}]\label{alg:bench-standard}
Input: independent commuting generator strings~\cite{gottesman1998heisenberg}.
\begin{enumerate}
\item Check pairwise commutation (\py!ValueError("generators ... anticommute")!).
\item X pivots: while some row $i\ge r$ has an $X$ on a qubit at position $\ge r$ of the current
  qubit order, take the first such row and its least such position, swap the row to $r$ and the
  qubit to position $r$, multiply every other row with an $X$ there by row $r$, and increment
  $r$.
\item Z pivots: for $s = r,\dots,m-1$, likewise with $Z$ (no pivot raises
  \py!ValueError("dependent generators")!).
\item The order is (first $r$ X pivots, then $m-r$ Z pivots, then $k=n-m$ carriers). The
  reference state $b$ has 1 at each Z pivot whose row has sign $-1$.
\item Per carrier $c$: $\bar X$ has $X$ on $c$ and on the Z pivots whose rows have $Z$ on $c$,
  and $Z$ on each X pivot $f$ where $[c\in Z_f]$ differs from the parity of those hits in
  $Z_f$; $\bar Z$ has $Z$ on $c$ and on every X pivot whose row has $X$ on $c$.
\item Verify that every generator commutes with every logical and that $\bar X_i$ anticommutes
  with $\bar Z_j$ exactly when $i=j$ (\py!ValueError! otherwise).
\end{enumerate}
\end{algorithm}

\py!Code! holds $n$, the reduced and defining generators, the order, $r$, the reference state,
and the logicals; $m$, $k$, \py!carriers!, and \py!is_css()! (every defining generator X type or
Z type) derive from them. Data wires are \code{q0}, \dots.

\begin{itemize}
\item \py!encoder_gates(code)!: \py!Init(q_j, b_j)! on every non-carrier; per logical, CX from
  its carrier to each other qubit of $\bar X$'s X support; per X pivot $f$ of generator $g$:
  $F(q_f)$, \py!ZPhase! $\{1: (p \bmod 4)/4\}$ when $p\not\equiv0$, then per other qubit of the
  support CZ (Z part) and CX (X part) from $q_f$.
\item \py!projector_spec(code, c)!: $2^{r/2}P_S\bar X^x\ket{b}$. Inputs are $Z$ spiders in
  carrier order, each copied onto $\bar X$'s X support with phase $(p + 2\sum_{t\in Z}b_t)/4$;
  each qubit an $X$ spider (phase $1/2$ when $b_j=1$) of its feeds; then per defining generator a
  parity gadget: a $Z$ port per support qubit in its letter's basis, a hub \azx!X(m -> 0)! with
  phase $1/2$ for a negative sign. The scalar exponent is $r + \sum(m_j-1) + \sum(|g|-2)$
  halves. Logical X operators with crossing Z parts raise \py!ValueError!.
\item \py!css_spec! / \py!css_into!: the phase-free CSS normal form: one $Z$ spider per logical
  X row (an input) and per defining X stabilizer (a state, \azx!Z(0 -> m)!; a negative one
  raises \py!ValueError!), one $X$ spider per qubit joined to the rows covering it (phase $1/2$
  where $b_j=1$), scalar $2^{(\sum(m_j-1) - \#X)/2}$.
\item \py!syndrome_gates(code, anc)!: per defining generator, Init, $F$, the phase, controlled
  Z/X parts, $F$ on its ancilla; the ancilla ends in $\ket{s}$ for eigenvalue $(-1)^s$.
  \py!syndrome(code, e)! is the anticommutation bit vector.
\item \py!correction_circuit(code, error)!: encode, apply the error ($Y$ as $Z$ then $X$),
  measure syndromes onto \code{a0}\dots, then for the error and every single-qubit error
  (\py!single_errors!: qubit-major, letters X, Z, Y) with a nonzero syndrome not seen before,
  MCX (X, Y) and MCZ (Z, Y) on the ancillas at that syndrome; then the inverse encoder.
  \py!correction_spec!: carriers as identities, other data qubits in $\ket{b}$, ancillas in the
  syndrome, scalar $2^{-c/2}$.
\item \py!logical_phase(code, turns)!: the logical diagonal of a transversal phase gate on a CSS
  code, the phase of each codeword coset checked constant (\py!ValueError! otherwise).
  \py!pre_unary(text)! and \py!pre_diagonal(f, k)! write a logical gate before the
  specification's inputs; \py!Logical(key, physical, pre)! and \py!logical_spec! compose it
  with \py!css_into! (CSS codes only).
\item \py!code_subject(id, name, source, generators, *, correct=None, logical=(), note)!: a
  fixed qubit code subject with equation \code{encoding} (CSS normal form for a CSS code,
  projector form otherwise), \code{correction} (literal) when \code{correct} names an error,
  and one literal equation per logical gate; \code{clifford=True}.
\item \py!code_family(id, name, source, make, size_parameter, domain)!: equation
  \code{encoding} in unrolled form, the code of \py!make(n)! cached per $n$.
\item \py!css_generators!, \py!support_rows!, \py!hadamard_conjugate!, \py!pauli_string!
  (product of single-qubit factors, phases dropped), and \py!independent! (drops a generator in
  the GF(2) span of the kept ones, by a pivot basis on the binary vector $(x,z)$).
\end{itemize}

\subsection{Clifford synthesis and contraction (\texttt{\_clifford.py})}

\py!conjugate(p, g)! returns $UpU^\dagger$ for $g\in\{H, S, S^\dagger, X, Z\}$ on qubit $a$ or
CX$(a,b)$, updating supports and phase (H: swap $x_a,z_a$, $p\mathrel{+}=2x_az_a$; S: $z_a
\mathrel{\oplus}= x_a$, $p\mathrel{+}=x_a$ ($3x_a$ for $S^\dagger$); X: $p\mathrel{+}=2z_a$; Z:
$p\mathrel{+}=2x_a$; CX: $x_b\mathrel{\oplus}=x_a$, $z_a\mathrel{\oplus}=z_b$).
\py!synthesize(images)! reduces the image rows $(P_j, Q_j)$ qubit by qubit to $(X_j, Z_j)$ by
these conjugations and returns the inverse of the reducing gates ($S\leftrightarrow S^\dagger$);
a row left unreduced raises \py!ValueError!. \py!ir_gates! maps the result onto named wires.
\py!contract(gens, u, v)! returns the stabilizer generators of $\bra{\Phi^+}_{uv}\ket{\psi}$:
the GF(2) kernel combinations commuting with $X_uX_v$ and $Z_uZ_v$, each cleared on $u$ by those
operators, then restricted to the other qubits (\py!ValueError("legs not cleared")! when $v$
stays occupied). \py!trivial_on! and \py!restrict! give subgroups trivial on, and restrictions
to, qubit lists.

\subsection{Qudit codes (\texttt{\_qudit\_codes.py})}

Rows are dense lists or sparse \texttt{\{column: value\}} maps over $\Zd{d}$.
\py!css_encoder(n, stabs, logicals, carriers)!: for each stabilizer row choose the pivot column
off the carriers and earlier pivots with entry $\pm1$ and the fewest logical rows touching it,
normalize the row to $+1$ there, eliminate that column from every other stabilizer and logical
row; check each logical row is 1 on its carrier; order the logicals so that a logical is ready
when no other pending logical has a nonzero entry on its carrier (none ready raises
\py!ValueError("carrier dependencies are cyclic")!); emit Init on non-carriers, carrier fans
($|a|$ CX of the sign of $a$), then $F$ and a fan from each pivot. \py!css_spec! writes the
integer CSS normal form $\ket{x}\mapsto d^{-r/2}\sum_s\ket{xL+sS}$: a $Z$ spider per logical row
(input, carrier order) and per stabilizer row (state), an $X$ spider per qudit, an entry $a$ as
$|a|$ wires through the antipode $F^2$ when negative, scalar $d^{(\sum(m_j-1) - r)/2}$.
\py!graph_encoder! and \py!graph_spec! do the same for graph codes
$\ket{x}\mapsto Z^{xN}\ket{G}$~\cite{schlingemann2001quantum}: $\ket{+}$ on non-carriers, CZ from
each carrier to its logical set, $F$ on carriers, the graph's CZ edges; specification $Z$
spiders joined by $F$ edges, scalar $d^{(\#F - n)/2}$.

\subsection{Lattices (\texttt{\_lattice.py})}

\py!Honeycomb(L)! ($L$ a multiple of 3): $2L^2$ sites $A(x,y)=2(xL+y)$, $B=A+1$ (coordinates
mod $L$); hexagon $(x,y)$ on $A(x,y), B(x,y), A(x,y{+}1), B(x{-}1,y{+}1), A(x{-}1,y{+}1),
B(x{-}1,y)$, colored $(x-y)\bmod3$; $3L^2$ edges with direction letters Z ($A(x,y)$--$B(x,y)$),
X ($A(x,y)$--$B(x{-}1,y)$), Y ($A(x,y)$--$B(x,y{-}1)$) and the color of neither bordering
hexagon. \py!Cubic(L)!: $3L^3$ edge qubits, \py!edge(p, d)! $= 3((xL+y)L+z)+d$, and
\py!step(p, d, s)!.

\section{The subject catalogue}\label{sec:bench-catalogue}

The catalogue holds 207 subjects: 47 algorithms, 47 codes, 42 protocols, and 71 circuits, with
231 equations and 234 $S$ instances (three restricted claims). \Cref{tab:bench-catalogue} lists every subject with its equations, its domain in
$n$ (what $n$ counts in parentheses) and $d$, and the paper of the algorithm, code, or protocol.
An equation's marks give its form (\code{b} boxed, \code{u} unrolled, \code{f} fixed), \code{L}
for a literal equation, and \code{R} for an equation with a reference circuit; $\dagger$ marks a
Clifford subject. Restricted claims follow the equations.

Each subject file states, in its own docstring, the circuit and the specification it is checked
against. The specification is written from the definition of the algorithm's result, never from
the circuit: a basis state, a spider normal form, a phase table, a CSS or projector normal form,
a defining gate, a dense weight table, or a closed-form probability distribution as a weighted
spider purified and traced (for example \code{chsh_game}: the bit with distribution
$(\cos^2(\pi/8), \sin^2(\pi/8))$ as \azx!Z(0 -> 1, D) weight{1: sqrt(2) - 1}! dephased, scalar
$\sqrt{2+\sqrt2}/2$).

{\small
\begin{xltabular}{\textwidth}{@{}>{\raggedright\arraybackslash}p{0.2\textwidth}
  >{\raggedright\arraybackslash}X
  >{\raggedright\arraybackslash}p{0.16\textwidth}
  >{\raggedright\arraybackslash}p{0.15\textwidth}
  >{\raggedright\arraybackslash}p{0.08\textwidth}
  >{\raggedright\arraybackslash}p{0.07\textwidth}@{}}
\caption{The subject catalogue.}\label{tab:bench-catalogue}\\
\toprule
Subject & Name & Equations & $n$ & $d$ & Source\\
\midrule\endfirsthead
\toprule
Subject & Name & Equations & $n$ & $d$ & Source\\
\midrule\endhead
\bottomrule\endfoot
\multicolumn{6}{@{}l}{\textbf{Category algorithm} (47 subjects)}\\
\midrule
\texttt{amplitude\_\allowbreak{}amplification} & Amplitude amplification, exact single round & \texttt{main}\,[f] & -- & $2$ & \cite{brassard2002quantum}\\
\texttt{amplitude\_\allowbreak{}estimation} & Amplitude estimation & \texttt{main}\,[f] & -- & $2$ & \cite{brassard2002quantum}\\
\texttt{bernstein\_\allowbreak{}vazirani} & Bernstein-Vazirani & \texttt{main}\,[bR]; restricted: n >= 1; d >= 2 & $n\ge 0$ (register qudits) & $d\ge 1$ & \cite{bernstein1997complexity}\\
\texttt{bernstein\_\allowbreak{}vazirani\_\allowbreak{}phase} & Bernstein-Vazirani, phase oracle & \texttt{main}\,[bR] & $n\ge 0$ (register qudits) & $d\ge 1$ & \cite{bernstein1997complexity}\\
\texttt{block\_\allowbreak{}encoding} & Block encoding by LCU & \texttt{main}\,[f] & -- & $2$ & \cite{gilyen2019quantum}\\
\texttt{bv\_\allowbreak{}secret\_\allowbreak{}pattern} & Bernstein-Vazirani, patterned secret & \texttt{main}\,[bR] & $n\ge 1$ (register qudits) & $d\ge 1$ & \cite{bernstein1997complexity}\\
\texttt{coined\_\allowbreak{}quantum\_\allowbreak{}walk} & Coined quantum walk on a cycle & \texttt{main}\,[bR] & $n\ge 1$ (position bits) & $2$ & \cite{aharonov1993quantum}\\
\texttt{cooley\_\allowbreak{}tukey\_\allowbreak{}qft} & Cooley-Tukey mixed-radix QFT & \texttt{main}\,[f] & -- & $6$ & \cite{cleve2000fast}\\
\texttt{deutsch} & Deutsch's algorithm & \texttt{f00}\,[f], \texttt{f01}\,[f], \texttt{f10}\,[f], \texttt{f11}\,[f] & -- & $d\ge 1$ & \cite{deutsch1985quantum}\\
\texttt{deutsch\_\allowbreak{}jozsa} & Deutsch-Jozsa & \texttt{constant}\,[b], \texttt{balanced}\,[b] & $n\ge 0$ (register qudits) & $d\ge 1$ & \cite{deutsch1992rapid}\\
\texttt{dihedral\_\allowbreak{}qft} & Quantum Fourier transform over the dihedral group D4 & \texttt{main}\,[f] & -- & $2$ & \cite{puschel1999fast}\\
\texttt{discrete\_\allowbreak{}log} & Shor's discrete logarithm, cyclic group & \texttt{main}\,[f] & -- & $d\ge 1$ & \cite{shor1997polynomial}\\
\texttt{dqc1} & DQC1 trace estimation & \texttt{main}\,[b] & $n\ge 1$ (mixed qubits) & $2$ & \cite{knill1998power}\\
\texttt{fixed\_\allowbreak{}point\_\allowbreak{}search} & Fixed-point quantum search & \texttt{main}\,[f] & -- & $2$ & \cite{grover2005fixed}\\
\texttt{forrelation} & Forrelation & \texttt{main}\,[b] & $n\ge 0$ (qudit pairs) & $d\ge 1$ & \cite{aaronson2015forrelation}\\
\texttt{grover} & Grover search, two qubits & \texttt{main}\,[fR] & -- & $2$ & \cite{grover1996fast}\\
\texttt{grover\_\allowbreak{}marked} & Grover search, every marked item & \texttt{m00}\,[fR], \texttt{m01}\,[fR], \texttt{m10}\,[fR], \texttt{m11}\,[fR] & -- & $2$ & \cite{grover1996fast}\\
\texttt{grover\_\allowbreak{}ququart} & Grover search on a ququart & \texttt{main}\,[f] & -- & $4$ & \cite{grover1996fast}\\
\texttt{grover\_\allowbreak{}rudolph} & Grover-Rudolph state preparation & \texttt{main}\,[f] & -- & $2$ & \cite{grover2002creating}\\
\texttt{grover\_\allowbreak{}two\_\allowbreak{}marked} & Grover search, two marked items & \texttt{main}\,[f] & -- & $2$ & \cite{boyer1998tight}\\
\texttt{hadamard\_\allowbreak{}test} & Hadamard test & \texttt{main}\,[f] & -- & $2$ & \cite{aharonov2009polynomial}\\
\texttt{hhl} & HHL linear-system solver, 2 x 2 & \texttt{main}\,[f] & -- & $2$ & \cite{harrow2009quantum}\\
\texttt{hidden\_\allowbreak{}linear\_\allowbreak{}function} & Hidden linear function & \texttt{main}\,[f] & -- & $2$ & \cite{bravyi2018quantum}\\
\texttt{hidden\_\allowbreak{}shift} & Hidden shift, inner-product bent function & \texttt{main}\,[bR] & $n\ge 0$ (qudit pairs) & $d\ge 1$ & \cite{vandam2003quantum}\\
\texttt{inverse\_\allowbreak{}qft} & Inverse quantum Fourier transform & \texttt{main}\,[uL] & $n\ge 1$ (qudits) & $d\ge 2$ & \cite{nielsen2010quantum}\\
\texttt{iqp\_\allowbreak{}circuit} & IQP circuit & \texttt{main}\,[b] & $n\ge 1$ (qubits after the first) & $2$ & \cite{shepherd2009temporally}\\
\texttt{iterative\_\allowbreak{}phase\_\allowbreak{}estimation} & Iterative phase estimation & \texttt{main}\,[f] & -- & $2$ & \cite{kitaev1995quantum}\\
\texttt{lcu\_\allowbreak{}shift\_\allowbreak{}average} & Linear combination of unitaries, shift average & \texttt{main}\,[f] & -- & $d\ge 1$ & \cite{childs2012hamiltonian}\\
\texttt{legendre\_\allowbreak{}hidden\_\allowbreak{}shift} & Shifted Legendre symbol (van Dam-Hallgren-Ip) & \texttt{main}\,[f] & -- & $2$ & \cite{vandam2006quantum}\\
\texttt{period\_\allowbreak{}finding} & Period finding & \texttt{main}\,[fL] & -- & $d\ge 2$ & \cite{shor1997polynomial}\\
\texttt{phase\_\allowbreak{}estimation} & Quantum phase estimation & \texttt{main}\,[uL] & $n\ge 1$ (clock qudits) & $d\ge 2$ & \cite{kitaev1995quantum}\\
\texttt{phase\_\allowbreak{}estimation\_\allowbreak{}digit} & Phase estimation, one digit & \texttt{main}\,[fR] & -- & $d\ge 1$ & \cite{kitaev1995quantum}\\
\texttt{qaoa\_\allowbreak{}maxcut} & QAOA MaxCut on a cycle & \texttt{main}\,[b] & $n\ge 1$ (qubits after the first two) & $2$ & \cite{farhi2014quantum}\\
\texttt{qft} & Quantum Fourier transform & \texttt{main}\,[uL] & $n\ge 1$ (qudits) & $d\ge 2$ & \cite{nielsen2010quantum}\\
\texttt{qft\_\allowbreak{}roundtrip} & QFT round trip & \texttt{main}\,[uLR] & $n\ge 1$ (qudits) & $d\ge 2$ & \cite{nielsen2010quantum}\\
\texttt{qft\_\allowbreak{}zero\_\allowbreak{}state} & QFT of the zero state & \texttt{main}\,[uLR] & $n\ge 1$ (qudits) & $d\ge 2$ & \cite{nielsen2010quantum}\\
\texttt{quantum\_\allowbreak{}counting} & Quantum counting & \texttt{main}\,[f] & -- & $2$ & \cite{brassard1998quantum}\\
\texttt{quantum\_\allowbreak{}walk\_\allowbreak{}hypercube} & Quantum walk on the hypercube, perfect state transfer & \texttt{main}\,[b] & $n\ge 0$ (cube dimension) & $2$ & \cite{childs2003exponential}\\
\texttt{qubitization\_\allowbreak{}chebyshev} & Qubitization, Chebyshev block encoding & \texttt{main}\,[b] & $n\ge 0$ (walk step pairs after the first step) & $2$ & \cite{low2019hamiltonian}\\
\texttt{semiclassical\_\allowbreak{}qft} & Semiclassical QFT & \texttt{main}\,[uL] & $n\ge 1$ (qudits) & $d\ge 2$ & \cite{griffiths1996semiclassical}\\
\texttt{shor\_\allowbreak{}factoring} & Shor's algorithm, N = 15 & \texttt{main}\,[f] & -- & $2$ & \cite{shor1997polynomial}\\
\texttt{simon} & Simon's algorithm & \texttt{main}\,[b] & $n\ge 0$ (register qudits after the first) & $d\ge 1$ & \cite{simon1997power}\\
\texttt{swap\_\allowbreak{}test} & Swap test & \texttt{main}\,[f] & -- & $2$ & \cite{buhrman2001quantum}\\
\texttt{swap\_\allowbreak{}test\_\allowbreak{}register} & Swap test on registers & \texttt{main}\,[u] & $n\ge 0$ (qubits per register) & $2$ & \cite{buhrman2001quantum}\\
\texttt{szegedy\_\allowbreak{}walk} & Szegedy walk, two-state Markov chain & \texttt{main}\,[f] & -- & $2$ & \cite{szegedy2004quantum}\\
\texttt{trotter\_\allowbreak{}ising} & Trotter step, transverse-field Ising chain & \texttt{main}\,[b] & $n\ge 1$ (qubits after the first) & $2$ & \cite{lloyd1996universal}\\
\texttt{vqe\_\allowbreak{}ucc} & VQE unitary coupled-cluster ansatz (H2) & \texttt{main}\,[f] & -- & $2$ & \cite{peruzzo2014variational}\\
\midrule
\multicolumn{6}{@{}l}{\textbf{Category code} (47 subjects)}\\
\midrule
\texttt{bacon\_\allowbreak{}shor\_\allowbreak{}code}${}^{\dagger}$ & Bacon-Shor subsystem code & \texttt{encoding}\,[f] & -- & $2$ & \cite{bacon2006operator}\\
\texttt{bacon\_\allowbreak{}shor\_\allowbreak{}family}${}^{\dagger}$ & Bacon-Shor code family & \texttt{encoding}\,[u] & $n\ge 2$ (grid side) & $2$ & \cite{bacon2006operator}\\
\texttt{bit\_\allowbreak{}flip\_\allowbreak{}correction} & Bit-flip code correction & \texttt{errornone}\,[f], \texttt{error0}\,[f], \texttt{error1}\,[f], \texttt{error2}\,[f] & -- & $2$ & \cite{nielsen2010quantum}\\
\texttt{bivariate\_\allowbreak{}bicycle\_\allowbreak{}code}${}^{\dagger}$ & Bivariate bicycle code & \texttt{encoding}\,[f] & -- & $2$ & \cite{bravyi2024high}\\
\texttt{color\_\allowbreak{}code}${}^{\dagger}$ & Color code (6.6.6, torus) & \texttt{encoding}\,[f] & -- & $2$ & \cite{bombin2006topological}\\
\texttt{cube\_\allowbreak{}code}${}^{\dagger}$ & [[8,3,2]] cube code & \texttt{encoding}\,[f], \texttt{transversal\_\allowbreak{}t}\,[fL] & -- & $2$ & \cite{campbell2016smallest}\\
\texttt{dual\_\allowbreak{}rail\_\allowbreak{}code} & Dual-rail code & \texttt{main}\,[f] & -- & $2$ & \cite{knill2001scheme}\\
\texttt{ea\_\allowbreak{}code} & Entanglement-assisted [[4,1,3;1]] code & \texttt{encoding}\,[f] & -- & $2$ & \cite{brun2006correcting}\\
\texttt{five\_\allowbreak{}qubit\_\allowbreak{}code}${}^{\dagger}$ & Five-qubit perfect code & \texttt{encoding}\,[f], \texttt{correction}\,[fL] & -- & $2$ & \cite{laflamme1996perfect}\\
\texttt{five\_\allowbreak{}qudit\_\allowbreak{}code}${}^{\dagger}$ & Five-qudit code, graph form & \texttt{encoding}\,[f] & -- & $d\ge 1$ & \cite{schlingemann2001quantum}\\
\texttt{four\_\allowbreak{}qubit\_\allowbreak{}code}${}^{\dagger}$ & [[4,2,2]] error-detecting code & \texttt{encoding}\,[f] & -- & $2$ & \cite{vaidman1996error}\\
\texttt{golay\_\allowbreak{}code}${}^{\dagger}$ & Quantum Golay [[23,1,7]] code & \texttt{encoding}\,[f] & -- & $2$ & \cite{steane2003overhead}\\
\texttt{gottesman\_\allowbreak{}code}${}^{\dagger}$ & Gottesman [[8,3,3]] code & \texttt{encoding}\,[f], \texttt{correction}\,[fL] & -- & $2$ & \cite{gottesman1996class}\\
\texttt{haah\_\allowbreak{}code}${}^{\dagger}$ & Haah's cubic code & \texttt{encoding}\,[f] & -- & $2$ & \cite{haah2011local}\\
\texttt{happy\_\allowbreak{}code}${}^{\dagger}$ & Holographic pentagon (HaPPY) code & \texttt{encoding}\,[f] & -- & $2$ & \cite{pastawski2015holographic}\\
\texttt{heavy\_\allowbreak{}hexagon\_\allowbreak{}code}${}^{\dagger}$ & Heavy-hexagon subsystem code & \texttt{encoding}\,[f] & -- & $2$ & \cite{chamberland2020topological}\\
\texttt{honeycomb\_\allowbreak{}code}${}^{\dagger}$ & Honeycomb Floquet code (instantaneous stabilizer group) & \texttt{encoding}\,[f] & -- & $2$ & \cite{hastings2021dynamically}\\
\texttt{hypergraph\_\allowbreak{}product\_\allowbreak{}code}${}^{\dagger}$ & Hypergraph product code [[13,1,3]] & \texttt{encoding}\,[f] & -- & $2$ & \cite{tillich2014quantum}\\
\texttt{iceberg\_\allowbreak{}code}${}^{\dagger}$ & Iceberg [[6,4,2]] code & \texttt{encoding}\,[f] & -- & $2$ & \cite{self2024protecting}\\
\texttt{iceberg\_\allowbreak{}family}${}^{\dagger}$ & Iceberg code family & \texttt{encoding}\,[u] & $n\ge 2$ (qubit pairs) & $2$ & \cite{self2024protecting}\\
\texttt{parity\_\allowbreak{}check} & Parity check & \texttt{syndrome}\,[b] & $n\ge 0$ (data qudits) & $d\ge 1$ & \cite{macwilliams1977theory}\\
\texttt{parity\_\allowbreak{}decoding} & Parity check uncomputation & \texttt{decoding}\,[bR] & $n\ge 0$ (data qudits) & $d\ge 1$ & \cite{macwilliams1977theory}\\
\texttt{phase\_\allowbreak{}flip\_\allowbreak{}code} & Qudit phase-flip code & \texttt{encoding}\,[b] & $n\ge 0$ (qudits after the first) & $d\ge 1$ & \cite{nielsen2010quantum}\\
\texttt{phase\_\allowbreak{}flip\_\allowbreak{}decoding} & Qudit phase-flip code decoding & \texttt{decoding}\,[bR] & $n\ge 0$ (qudits after the first) & $d\ge 1$ & \cite{nielsen2010quantum}\\
\texttt{polynomial\_\allowbreak{}code} & Quantum polynomial code & \texttt{main}\,[f] & -- & $5$ & \cite{aharonov1997fault}\\
\texttt{quantum\_\allowbreak{}hamming\_\allowbreak{}code}${}^{\dagger}$ & Quantum Hamming [[15,7,3]] code & \texttt{encoding}\,[f] & -- & $2$ & \cite{steane1996simple}\\
\texttt{quantum\_\allowbreak{}hamming\_\allowbreak{}family}${}^{\dagger}$ & Quantum Hamming code family & \texttt{encoding}\,[u] & $n\ge 3$ (check rows) & $2$ & \cite{steane1996simple}\\
\texttt{qudit\_\allowbreak{}four\_\allowbreak{}code}${}^{\dagger}$ & Qudit [[4,2,2]] code & \texttt{encoding}\,[f] & -- & $d\ge 1$ & \cite{grassl1997codes}\\
\texttt{qudit\_\allowbreak{}toric\_\allowbreak{}code}${}^{\dagger}$ & Qudit toric code, 2 x 2 torus & \texttt{encoding}\,[f] & -- & $d\ge 1$ & \cite{bullock2007qudit}\\
\texttt{qudit\_\allowbreak{}toric\_\allowbreak{}family}${}^{\dagger}$ & Qudit toric code family & \texttt{encoding}\,[u] & $n\ge 2$ (torus side) & $d\ge 1$ & \cite{bullock2007qudit}\\
\texttt{qutrit\_\allowbreak{}code}${}^{\dagger}$ & Qutrit [[3,1,2]] code & \texttt{encoding}\,[f] & -- & $3$ & \cite{cleve1999share}\\
\texttt{reed\_\allowbreak{}muller\_\allowbreak{}code}${}^{\dagger}$ & Quantum Reed-Muller [[15,1,3]] code & \texttt{encoding}\,[f], \texttt{transversal\_\allowbreak{}t}\,[fL] & -- & $2$ & \cite{knill1996threshold}\\
\texttt{repetition\_\allowbreak{}code} & Qudit repetition code & \texttt{encoding}\,[bR], \texttt{syndrome}\,[b] & $n\ge 0$ (qudits after the first) & $d\ge 1$ & \cite{nielsen2010quantum}\\
\texttt{repetition\_\allowbreak{}decoding} & Qudit repetition code decoding & \texttt{decoding}\,[bR] & $n\ge 0$ (qudits after the first) & $d\ge 1$ & \cite{nielsen2010quantum}\\
\texttt{repetition\_\allowbreak{}error\_\allowbreak{}syndrome} & Repetition code error syndrome & \texttt{main}\,[u] & $n\ge 0$ (qudits beyond two) & $d\ge 1$ & \cite{nielsen2010quantum}\\
\texttt{rotated\_\allowbreak{}surface\_\allowbreak{}code}${}^{\dagger}$ & Rotated surface code, distance 3 & \texttt{encoding}\,[f], \texttt{correction}\,[fL] & -- & $2$ & \cite{horsman2012surface}\\
\texttt{shor\_\allowbreak{}code} & Qudit Shor code & \texttt{encoding}\,[f] & -- & $d\ge 1$ & \cite{shor1995scheme}\\
\texttt{steane\_\allowbreak{}code}${}^{\dagger}$ & Steane [[7,1,3]] code & \texttt{encoding}\,[f], \texttt{correction}\,[fL], \texttt{transversal\_\allowbreak{}h}\,[fL] & -- & $2$ & \cite{steane1996multiple}\\
\texttt{surface\_\allowbreak{}code\_\allowbreak{}family}${}^{\dagger}$ & Rotated surface code family & \texttt{encoding}\,[u] & $n\ge 3$, step 2 (code distance) & $2$ & \cite{horsman2012surface}\\
\texttt{tesseract\_\allowbreak{}code}${}^{\dagger}$ & Tesseract [[16,6,4]] code & \texttt{encoding}\,[f] & -- & $2$ & \cite{reichardt2024demonstration}\\
\texttt{toric\_\allowbreak{}code}${}^{\dagger}$ & Toric code, 2 x 2 torus & \texttt{encoding}\,[f] & -- & $2$ & \cite{kitaev2003fault}\\
\texttt{toric\_\allowbreak{}code\_\allowbreak{}3d}${}^{\dagger}$ & 3D toric code & \texttt{encoding}\,[f] & -- & $2$ & \cite{castelnovo2008topological}\\
\texttt{toric\_\allowbreak{}code\_\allowbreak{}family}${}^{\dagger}$ & Toric code family & \texttt{encoding}\,[u] & $n\ge 2$ (torus side) & $2$ & \cite{kitaev2003fault}\\
\texttt{transversal\_\allowbreak{}cnot} & Transversal CNOT, Steane code & \texttt{main}\,[f] & -- & $2$ & \cite{steane1996multiple}\\
\texttt{x\_\allowbreak{}cube\_\allowbreak{}code}${}^{\dagger}$ & X-cube fracton code & \texttt{encoding}\,[f] & -- & $2$ & \cite{vijay2016fracton}\\
\texttt{xy\_\allowbreak{}surface\_\allowbreak{}code}${}^{\dagger}$ & XY surface code, distance 3 & \texttt{encoding}\,[f] & -- & $2$ & \cite{tuckett2018ultrahigh}\\
\texttt{xzzx\_\allowbreak{}surface\_\allowbreak{}code}${}^{\dagger}$ & XZZX surface code, distance 3 & \texttt{encoding}\,[f] & -- & $2$ & \cite{bonillaataides2021xzzx}\\
\midrule
\multicolumn{6}{@{}l}{\textbf{Category protocol} (42 subjects)}\\
\midrule
\texttt{b92} & B92 key distribution & \texttt{main}\,[f] & -- & $2$ & \cite{bennett1992quantumb92}\\
\texttt{bb84} & BB84 key distribution & \texttt{main}\,[f] & -- & $2$ & \cite{bennett1984quantum}\\
\texttt{bb84\_\allowbreak{}intercept\_\allowbreak{}resend} & BB84 under intercept-resend & \texttt{main}\,[f] & -- & $2$ & \cite{bennett1984quantum}\\
\texttt{bbm92} & BBM92 key distribution & \texttt{main}\,[f] & -- & $2$ & \cite{bennett1992quantum}\\
\texttt{blind\_\allowbreak{}computation} & Universal blind quantum computation & \texttt{main}\,[f] & -- & $2$ & \cite{broadbent2009universal}\\
\texttt{chsh\_\allowbreak{}game} & CHSH game & \texttt{main}\,[f] & -- & $2$ & \cite{clauser1969proposed}\\
\texttt{dejmps\_\allowbreak{}purification} & DEJMPS entanglement purification & \texttt{main}\,[f] & -- & $2$ & \cite{deutsch1996quantum}\\
\texttt{dephasing\_\allowbreak{}twirl} & Dephasing twirl & \texttt{main}\,[f] & -- & $d\ge 1$ & \cite{nielsen2010quantum}\\
\texttt{e91} & E91 key distribution & \texttt{main}\,[f] & -- & $2$ & \cite{ekert1991quantum}\\
\texttt{entanglement\_\allowbreak{}purification} & BBPSSW entanglement purification & \texttt{main}\,[f] & -- & $2$ & \cite{bennett1996purification}\\
\texttt{entanglement\_\allowbreak{}swapping} & Entanglement swapping & \texttt{main}\,[f] & -- & $d\ge 1$ & \cite{zukowski1993event}\\
\texttt{gate\_\allowbreak{}teleportation} & Gate teleportation & \texttt{main}\,[f] & -- & $d\ge 1$ & \cite{gottesman1999demonstrating}\\
\texttt{ghz\_\allowbreak{}fusion} & GHZ fusion & \texttt{main}\,[f] & -- & $d\ge 1$ & \cite{browne2005resource}\\
\texttt{ghz\_\allowbreak{}game} & GHZ game & \texttt{xxx}\,[f], \texttt{xyy}\,[f] & -- & $2$ & \cite{mermin1990quantum}\\
\texttt{ghz\_\allowbreak{}key\_\allowbreak{}agreement} & GHZ conference key agreement & \texttt{main}\,[b] & $n\ge 0$ (parties after the first) & $d\ge 1$ & \cite{epping2017multi}\\
\texttt{ghz\_\allowbreak{}metrology} & GHZ phase estimation & \texttt{main}\,[b] & $n\ge 0$ (qubits after the first) & $2$ & \cite{bollinger1996optimal}\\
\texttt{hbb\_\allowbreak{}secret\_\allowbreak{}sharing} & HBB quantum secret sharing & \texttt{main}\,[f] & -- & $d\ge 1$ & \cite{hillery1999quantum}\\
\texttt{magic\_\allowbreak{}square\_\allowbreak{}game} & Mermin-Peres magic square game & \texttt{r1c3}\,[f], \texttt{r3c1}\,[f], \texttt{r3c3}\,[f] & -- & $2$ & \cite{mermin1990simple}\\
\texttt{magic\_\allowbreak{}state\_\allowbreak{}distillation} & 15-to-1 magic state distillation & \texttt{main}\,[f] & -- & $2$ & \cite{bravyi2005universal}\\
\texttt{magic\_\allowbreak{}state\_\allowbreak{}injection} & Magic state injection & \texttt{main}\,[f] & -- & $2$ & \cite{bravyi2005universal}\\
\texttt{mbqc\_\allowbreak{}cnot} & Measurement-based CNOT & \texttt{main}\,[f] & -- & $d\ge 1$ & \cite{raussendorf2003measurement}\\
\texttt{mbqc\_\allowbreak{}j\_\allowbreak{}gate} & Measurement-based J gate & \texttt{main}\,[f] & -- & $d\ge 1$ & \cite{raussendorf2001one}\\
\texttt{mbqc\_\allowbreak{}wire} & Measurement-based quantum wire & \texttt{main}\,[b] & $n\ge 0$ (cluster qudits) & $d\ge 1$ & \cite{raussendorf2003measurement}\\
\texttt{mdi\_\allowbreak{}qkd} & Measurement-device-independent key distribution & \texttt{main}\,[f] & -- & $2$ & \cite{lo2012measurement}\\
\texttt{nonlocal\_\allowbreak{}cnot} & Nonlocal CNOT & \texttt{main}\,[f] & -- & $d\ge 1$ & \cite{eisert2000optimal}\\
\texttt{nonlocal\_\allowbreak{}cz} & Nonlocal CZ & \texttt{main}\,[f] & -- & $d\ge 1$ & \cite{eisert2000optimal}\\
\texttt{one\_\allowbreak{}bit\_\allowbreak{}teleportation} & One-bit teleportation & \texttt{main}\,[f] & -- & $d\ge 1$ & \cite{zhou2000methodology}\\
\texttt{one\_\allowbreak{}time\_\allowbreak{}pad} & Quantum one-time pad & \texttt{correctness}\,[f], \texttt{secrecy}\,[f] & -- & $d\ge 1$ & \cite{ambainis2000private}\\
\texttt{random\_\allowbreak{}access\_\allowbreak{}code} & Quantum random access code & \texttt{main}\,[f] & -- & $2$ & \cite{ambainis2002dense}\\
\texttt{randomized\_\allowbreak{}compiling} & Randomized compiling & \texttt{main}\,[f] & -- & $d\ge 1$ & \cite{wallman2016noise}\\
\texttt{remote\_\allowbreak{}state\_\allowbreak{}preparation} & Remote state preparation & \texttt{main}\,[f] & -- & $d\ge 1$ & \cite{bennett2001remote}\\
\texttt{repeater\_\allowbreak{}chain} & Quantum repeater chain & \texttt{main}\,[b] & $n\ge 0$ (swaps) & $d\ge 1$ & \cite{briegel1998quantum}\\
\texttt{six\_\allowbreak{}state\_\allowbreak{}qkd} & Six-state key distribution & \texttt{main}\,[f] & -- & $2$ & \cite{bruss1998optimal}\\
\texttt{superdense\_\allowbreak{}coding} & Superdense coding & \texttt{main}\,[f] & -- & $d\ge 1$ & \cite{bennett1992communication}\\
\texttt{superdense\_\allowbreak{}register} & Register superdense coding & \texttt{main}\,[b] & $n\ge 0$ (dit pairs) & $d\ge 1$ & \cite{bennett1992communication}\\
\texttt{teleportation} & Quantum teleportation & \texttt{main}\,[f] & -- & $d\ge 1$ & \cite{bennett1993teleporting}\\
\texttt{teleportation\_\allowbreak{}chain} & Teleportation chain & \texttt{main}\,[b] & $n\ge 0$ (relays) & $d\ge 1$ & \cite{bennett1993teleporting}\\
\texttt{teleportation\_\allowbreak{}register} & Register teleportation & \texttt{main}\,[b] & $n\ge 0$ (qudits) & $d\ge 1$ & \cite{bennett1993teleporting}\\
\texttt{threshold\_\allowbreak{}secret\_\allowbreak{}sharing} & Threshold quantum secret sharing & \texttt{main}\,[f] & -- & $3$ & \cite{cleve1999share}\\
\texttt{universal\_\allowbreak{}cloner} & Universal quantum cloner & \texttt{main}\,[f] & -- & $2$ & \cite{werner1998optimal}\\
\texttt{wiesner\_\allowbreak{}money} & Wiesner quantum money & \texttt{main}\,[b] & $n\ge 0$ (qubits) & $2$ & \cite{wiesner1983conjugate}\\
\texttt{y\_\allowbreak{}state\_\allowbreak{}injection} & Y-state injection & \texttt{main}\,[f] & -- & $2$ & \cite{aliferis2006quantum}\\
\midrule
\multicolumn{6}{@{}l}{\textbf{Category circuit} (71 subjects)}\\
\midrule
\texttt{amm\_\allowbreak{}barenco\_\allowbreak{}tof\_\allowbreak{}3} & 3-controlled Toffoli, dirty ancillas (Amy-Maslov-Mosca benchmark barenco\_tof\_3) & \texttt{main}\,[f] & -- & $2$ & \cite{amy2014polynomial}\\
\texttt{amm\_\allowbreak{}barenco\_\allowbreak{}tof\_\allowbreak{}4} & 4-controlled Toffoli, dirty ancillas (Amy-Maslov-Mosca benchmark barenco\_tof\_4) & \texttt{main}\,[f] & -- & $2$ & \cite{amy2014polynomial}\\
\texttt{amm\_\allowbreak{}barenco\_\allowbreak{}tof\_\allowbreak{}5} & 5-controlled Toffoli, dirty ancillas (Amy-Maslov-Mosca benchmark barenco\_tof\_5) & \texttt{main}\,[f] & -- & $2$ & \cite{amy2014polynomial}\\
\texttt{amm\_\allowbreak{}gf24\_\allowbreak{}mult} & GF(2\^{}4) multiplier (Amy-Maslov-Mosca benchmark gf2\^{}4\_mult) & \texttt{main}\,[f] & -- & $2$ & \cite{amy2014polynomial}\\
\texttt{amm\_\allowbreak{}hwb6} & Hidden weighted bit (Amy-Maslov-Mosca benchmark hwb6) & \texttt{main}\,[f] & -- & $2$ & \cite{amy2014polynomial}\\
\texttt{amm\_\allowbreak{}mod5\_\allowbreak{}4} & Divisibility by five (Amy-Maslov-Mosca benchmark mod5\_4) & \texttt{main}\,[f] & -- & $2$ & \cite{amy2014polynomial}\\
\texttt{amm\_\allowbreak{}tof\_\allowbreak{}3} & 3-controlled Toffoli, clean ancillas (Amy-Maslov-Mosca benchmark tof\_3) & \texttt{main}\,[f] & -- & $2$ & \cite{amy2014polynomial}\\
\texttt{amm\_\allowbreak{}tof\_\allowbreak{}4} & 4-controlled Toffoli, clean ancillas (Amy-Maslov-Mosca benchmark tof\_4) & \texttt{main}\,[f] & -- & $2$ & \cite{amy2014polynomial}\\
\texttt{amm\_\allowbreak{}tof\_\allowbreak{}5} & 5-controlled Toffoli, clean ancillas (Amy-Maslov-Mosca benchmark tof\_5) & \texttt{main}\,[f] & -- & $2$ & \cite{amy2014polynomial}\\
\texttt{amm\_\allowbreak{}vbe\_\allowbreak{}adder\_\allowbreak{}3} & VBE ripple-carry adder (Amy-Maslov-Mosca benchmark vbe\_adder\_3) & \texttt{main}\,[f] & -- & $2$ & \cite{amy2014polynomial}\\
\texttt{bell\_\allowbreak{}basis\_\allowbreak{}change} & Bell basis change & \texttt{main}\,[fR] & -- & $d\ge 1$ & \cite{nielsen2010quantum}\\
\texttt{bell\_\allowbreak{}register} & Register of Bell pairs & \texttt{main}\,[b] & $n\ge 0$ (pairs) & $d\ge 1$ & \cite{nielsen2010quantum}\\
\texttt{bell\_\allowbreak{}stabilizers} & Bell state stabilizers & \texttt{main}\,[fR] & -- & $d\ge 1$ & \cite{gottesman1999fault}\\
\texttt{bell\_\allowbreak{}state} & Bell state & \texttt{main}\,[fR] & -- & $d\ge 1$ & \cite{nielsen2010quantum}\\
\texttt{bridge\_\allowbreak{}gate} & Bridge gate & \texttt{main}\,[f] & -- & $d\ge 1$ & \cite{itoko2020optimization}\\
\texttt{ccz\_\allowbreak{}decomposition} & CCZ Clifford+T decomposition & \texttt{main}\,[fR] & -- & $2$ & \cite{nielsen2010quantum}\\
\texttt{ccz\_\allowbreak{}t\_\allowbreak{}depth\_\allowbreak{}one} & CCZ with T-depth one & \texttt{main}\,[f] & -- & $2$ & \cite{selinger2013quantum}\\
\texttt{clifford\_\allowbreak{}relation} & Clifford relation (SH)\^{}3 & \texttt{main}\,[f] & -- & $2$ & \cite{gottesman1998heisenberg}\\
\texttt{cluster\_\allowbreak{}stabilizer} & Cluster state stabilizer & \texttt{main}\,[bR] & $n\ge 0$ (qudits after the first two) & $d\ge 1$ & \cite{raussendorf2001one}\\
\texttt{cluster\_\allowbreak{}state} & Linear cluster state & \texttt{main}\,[b] & $n\ge 0$ (qudits after the first) & $d\ge 1$ & \cite{raussendorf2001one}\\
\texttt{cluster\_\allowbreak{}state\_\allowbreak{}2d} & Two-dimensional cluster state & \texttt{main}\,[bR] & $n\ge 1$ (grid side) & $d\ge 1$ & \cite{raussendorf2001one}\\
\texttt{cnot\_\allowbreak{}from\_\allowbreak{}ms} & CNOT from the Molmer-Sorensen gate & \texttt{main}\,[fR] & -- & $2$ & \cite{maslov2017basic}\\
\texttt{cnot\_\allowbreak{}reversal} & Controlled-add reversal & \texttt{main}\,[f] & -- & $d\ge 1$ & \cite{nielsen2010quantum}\\
\texttt{complete\_\allowbreak{}graph\_\allowbreak{}state} & Complete graph state & \texttt{main}\,[b] & $n\ge 0$ (qudits after the first) & $d\ge 1$ & \cite{hein2004multiparty}\\
\texttt{controlled\_\allowbreak{}phase\_\allowbreak{}decomposition} & Controlled phase decomposition & \texttt{cs}\,[f], \texttt{ct}\,[f] & -- & $2$ & \cite{barenco1995elementary}\\
\texttt{cuccaro\_\allowbreak{}adder} & Cuccaro ripple-carry adder & \texttt{main}\,[u] & $n\ge 1$ (bits per register) & $2$ & \cite{cuccaro2004new}\\
\texttt{cx\_\allowbreak{}from\_\allowbreak{}cz} & CX from CZ and Fourier & \texttt{main}\,[fR] & -- & $d\ge 1$ & \cite{gottesman1999fault}\\
\texttt{cx\_\allowbreak{}inverse} & Controlled add and its inverse & \texttt{main}\,[bR] & $n\ge 0$ (qudit pairs) & $d\ge 1$ & \cite{gottesman1999fault}\\
\texttt{cyclic\_\allowbreak{}shift} & Cyclic shift by SWAP chain & \texttt{main}\,[bR] & $n\ge 0$ (wires after the first) & $d\ge 1$ & \cite{nielsen2010quantum}\\
\texttt{cz\_\allowbreak{}symmetry} & Controlled-phase symmetry & \texttt{main}\,[fR] & -- & $d\ge 1$ & \cite{gottesman1999fault}\\
\texttt{draper\_\allowbreak{}adder} & Draper QFT adder & \texttt{main}\,[uL] & $n\ge 1$ (digits per register) & $d\ge 2$ & \cite{draper2000addition}\\
\texttt{fermionic\_\allowbreak{}swap} & Fermionic swap & \texttt{main}\,[f] & -- & $2$ & \cite{bravyi2002fermionic}\\
\texttt{fourier\_\allowbreak{}basis\_\allowbreak{}state} & Fourier basis state & \texttt{main}\,[b] & $n\ge 0$ (qudits) & $d\ge 1$ & \cite{nielsen2010quantum}\\
\texttt{fourier\_\allowbreak{}conjugation} & Fourier conjugation of the shift & \texttt{main}\,[bR] & $n\ge 0$ (qudits) & $d\ge 1$ & \cite{gottesman1999fault}\\
\texttt{fourier\_\allowbreak{}constant\_\allowbreak{}adder} & Fourier constant adder & \texttt{main}\,[uL] & $n\ge 1$ (qudits) & $d\ge 2$ & \cite{draper2000addition}\\
\texttt{fourier\_\allowbreak{}order\_\allowbreak{}four} & Fourier gate of order four & \texttt{main}\,[bR] & $n\ge 0$ (qudits) & $d\ge 1$ & \cite{wang2023qufinitezxw}\\
\texttt{fredkin\_\allowbreak{}decomposition} & Fredkin gate decomposition & \texttt{main}\,[fR] & -- & $2$ & \cite{fredkin1982conservative}\\
\texttt{full\_\allowbreak{}adder} & Reversible full adder & \texttt{main}\,[f] & -- & $2$ & \cite{vedral1996quantum}\\
\texttt{ghz} & GHZ state & \texttt{main}\,[bR]; restricted: n >= 2 & $n\ge 0$ (qudits after the first) & $d\ge 1$ & \cite{ghz1989beyond}\\
\texttt{ghz\_\allowbreak{}disentangle} & GHZ disentangling & \texttt{main}\,[bR] & $n\ge 0$ (qudits after the first) & $d\ge 1$ & \cite{ghz1989beyond}\\
\texttt{ghz\_\allowbreak{}log\_\allowbreak{}depth} & GHZ state, logarithmic depth & \texttt{main}\,[b] & $0\le n\le 3$ (tree levels (2\^{}m qudits)) & $d\ge 1$ & \cite{cruz2019efficient}\\
\texttt{ghz\_\allowbreak{}x\_\allowbreak{}basis} & GHZ state in the Fourier basis & \texttt{main}\,[b] & $n\ge 0$ (qudits after the first) & $d\ge 1$ & \cite{ghz1989beyond}\\
\texttt{gidney\_\allowbreak{}and} & Temporary logical AND & \texttt{main}\,[f] & -- & $2$ & \cite{gidney2018halving}\\
\texttt{givens\_\allowbreak{}rotation} & Givens rotation & \texttt{main}\,[f] & -- & $2$ & \cite{kivlichan2018quantum}\\
\texttt{hadamard\_\allowbreak{}euler} & Euler decomposition of H & \texttt{main}\,[f] & -- & $2$ & \cite{coecke2008interacting}\\
\texttt{hypergraph\_\allowbreak{}state} & Hypergraph state & \texttt{main}\,[f] & -- & $2$ & \cite{rossi2013quantum}\\
\texttt{incrementer} & Quantum incrementer & \texttt{main}\,[b] & $n\ge 1$ (bits) & $2$ & \cite{barenco1995elementary}\\
\texttt{iswap\_\allowbreak{}decomposition} & iSWAP decomposition & \texttt{main}\,[fR] & -- & $2$ & \cite{schuch2003natural}\\
\texttt{local\_\allowbreak{}complementation} & Local complementation & \texttt{main}\,[f] & -- & $2$ & \cite{vandennest2004graphical}\\
\texttt{margolus\_\allowbreak{}gate} & Margolus gate & \texttt{main}\,[fR] & -- & $2$ & \cite{barenco1995elementary}\\
\texttt{measurement\_\allowbreak{}uncomputation} & Measurement-based uncomputation & \texttt{main}\,[f] & -- & $2$ & \cite{gidney2018halving}\\
\texttt{modular\_\allowbreak{}multiplication} & Modular multiplication by 7 mod 15 & \texttt{main}\,[f] & -- & $2$ & \cite{markov2012constant}\\
\texttt{multi\_\allowbreak{}controlled\_\allowbreak{}toffoli} & Multi-controlled Toffoli, V-chain & \texttt{main}\,[uR] & $n\ge 2$ (controls) & $2$ & \cite{barenco1995elementary}\\
\texttt{pauli\_\allowbreak{}gadget} & Pauli exponential synthesis & \texttt{main}\,[f] & -- & $2$ & \cite{cowtan2020phase}\\
\texttt{pauli\_\allowbreak{}propagation} & Pauli propagation through CX & \texttt{x}\,[fR], \texttt{z}\,[fR] & -- & $d\ge 1$ & \cite{gottesman1999fault}\\
\texttt{peres\_\allowbreak{}gate} & Peres gate & \texttt{main}\,[fR] & -- & $2$ & \cite{peres1985reversible}\\
\texttt{phase\_\allowbreak{}kickback} & Phase kickback & \texttt{main}\,[f] & -- & $d\ge 1$ & \cite{cleve1998quantum}\\
\texttt{phase\_\allowbreak{}polynomial\_\allowbreak{}synthesis} & Phase-polynomial synthesis & \texttt{main}\,[f] & -- & $2$ & \cite{amy2014polynomial}\\
\texttt{plus\_\allowbreak{}state} & Uniform superposition & \texttt{main}\,[b] & $n\ge 0$ (qudits) & $d\ge 1$ & \cite{nielsen2010quantum}\\
\texttt{qram} & Quantum RAM lookup & \texttt{main}\,[bR] & $n\ge 1$ (address bits) & $2$ & \cite{giovannetti2008quantum}\\
\texttt{qudit\_\allowbreak{}swap} & SWAP from controlled adds & \texttt{main}\,[fR] & -- & $d\ge 1$ & \cite{gottesman1999fault}\\
\texttt{ring\_\allowbreak{}cluster\_\allowbreak{}state} & Ring cluster state & \texttt{main}\,[b] & $n\ge 1$ (qudits after the first two) & $d\ge 1$ & \cite{raussendorf2001one}\\
\texttt{sleator\_\allowbreak{}weinfurter\_\allowbreak{}toffoli} & Sleator-Weinfurter Toffoli & \texttt{main}\,[f] & -- & $2$ & \cite{barenco1995elementary}\\
\texttt{sqrt\_\allowbreak{}swap} & Square root of SWAP & \texttt{main}\,[f] & -- & $2$ & \cite{loss1998quantum}\\
\texttt{star\_\allowbreak{}graph\_\allowbreak{}ghz} & GHZ from the star graph & \texttt{main}\,[b] & $n\ge 0$ (leaves) & $d\ge 1$ & \cite{hein2004multiparty}\\
\texttt{star\_\allowbreak{}graph\_\allowbreak{}state} & Star graph state & \texttt{main}\,[b] & $n\ge 0$ (leaves) & $d\ge 1$ & \cite{hein2004multiparty}\\
\texttt{swap\_\allowbreak{}from\_\allowbreak{}cz} & SWAP from controlled phases & \texttt{main}\,[f] & -- & $d\ge 1$ & \cite{nielsen2010quantum}\\
\texttt{toffoli\_\allowbreak{}decomposition} & Toffoli Clifford+T decomposition & \texttt{main}\,[fR] & -- & $2$ & \cite{nielsen2010quantum}\\
\texttt{toffoli\_\allowbreak{}temporary\_\allowbreak{}and} & Toffoli via temporary AND & \texttt{main}\,[f] & -- & $2$ & \cite{jones2013low}\\
\texttt{w\_\allowbreak{}state} & W state preparation & \texttt{main}\,[b] & $0\le n\le 3$ (tree levels (2\^{}m qubits)) & $2$ & \cite{dur2000three}\\
\texttt{weyl\_\allowbreak{}commutation} & Weyl commutation relation & \texttt{main}\,[f] & -- & $d\ge 1$ & \cite{weyl1927quantenmechanik}\\
\midrule
\end{xltabular}
}

\section{The false set}\label{sec:bench-false}

\module{benchmarks.breaks} builds the false set $F$: for each $S$ instance, at most one broken
copy whose circuit no longer equals the specification. A \term{break kind} lists candidate edits
(\term{variants}) of the circuit, each with probe points where it should differ and a record of
where it stays equal. Generation applies the first variant whose counterexample search finds an
entry that differs; a kind with no such variant does not apply, and the assignment moves to the
next kind. No break edits the specification.

\subsection{Edits}

\begin{signature}
\py!Edit(kind: str, label: str, params: dict)!\\
\py!  apply(c, inst, nv, dv) -> Circuit!  \qquad \py!apply_spec(spec, inst, nv, dv, ctx) -> spec!
\end{signature}

\py!apply! dispatches on \code{kind} through the registry \py!EDITS!; \py!KINDS! is its key tuple
in the order of \cref{tab:bench-kinds}. Helpers: the outputs of a circuit are its wires with no
\py!Post! or \py!Discard!, its inputs those with no \py!Init! (chain bodies flattened); the
\emph{phase wires} are the open, unmeasured outputs, the single wires first in reverse order,
then the registers. An edit names a wire by name (\code{wire}), or, on an unrolled equation, by
\code{role} (\code{phase}, \code{inputs}, \code{all}) and \code{slot} in that list. It names a
gate by \code{index}, or on an unrolled equation by type, wires, occurrence \code{occ} and total
\code{of} (\py!NotApplicable! when the count of such gates is not \code{of}).

\begin{xltabular}{\textwidth}{@{}>{\ttfamily}p{0.27\textwidth} X@{}}
\caption{Break kinds, their edits, and their variants (in registry order).}\label{tab:bench-kinds}\\
\toprule
\normalfont Kind & Edit and variants\\
\midrule\endfirsthead
\toprule
\normalfont Kind & Edit and variants\\
\midrule\endhead
\bottomrule\endfoot
beyond\_oracle & Appends \py!ZPhase! (or \py!XPhase!) with a rule on a phase wire. General $d$:
  rule $k(d-1)(d-2)(d-3)(d-4)/d$, probes at the first $n$ and two $d\ge5$, labels
  \code{Z|X phase on q}. Boxed with size parameter: rule
  $k\,n(n-1)(n-2)(n-3)/(24d)$, probes at the first $n\ge4$ and two $d\ge2$, labels \code{Z|X n-phase on q}. The rule is an integer at every $k$ inside the oracle grid $d\le4$, $n\le3$;
  equal-where \texttt{\{grid: oracle\}}.\\
beyond\_N & Unrolled: $X$ on output $-i$ when $n > N$ ($N$ = \code{false_set.beyond_N}),
  for every $i\ge1$ whose output $-i$ is a single wire of the subject's dimension at the probe
  $n$ and the next 3 domain $n$; label \code{X on output -i beyond N}. Boxed: a gadget
  multiplying by $-1$ when $n \equiv N+1 \pmod{N+2}$ (\cref{alg:bench-gadget}), label
  \code{sign at n = N + 1 mod N + 2}. Probe at the first $n\ge N+1$; equal-where
  \texttt{\{n\_max: N\}}. Not for fixed equations.\\
prime\_only & Appends a rule $(k^d-k)/d$ (Fermat) on each phase wire, $Z$ or $X$; general $d$
  only; probes at two composite $d\ge4$; equal-where \texttt{\{d\_prime: true, also: [561, 1105, 1729]\}}.\\
qubit\_only & Appends $F^2$ (the antipode) on each phase wire; general $d$ only; probes at two
  $d\ge3$; equal-where \texttt{\{d\_max: 2\}}.\\
coincident\_dims & When the circuit has wires of at least two dimension expressions: a rule
  $k(D_q - D_o)/D_q$ on each phase wire, $D_o$ the first other dimension; equal-where
  \texttt{\{dims\_equal: [...]\}}.\\
scalar\_small\_one & Multiplies the scalar by $1+(d-1)(d-2)(d-3)(d-4)$ (general $d$; probe at
  one $d\ge5$) or by $1+n(n-1)(n-2)(n-3)$ (boxed; probe at one $n\ge4$).\\
global\_phase & Multiplies the scalar by $e^{2\pi i/7}$; equal-where
  \texttt{\{up\_to\_global\_phase: true\}}.\\
tiny\_phase & Appends \py!ZPhase! $\{1: 1/d^k\}$, $k$ = \code{false_set.tiny_k}, on each phase
  wire; parameter \code{d_min = 2}; a variant whose counterexample is proportional (a
  scalar-only break) ends the kind.\\
leg\_order & Appends \py!Swap! of the first two output registers of one count (label
  \code{swap a, b per copy}); for a non-boxed equation with at least two single outputs, of the
  last two outputs (\code{swap last two outputs}).\\
box\_boundary & Boxed only: multiplies the scalar by $1 + 0^{(n - n_0)}$, $n_0$ the domain's
  $n_{\min}$; probes at $n_0$; equal-where \texttt{\{n\_gt: n0\}}.\\
state\_not\_map & Prepends $Z$ on each unmeasured input; equal-where \texttt{\{input\_zero: q\}}.\\
single\_entry & When every output is a single wire: appends a \py!PhaseFn! over the unmeasured
  single outputs with phase $1/2$ where every output is 1 (support that point only); parameter
  \code{d_min = 2}. A product dimension above 18 decimal digits (the engine loader's numeral
  bound) raises \py!TableUnsupported!, a subclass of both \py!NotApplicable! and
  \py!Unsupported!.\\
wrong\_phase & Per top-level gate $X$, $Z$, \py!ZPhase!, \py!XPhase!, CX, CZ: $a\mapsto a+1$;
  rule $\mapsto$ rule $+k/(3d)$; the least entry key (1 when none) $\mapsto$ entry $+1/3$;
  $s\mapsto -s$.\\
missing\_gate & Drops one top-level gate other than Init, Post, Discard, Measure, Swap, last
  first; refused when the interface changes.\\
extra\_gate & Appends $X$ on each phase wire.\\
swapped\_control\_target & Exchanges control and target of one top-level CX.\\
off\_by\_one\_n & Size parameter and non-fixed form: unrolled, the circuit at $n+1$; boxed,
  every register count $c\mapsto c+1$ and every Chain or Tree whose count is a register count or
  \code{n} likewise.\\
wrong\_dimension & General $d$: per wire touched only by single-wire gates, its dimension
  $D\mapsto D+1$.\\
\end{xltabular}

Unless stated otherwise, a variant probes the \emph{small points}: the first three domain $n$
from $n_{\min}$ (within 40 steps) and the first two domain $d\ge2$ (at most 64), all pairs, $d$
outer; its equal-where record is empty. Candidate $d$ lists stop at $\min(d_{\max}, 64)$;
candidate $n$ lists scan at most 200 values. The probe circuit is built at the first domain $n$
and $d$ (the first domain $d\ge2$, else $\ge1$).

\begin{algorithm}[Boxed \texttt{beyond\_N} gadget]\label{alg:bench-gadget}
With every register count concrete and $n$ known, the edit multiplies the scalar by $-1$ exactly
when $n \bmod (N+2) = N+1$. Otherwise it adds a register \code{_cnt} of count \code{n} (after the
last register of count \code{n}, or last when none) and a single wire \code{_acc} at the end, both
of dimension $M = N+2$, and appends Init$(\_cnt, 1)$, Init$(\_acc, 0)$, CX$(\_cnt, \_acc)$,
\py!ZPhase(_acc, {N+1: 1/2})!, CX$^{-1}(\_cnt,\_acc)$, Post$(\_acc,0)$, Post$(\_cnt,1)$. The
accumulator holds $n \bmod M$ after the fan, so the gadget's value is
\py!gadget_value(params, n)! $=-1$ iff $n\equiv N+1 \pmod M$, else 1.
\end{algorithm}

On an unrolled equation, \py!variants! rewrites each wire variant to (role, slot), and each
gate variant of \code{wrong_phase}, \code{missing_gate}, \code{swapped_control_target} to
(type, wires, occurrence, total), keeping it only when the next 3 domain $n$ have the same total
of such gates; its label becomes \code{[drop |reverse ]NAME OCC of TOTAL on WIRES}.

\subsection{Counterexample search}

\begin{definition}[Counterexample]
A \py!Witness! records the point $(n,d)$, the output and input indices \code{outs}, \code{ins}
(of the doubled diagram for a mixed equation), the two values \code{left} and \code{right} as
\code{[real, imag]}, or an \code{interface} text when the boundaries differ; \code{method}
(\code{entry} or \code{factored}); and \code{proportional} (whether the two tensors differ by
one scalar; \code{None} when no dense check ran). It serializes to a JSON object with exactly
these nine keys.
\end{definition}

Constants: tolerance \py!TOL! $=10^{-8}$; dense limit \py!DENSE_LIMIT! $=2^{18}$ tensor entries.

\begin{algorithm}[\texttt{differing\_entry}]\label{alg:bench-differing}
Input: the two loaded diagrams (circuit, specification) of the broken instance at $(n,d)$.
\begin{enumerate}
\item Instantiate both; when either is mixed, double both (\cref{ch:mixed}).
\item When the boundary dimensions differ, return an interface counterexample
  \code{left [lo]->[li], right [ro]->[ri]} with \code{proportional = False}.
\item When the tensor has at most $2^{18}$ entries: contract both densely; when the largest
  entrywise gap is at most $10^{-8}$ return none, else return the argmax entry with
  \code{proportional} by \py!proportional(a, b)! (both nonzero and $a = \lambda b$ within
  \py!TOL!, $\lambda$ read at $b$'s largest entry).
\item Otherwise evaluate candidate index tuples: for fill values 0, 1, $-1$ (mod each
  dimension), the tuples (fill, fill), (fill, 0), (0, fill); then 24 uniform random tuples from
  the seeded generator; duplicates skipped. A diagram whose largest node arity is at most 12 is
  evaluated by plugging basis effects $X$ states and contracting numerically (\py!plugged!),
  otherwise by symbolic contraction at the bound values. Return the first tuple whose values
  differ by more than \py!TOL!.
\end{enumerate}
\end{algorithm}

The boxed \code{beyond_N} kind uses the \emph{factored} search instead: with $g$ the gadget
value ($|g|^2$ when mixed; $g=1$ means no counterexample), find a nonzero specification entry $e$
among the extrapolated tuples and then the candidates, and return \code{left} $=ge$,
\code{right} $=e$, method \code{factored}, \code{proportional = True}. The extrapolated tuples
take the first nonzero dense entries at the two smallest $n\ge1$ below the target, find a block
of $b>0$ indices that repeats once per copy in both, and repeat it $n$ times.

\py!dense_differs(left, right)! decides whether two loaded diagrams differ: \code{True} on a
boundary mismatch ignoring dimension-1 ports, \code{None} above the dense limit, else whether the
largest entrywise gap exceeds \py!TOL!.

\subsection{Generation}

\begin{algorithm}[Kind order]\label{alg:bench-order}
The allowed kinds are the registry kinds matching a glob of \code{false_set.kinds}, in
registry order. For the $i$-th $S$ instance, with \code{round_robin} the order is the rotation
starting at $i \bmod |\mathit{kinds}|$; with \code{random} it is a shuffle by a generator seeded
with the first 8 bytes (big-endian) of SHA-256(\code{"SEED|EQUATION"}).
\end{algorithm}

\begin{algorithm}[\texttt{generate\_one}]
\begin{enumerate}
\item Seed the candidate generator with the first 8 bytes of SHA-256(\code{"SEED|EQUATION"}).
  Set the deadline \code{equation_seconds} ahead (none when not positive).
\item For each kind in order, try its variants in order (\py!try_kind!): for each probe point,
  run the search under a wall limit of $\min(\mathit{probe\_seconds}, \text{time left})$ (a
  \code{SIGALRM} interval timer); a probe that overruns or raises is skipped (a
  \code{ContractSizeError} or \py!MemoryError! marks the equation as overflowing); passing the
  deadline aborts the equation. A counterexample ends the search, except that a proportional
  \code{tiny_phase} counterexample ends that kind without a record.
\item On success return the $F$ record annotated (\py!annotate!) plus \code{skipped}, the kinds
  tried before. On abort return \texttt{\{equation, kind: null, skipped, reason: "time budget"\}};
  when no kind applies, \texttt{\{equation, kind: null, skipped\}}, plus \code{reason: "oracle size"} for an overflowing equation.
\end{enumerate}
\end{algorithm}

An $F$ record (\py!FalseEquation.to_dict!) has the keys \code{equation} (the $S$ id),
\code{kind}, \code{label}, \code{params}, \code{witness}, \code{equal_where}, and after
annotation \code{differs_at} and \code{skipped}.

\py!generate(cfg, path, jobs)! returns the records of \code{path} when the file holds a readable
false set; otherwise it generates every record in a process pool of \code{jobs} processes (chunk
size 1), writes \texttt{\{"fingerprint", "records"\}} (JSON, indent 1) to the temporary file
\code{false_set.<task>.<pid>.tmp} (task from \code{SLURM_ARRAY_TASK_ID}, else \code{local}), and
renames it over \code{path}. \py!load(cfg, path)! returns the records whenever the file exists,
parses, and holds a list under \code{records}; it does not compare fingerprints, so a false set
is regenerated by removing \code{path}.

\begin{definition}[False-set fingerprint]
The first 16 hexadecimal digits of SHA-256 over: the JSON (sorted keys) of the \code{subjects}
and \code{false_set} sections; then, for every \file{*.py} of the catalogue directory (sorted,
helpers included), its name and bytes; then the bytes of \file{ir.py}, \file{translate.py},
\file{breaks.py}, \file{subject.py}. It is written into the file for reference only.
\end{definition}

\subsection{The differs table}

\py!table_grid(inst)! lists the domain points with $d\le 4$ (only the fixed $d$ for a fixed
subject) and $n\le4$ ($n\le12$ at a fixed $d$), $d$ outer. \py!differs_table(broken, seconds)!
returns rows \code{[n, d, differs]} at the grid points the dense check settles, each within
\code{seconds}; the line of a $d$ stops at its first unsettled $n$. \py!annotate(rec, insts, seconds)! adds \code{differs_at}. When the counterexample is not a doubling interface mismatch
(\cref{def:bench-doubling}), its own point is appended as differing when absent; when it is one,
the counterexample is replaced by a dense \py!differing_entry! at the first differing table
point (kept when none succeeds). \py!tabled(records, seconds, jobs)! annotates, in a process
pool, every record with a kind and no \code{differs_at}; an unknown equation or a raising check
yields an empty table. \py!annotate_file(src, dst, cfg, jobs)! writes an annotated copy.

\py!oracle_grid(inst)! lists the points with $d\le4$ and $n\le3$, $n$ outer.

\subsection{Record validation}

\py!mismatch(rec, inst)! returns \code{None} when the record's (label, params) is among today's
variants of its kind on the $S$ instance (computed with the record's own $N$ and $k$, cached), or
when a positional record still names the same thing (\py!_positional_ok!: a fixed
\code{swap last two outputs} whose $i, j$ are the last two output positions; or an unrolled gate
index that is a raw variant of the probe circuit with the same type, wires, and occurrence at the
next 3 domain $n$). Otherwise it returns \code{KIND edit 'LABEL' PARAMS is not a variant of the FORM form of ID}. \py!false_instance(rec, insts)! returns the $F$ instance.

\section{Where an \texorpdfstring{$F$}{F} pair differs}\label{sec:bench-differs}

\module{benchmarks.differs} decides, by exact integer arithmetic, whether the $F$ pair of a
record differs at a point. A pair that is equal at a point measures nothing there; the schedule
avoids such points and the classifier marks runs at them.

\begin{definition}[Doubling mismatch]\label{def:bench-doubling}
An interface text \code{left [lo]->[li], right [ro]->[ri]} is a doubling mismatch when the
multiset of one side's dimensions is the other's taken twice, on both outputs and inputs.
\end{definition}

\begin{algorithm}[\texttt{equal\_at(false, n, d)}]\label{alg:bench-equal-at}
Returns \code{True} (equal), \code{False} (differs), or \code{None} (only a doubling
counterexample speaks).
\begin{enumerate}
\item No record (an $S$ equation): \code{False}. A record with a kind and no \code{differs_at}
  key raises \py!ValueError!. A symbolic point ($d$ \code{None}): \code{False}.
\item A tabled point: the negation of its table value.
\item \code{beyond_oracle}, \code{prime_only}, \code{coincident_dims} with a rule: equal iff the
  phase is an integer at every $k<d$: for the Fermat rule, iff $k^d\equiv k\pmod d$ for every
  $k<d$; otherwise the rule must be $k\cdot r(n,d)$ with $r$ free of $k$ (else not equal), and
  equal iff the numerator of $r$ is divisible by its nonzero denominator.
\item \code{scalar_small_one}, \code{box_boundary} with a factor: equal iff factor $-1$ is 0 at
  $(n,d)$ ($n$ read as 0 when absent).
\item \code{beyond_N}: unrolled, equal iff $n\le N$ or $d=1$; boxed, equal iff
  $n\bmod(N+2)\ne N+1$.
\item \code{qubit_only}: equal iff $d\le2$.
\item The table's lines (\cref{alg:bench-table-equal}) show it equal: \code{True}.
\item A gate kind (\code{wrong_phase}, \code{missing_gate}, \code{swapped_control_target}):
  \code{None} when the edited gate differs between the unrolled probe $n$ and the next 3 domain
  $n$; \code{True} when the gate is trivial at the point (\cref{alg:bench-gate-trivial}).
\item \code{off_by_one_n}: \code{False}.
\item A doubling counterexample of method \code{entry}: \code{None}. Otherwise \code{False}.
\end{enumerate}
\end{algorithm}

\begin{algorithm}[Table lines]\label{alg:bench-table-equal}
With \code{TAIL} $=3$ and \code{PERIODS} $=3$, the table shows the pair equal at an untabled
$(n,d)$ when: no tabled point differs (a no-op edit); or $d=1$ and every tabled $d=1$ point is
equal; or $n=0$ and every tabled $n=0$ point is equal; or $n$ exceeds the largest tabled $n$ and
either the last 3 tabled $n$ are equal at every tabled $d$, or the tabled $n$ line of $d$ has a
least period $p$ (at least 3 checked pairs $v(m) = v(m+p)$, $p\le(\max-\min)/2$) and the tabled
$m\equiv n \pmod p$ is equal; or $d$ exceeds the largest tabled $d$ and the last 3 tabled $d$ are
equal at every tabled $n$.
\end{algorithm}

\begin{algorithm}[Trivial edited gate]\label{alg:bench-gate-trivial}
The edited gate is looked up in the $S$ instance's circuit built at the probe $n$ and symbolic
$d$ (by \code{index}, or by type, wires, and occurrence), with the dimension expressions of its
wires; the lookup is cached per (equation, parameters), and a failed lookup counts as not
trivial. The dimensions and amounts are evaluated at $(n,d)$ ($n$ read as 0 when absent; a
non-rational value means not trivial). With sizes $D_i$:
\begin{itemize}
\item \code{wrong_phase}: an $X$ or $Z$ on a one-level wire; a \py!ZPhase!/\py!XPhase! on a
  one-level wire, or, for an entry edit, whose edited level (the least entry key, else 1) is at
  least $D$; a CX or CZ on two wires of equal size $D$ with $D \mid 2s$.
\item \code{missing_gate}: an $X^a$ or $Z^a$ with $a$ an integer divisible by $D$; a
  \py!ZPhase!/\py!XPhase! whose phase is an integer at every level $k<D$ (rule) or at every entry
  key below $D$ (entries), only for $D\le4096$; a CX, CZ (equal sizes), or $F$ with $D=1$.
\item \code{swapped_control_target}: a CX with some wire of one level.
\end{itemize}
\end{algorithm}

\py!differs(false, n, d)! is \code{equal_at(...) is False}.

\section{The point schedule}\label{sec:bench-points}

\module{benchmarks.points} chooses the points of one unit. A point is $(n, d)$, $n$ \code{None}
for a subject without a size parameter. Phase 1 is a fixed growing ladder ended by its first
counted failure; phase 2 draws log-uniform points from the region not dominated by a failed point
until $K$ successes or a cap on failures.

\subsection{Axes}

\py!axes(inst, cfg, false=None, dims=None) -> Axes! sets the ranges
\[
  n\in\bigl[\max(n^{\mathrm{cfg}}_{\min}, n^{\mathrm{dom}}_{\min}),\,
    \min(n^{\mathrm{cfg}}_{\max}, n^{\mathrm{dom}}_{\max})\bigr],\qquad
  d\in\bigl[\max(d^{\mathrm{cfg}}_{\min}, d^{\mathrm{dom}}_{\min}),\,
    \min(d^{\mathrm{cfg}}_{\max}, d^{\mathrm{dom}}_{\max})\bigr],
\]
the configured bounds from \code{points}; there is no $n$ axis without a size parameter and the
$d$ range is the fixed $d$ of a fixed subject. With a tool's dimensions, $d$ narrows to
$[\min \mathit{dims}, \max\mathit{dims}]$. The axes keep the domain tests and a \code{differs}
predicate from \cref{sec:bench-differs}; \py!contains(p)! checks the ranges, the domain, and \code{differs}.

\py!stated(inst, cfg, n, d)! maps a published proof's stated size into the domain: $n$ dropped
without an $n$ axis, \code{None} read as the least domain $n$ of the range; \code{None} when the
point lies outside.

\subsection{Phase 1: the ladder}

\begin{algorithm}[\texttt{ladder}]\label{alg:bench-ladder}
For $k = 1,\dots,K$ with $u = k/K$:
\begin{enumerate}
\item $n = \lfloor (n_{\max}+1)^u\rfloor - 1$ and $d=\lfloor (d_{\max}+1)^u\rfloor$, with
  $n_{\max}$, $d_{\max}$ from \code{points}, clamped to the axes' upper bounds and moved up to
  the least domain value (falling back to the least domain value of the whole range).
\item Skip the step when no $d$ (or no required $n$) exists.
\item Make the point distinct (\py!_distinct!): keep it when unused and differing; else raise
  $n$ through the next domain values (at most 4096); else raise $d$ likewise; else the first
  unused point of the domain enumeration (at most 10\,000).
\end{enumerate}
\end{algorithm}

\subsection{Phase 2: the fill}

\py!seed(base, equation, tool)! is the first 8 bytes (big-endian) of
SHA-256(\code{"BASE|EQUATION|TOOL"}). \py!Phase2.of(ax, cfg, equation, tool)! creates a generator
with that seed and \code{points.seed}. \py!draw()! draws $u$ uniform and sets
$n=\lfloor(n_{\max}+1)^u\rfloor-1$, and $d$ fixed when the $d$ range is one value, else
$\lfloor(d_{\max}+1)^{u'}\rfloor$; it redraws until \py!contains!, stopping after 200\,000 draws.
\py!next(failed, used)! returns the next drawn point that is unused and not \term{dominated}
(at least as large in $n$ and $d$ as some failed point, a missing $n$ comparing as larger), or
\code{None} when the draws run out.

\py!enumerate_domain(ax, limit, differing=True)! yields domain points, $n$ outer, at most
\code{limit}; with \code{differing}, only differing points among the first $4096\cdot
\mathit{limit}$ domain points. \py!domain_size(ax, cap)! counts domain points up to \code{cap}.

\section{Stages and timing spans}\label{sec:bench-stages}

\module{benchmarks.stages} runs the ArchytasZX stages of one run inside the worker process.
The stage names (\py!STAGES!) are \code{decide}, \code{certificate}, \code{lean_export},
\code{lean_check}; the adapter adds a preceding \code{generate} stage (script generation) and is
specified in \cref{ch:bench-tools}.

\begin{signature}
\py!timed(fn, limit: float | None) -> dict!
\end{signature}
\noindent runs \code{fn} under a CPU limit: it installs a \code{SIGPROF} handler, arms the
profiling interval timer \code{ITIMER_PROF} at \code{limit} seconds, resets the engine trace,
and records \py!process_time! and \py!perf_counter!. On expiry the handler stores a span snapshot
(when tracing is on), marks expiry, re-arms the timer at 0.5\,s, and raises
\py!StageTimeout!, so code that swallows one exception is interrupted again. The result is
\code{fn}'s fields plus \code{status}: \code{ok}, \code{timeout}, \code{oom} (\py!MemoryError!,
\code{detail = "MemoryError"}), or \code{error} (\code{detail} the last line of the exception,
at most 500 characters). A stage that expired is \code{timeout} whatever it returned. It adds
\code{cpu}, \code{wall}, and \code{peak_mb} (the process's peak resident memory in MiB, one
decimal); with tracing on, \code{spans} (the span tree), and on a timeout
\code{spans_at_timeout} and \code{open_span} (the open span path joined by \code{/}).

\paragraph{Spans.} Engine tracing is on when \code{ARCHYTASZX_TRACE=1} at import. A span node is
a JSON object with \code{name}, \code{cpu}, \code{wall} (seconds, six decimals), \code{calls},
optional \code{label}, and optional \code{children}; a snapshot marks unfinished nodes with
\code{"open": true} and their partial times. The open path names each open node as
\code{name[label]}.

\paragraph{Sessions.} \py!new_run(base="")! creates a fresh shell session with concrete axioms on
(as the shell starts) and an empty outcome list; \py!execute(run, text)! runs one command line;
\py!logical_lines(script)! splits a script into the shell's logical lines, dropping lines that are
empty or comment only. \py!diagrams(script)! executes the lines before the first \code{decide} or
\code{prove} and returns the current diagram and the name \code{rhs} (a failing line raises
\py!RuntimeError!).

\paragraph{Stage functions.}
\begin{itemize}
\item \py!decide_stage(run, script)!: executes every line; a failing non-decide line raises
  \py!RuntimeError("LINE: TEXT")!; no decide line raises \py!RuntimeError("the script has no decide line")!. It returns \py!verdict_fields! of the decide outcome: \code{verdict}
  (\code{equal}, \code{unequal}, \code{equal_uncertified}, \code{unequal_uncertified},
  \code{unknown}), \code{certified}, \code{method}, \code{inner_method}, \code{reason} (clipped by
  \py!clip_reason!, \cref{sec:bench-outcomes}), \code{rungs} (\py!stop_codes! of the decision's stops, else \py!rung_codes!
  of the reason), \code{derivation_steps} (the certificates' step total, else the induction
  length), \code{samples}, \code{assumptions} (a count), and \code{family_holds}.
\item \py!certificate_stage(run)!: \code{certificate kernel}; \code{kernel} is \code{rejected}
  (with \code{detail}, the last 300 characters), \code{theorem} (with \code{inferences}, the
  theorem's size; \code{inference_kinds}, a count per inference kind over the premise tree; and
  \code{rules}, the number of axioms used), \code{symbolic}, \code{refutation}, or the payload's
  class name. \py!certificate_bytes(run)! is the UTF-8 length of \code{certificate show}.
\item \py!lean_export_stage(run, path)!: \code{certificate lean PATH}; returns
  \code{lean_semantics} (the first \code{lean: } line not starting \code{lean: wrote}, at most
  200 characters), \code{lean_bytes}, \code{lean_lines}, and \code{lean_encoding_only}
  (whether the output holds the encoding-level line of \code{certificate lean}); a failure
  or missing file raises.
\item \py!lean_check_stage(path, limit, memory_fraction)!: the Lean runner
  (\cref{sec:bench-lean}); a Lean timeout raises \py!StageTimeout!; a status other than
  \code{accepted} sets \code{status} to \code{oom} or \code{error} and \code{detail}.
\item \py!not_run(reason)!: \texttt{\{status: "not\_run", detail: reason\}}.
\end{itemize}

\section{The runner}\label{sec:bench-runner}

\module{benchmarks.runner} takes every (equation, tool) unit through its symbolic run, the
phase 1 ladder, and the phase 2 fill. Every run is a fresh worker process under a memory cap,
and every finished run is appended to the log before the next starts. Units run
\code{platform.jobs} at a time in a thread pool; the runs of one unit go one after another,
since each step depends on the outcomes before it.

\subsection{Start-up}

\begin{algorithm}[\texttt{run(cfg, shard, only)}]\label{alg:bench-run}
\begin{enumerate}
\item Create the results directory; write the configuration to
  \file{configs/<hash>.json} (indent 1).
\item Preflight the enabled tools (\py!tools.preflight!, \cref{ch:bench-tools}); on any problem
  print \code{benchmarks: preflight failed, not starting:} and one indented line per problem, and
  exit with status 2.
\item Record the environment (\cref{tab:bench-environment}) with the calibration and write it to
  \file{environment/<host>-<job>-<task>-<unix time>.json}; the host is the short host name, the
  job \code{SLURM_ARRAY_JOB_ID} (else \code{local}), the task \code{SLURM_ARRAY_TASK_ID} (else
  \code{0}).
\item Obtain the false set (\cref{alg:bench-falseset}) when enabled.
\item Build the context: the log index of this configuration hash, the machine record, the
  calibration median and tag, the per-worker memory cap, the cgroup flag (whether
  \py!memcap.mechanism()! is \code{systemd}), the git
  commit and dirty flag, the Lean slots and budget, and the stale units
  (\cref{alg:bench-stale}).
\item Build the units (\cref{alg:bench-units}); keep those whose id hashes into this shard
  ($\text{SHA-256}(\mathit{id})$ as an integer $\bmod N = I$); with \code{--only}, keep those
  matching a glob.
\item Print \code{benchmarks: U units, config H, J workers, M MiB each, S Lean slots of B MiB, R runs already logged}.
\item Install \code{SIGINT} and \code{SIGTERM} handlers that set the stop flag and terminate
  every live worker. Run the units in the pool; after each, print \code{[k/U] UNIT}; an exception
  prints \code{unit UNIT: TYPE: MESSAGE} and moves on.
\item When stopped, print \code{benchmarks: interrupted; finished runs are logged, rerun to resume} and exit with status 130.
\end{enumerate}
\end{algorithm}

\begin{xltabular}{\textwidth}{@{}>{\ttfamily}p{0.2\textwidth} X@{}}
\caption{The environment record.}\label{tab:bench-environment}\\
\toprule
\normalfont Field & Content\\
\midrule\endfirsthead
\bottomrule\endfoot
config, config\_hash & The resolved configuration and its hash.\\
git\_commit, git\_dirty & \code{git rev-parse HEAD}; whether \code{git status --porcelain --untracked-files=no} prints anything.\\
machine & \code{node} (host name), \code{cpu_model} (first \code{model name} of
  \file{/proc/cpuinfo}), \code{governor} (cpu0's scaling governor), \code{mem_total_mb}
  (\cref{sec:bench-lean}), \code{os}, \code{python}, \code{cpus}.\\
software & Installed versions of \code{archytaszx}, \code{sympy}, \code{numpy}, \code{matplotlib}
  (\code{null} when absent).\\
lean\_toolchain, mathlib & The Lean project's \file{lean-toolchain} text and the Mathlib
  \code{rev} of \file{lake-manifest.json}.\\
started & Local time \code{YYYY-MM-DDTHH:MM:SS}.\\
slurm & Every environment variable whose name starts \code{SLURM_}.\\
calibration & \texttt{\{script: "GHZ symbolic decide", cpu: [...], median\}}.\\
\end{xltabular}

\paragraph{Calibration.} The runner executes a pinned script three times, each in a fresh
Python process, and records the decide stage's CPU seconds; the median is the middle sorted
value (\code{None} with no time). The script decides, at symbolic $n$ and $d$, the
\bbox{}ed Bell-pair circuit (an \azx!X! state through $F$, copied by \azx!Z(1 -> 2)! onto an
\azx!X(2 -> 1)! fed by a second \azx!X! state, the second state and the sum node in
\azx!repeat n!, scalar \code{1/sqrt(d)}) against \azx!g = Z(0 -> 2, d)! with
\azx!repeat n: g.out1! and scalar $d^{-1/2}$. Point rows divide stage CPU medians by this
median (\cref{sec:bench-aggregate}).

\begin{algorithm}[False set of a run]\label{alg:bench-falseset}
Shard 0 takes an exclusive \code{flock} on \file{false_set.lock} and calls \py!generate!, then
\py!tabled! with \code{probe_seconds} per point. Every other shard polls \py!load! every 30
seconds (printing \code{benchmarks: shard I waits for PATH} every 600 seconds) until the file
holds this configuration's set, then calls \py!tabled!.
\end{algorithm}

\begin{algorithm}[Units]\label{alg:bench-units}
For every $S$ instance and every enabled tool known to the adapter registry, a unit
\py!Unit(equation, tool, None, restricted)!. When the false set is enabled, for every record with
a kind: when \py!mismatch! returns a reason, print \code{benchmarks: F record refused: REASON}
and skip it; otherwise one unit per tool on the $F$ instance. A unit's id is
\code{EQUATION[!]|TOOL}.
\end{algorithm}

\begin{algorithm}[Stale units]\label{alg:bench-stale}
A logged unit record is stale, and its unit runs again, when it lacks \code{differs_checked}
and either (a) its tool is check-once and its \code{K} is not 1, or it has no
\code{check_phase} and the adapter's \code{symbolic_covered} refuses the subject; or (b) it is
an $F$ unit with a logged concrete run at a point where the pair does not differ.
\end{algorithm}

\subsection{The unit schedule}

\py!UnitRun(unit, inst, ctx)! computes the source digest of the instance once. A run's version
is \code{restricted} for a restricted instance, else \code{symbolic} for the point \code{None},
else \code{concrete}. \py!run(point, repeat, trace, phase)! returns the logged record of its key
(\cref{def:bench-key}) when one exists, reclassified when it has no \code{cause}; a logged
check-once refusal (unsupported code \code{published_once}) runs again. Otherwise it launches a
worker, builds the run record (\cref{tab:bench-run-record}), classifies it
(\cref{sec:bench-outcomes}), appends it, and indexes it.

\begin{algorithm}[\texttt{UnitRun.execute}]\label{alg:bench-execute}
\begin{enumerate}
\item Skip the unit when its unit key \code{HASH|EQ@SOURCE|unit|TOOL} is logged and not stale.
\item Ask the adapter \py!supports(inst, None)!. A check-once adapter
  (\code{CHECK_ONCE}) with support runs \cref{alg:bench-check-once}. A refusal writes the unit
  record with \code{unsupported}, its code, and \code{successes = 0}.
\item With \code{points.symbolic} and a symbolic version, run the symbolic point (phase
  \code{symbolic}); on success run its repeats.
\item Build the axes with the $F$ record and the adapter's \code{DIMENSIONS}; count the domain up
  to $K$. With adapter dimensions and an empty domain, write the unit as unsupported with the
  reason \code{qubits only (d = 2)} (dimensions \code{(2,)}) or \code{dimensions [...] only}.
\item When the domain has fewer than $K$ points, run every one (phase \code{all}), repeats after
  successes, and write the unit with \code{domain} and \code{region_exhausted = False}.
\item Phase 1: for each ladder point run it (phase \code{ladder}); an uncounted failure is
  skipped; a counted failure is recorded and ends the ladder; a success runs its repeats.
\item Phase 2, while successes are below $K$: draw the next point (phase \code{random}); none
  left sets \code{region_exhausted}; a success runs its repeats; an uncounted failure is
  skipped; a counted failure joins the failed list, and after \code{phase2_extra_failures} of
  them \code{region_exhausted} is set and the phase ends.
\item Write the unit record.
\end{enumerate}
\end{algorithm}

\begin{algorithm}[Check-once units]\label{alg:bench-check-once}
A check-once tool (published proofs) makes one check with $K=1$: the symbolic run when
\code{points.symbolic}, the instance has a symbolic version, and the adapter's
\code{symbolic_covered} accepts the subject; otherwise the adapter's \py!stated(inst, cfg)!
point (phase \code{stated}); when that is \code{None}, the unit is unsupported with \code{the published proof's stated size lies outside the domain}. The unit record carries \code{K: 1} and
\code{check_phase}.
\end{algorithm}

\begin{algorithm}[Repeats]\label{alg:bench-repeats}
After a successful repeat 0 at a point: in \code{fixed} mode run repeats $1,\dots,\mathit{runs}-1$.
In \code{adaptive} mode collect the main CPU of each success (the \code{decide} stage, else
\code{check}, else the first stage) and run repeats $r=1,2,\dots$ while $r<\mathit{max\_runs}$,
stopping before a repeat once $r\ge\max(\mathit{min\_runs},1)$, $r\ge$ \py!stats.min_runs!, and
the sample is precise (\cref{sec:bench-stats}). With tracing on and a positive
\code{overhead_sample}, a point whose key \code{"EQ|TOOL|POINT"} hashes below the share (first 8
bytes of SHA-256 over $2^{64}$) also runs repeat 0 with tracing off.
\end{algorithm}

\subsection{Records}

\begin{xltabular}{\textwidth}{@{}>{\ttfamily}p{0.2\textwidth} X@{}}
\caption{Fields the runner writes on every run record.}\label{tab:bench-run-record}\\
\toprule
\normalfont Field & Content\\
\midrule\endfirsthead
\toprule
\normalfont Field & Content\\
\midrule\endhead
\bottomrule\endfoot
key & The run key (\cref{def:bench-key}).\\
config & The configuration hash.\\
equation & The instance id (with \code{!} for $F$).\\
source & The source digest.\\
truth & True for $S$.\\
kind & The break kind, or \code{null}.\\
subject, category & Of the subject.\\
version & \code{symbolic}, \code{concrete}, \code{restricted}.\\
tool & The tool name.\\
n, d & The point (\code{null} when symbolic).\\
repeat, trace & The repeat index; whether tracing was on.\\
phase & \code{symbolic}, \code{all}, \code{ladder}, \code{random}, \code{stated}.\\
started, finished & Unix times around the launch.\\
machine, calibration, git & The machine record; \texttt{\{median, tag\}}; \texttt{\{commit, dirty\}}.\\
memory\_cap\_mb & The worker memory cap.\\
(worker fields) & The worker's result, or the kill record (\cref{sec:bench-launch}), plus
  \code{process_wall}.\\
outcome, stage, cause, cause\_message & From the classifier; also \code{raw_outcome},
  \code{tags}, \code{point_invalid}, \code{wrong} where it sets them.\\
\end{xltabular}

A unit record has \code{key}, \code{config}, \code{unit: true}, \code{equation}, \code{source},
\code{truth}, \code{subject}, \code{category}, \code{tool}, \code{kind}, \code{restricted},
\code{K}, \code{differs_checked} (true for $F$ units), and the schedule fields:
\code{successes}, \code{failures}, \code{skipped}, \code{succeeded} (lists of rows
\texttt{\{point, outcome, stage, cause, cause\_message, key[, point\_invalid][, phase]\}}, phase 1
or 2 for ladder and fill rows), \code{symbolic} (the symbolic row), \code{region_exhausted},
\code{domain} (small domains), \code{unsupported} and \code{unsupported_code}, and for check-once
units \code{K: 1} and \code{check_phase}.

\subsection{Worker launch and memory caps}\label{sec:bench-launch}

\py!memory_cap_mb(cfg, jobs)! is $\lfloor \mathit{memory\_fraction}\cdot A / \max(1,
\mathit{jobs})\rfloor$ MiB, $A$ the \code{MemAvailable} of \file{/proc/meminfo}, at most the Slurm
allocation.

\begin{algorithm}[\texttt{launch(job, ctx)}]
\begin{enumerate}
\item The wall limit is $\mathit{timeout\_seconds}\cdot(s+1+[\mathit{lean\_check}])\cdot
  \mathit{wall\_factor}$, $s$ the number of enabled stages among decide, certificate,
  lean\_export, lean\_check (the Lean check counts twice: its light and full attempts,
  \cref{sec:bench-lean}).
\item Choose the worker's cap mechanism (\py!worker_memcap!): \code{systemd} when the context's
  cgroup flag is set (\py!memcap.mechanism()! is \code{systemd} at start-up); else
  \code{cgroup_v2} when the context's \code{memcap} override, or \py!memcap.mechanism()!, is
  \code{cgroup_v2} and a child cgroup \code{azx-worker-*} of the cap can be made
  (\cref{sec:bench-memcap}); else \code{srun} when \code{BENCH_WORKER_STEP} is \code{srun}
  and \code{srun} is usable (inside a Slurm job, \code{srun} on \code{PATH}); else
  \code{rlimit_as}. (The \code{srun} test comes before the \code{cgroup_v2} test.)
\item Write the job JSON to a temporary file \code{bench-job-*.json}. The command is
  \code{python -m benchmarks.worker FILE}; under \code{systemd} it runs under
  \code{systemd-run --user --scope --quiet --collect -p MemoryMax=<cap>M -p MemorySwapMax=0};
  under \code{cgroup_v2} the child joins its cgroup before exec; under \code{srun} it is wrapped
  by a \py!memcap.Cap! of the cap with mechanism \code{srun} and tag \code{worker}, a Slurm step
  of its own.
\item The environment adds \code{OMP_NUM_THREADS=1}, \code{OPENBLAS_NUM_THREADS=1},
  \code{PYTHONHASHSEED=0}, the repository on \code{PYTHONPATH}, \code{BENCH_MEMORY_MB} (the cap
  under \code{rlimit_as}, else empty), \code{BENCH_MEMCAP_WORKER} (the mechanism),
  \code{BENCH_SPAN_FILE} (the job file plus \code{.spans}), and, with a Lean budget,
  \code{BENCH_LEAN_SLOTS} and \code{BENCH_LEAN_BUDGET_MB}.
\item Start the process in a new session; a timer kills it at the wall limit (\py!terminate!:
  \code{SIGTERM} to the process group, \code{SIGKILL} after 5 seconds). After exit, a
  \code{cgroup_v2} cap reports an oom kill when its cgroup's \code{oom_kill} counter rose, an
  \code{srun} cap when the exit is nonzero and standard error matches a memory-kill message
  (\py!Cap.oom!); the cgroup and the \code{srun} stats file are removed.
\item Read standard output line by line: \code{@@STAGE NAME} sets the current stage,
  \code{@@RESULT JSON} the result, \code{@@OOMKILL K} the scope's oom-kill rise, \code{@@SPANS JSON} the spans (source \code{sigterm}).
\item Without a result or spans, read the spans from the side file (source \code{snapshot});
  remove the job file and side files.
\item When stopping and no result arrived, raise \py!InterruptedError!. A result receives
  \code{process_wall} and \code{memcap} (the mechanism) and is returned.
\item Otherwise classify the kill: \code{timeout} by the stage timer; \code{oom} by exit code
  $-9$ or 137 (\code{SIGKILL}); \code{oom} by a positive oom-kill rise reported by the worker or
  by the \code{cgroup_v2} or \code{srun} cap (\code{cgroup oom_kill});
  else \code{error} (\code{exit CODE}). An error with exit 143 or $-15$ under cgroups becomes
  \code{oom} when the user journal shows an oom-kill failure of scope \code{run-p<pid>-*} since
  the start (polled every 0.5\,s for up to 3\,s). Return \texttt{\{outcome, killed\_in, exit, kill, detail (the last 400 characters of standard error), process\_wall, memcap[, spans, spans\_source, open\_span]\}}.
\end{enumerate}
\end{algorithm}

\section{The worker}\label{sec:bench-worker}

\code{python -m benchmarks.worker JOB_JSON_FILE} executes one run. The job holds
\code{equation}, \code{false}, \code{tool}, \code{point}, \code{repeat}, \code{trace}, and the
resolved \code{config}.

\begin{algorithm}[Worker main]
\begin{enumerate}
\item Reserve 32\,MiB in a byte array (released before encoding the result); lift the integer
  string-conversion limit; without \code{BENCH_MEMORY_MB}, read the scope's \code{oom_kill}
  counter from cgroup v2 \file{memory.events} (\code{None} outside a \file{.scope}); install the
  \code{SIGTERM} handler.
\item Set \code{ARCHYTASZX_TRACE} to 1 or 0 from the job. For a traced ArchytasZX run with
  \code{BENCH_SPAN_FILE}, start a daemon thread that writes \texttt{\{"spans": ...\}} to the side
  file every 30 seconds (atomically through \file{.tmp}).
\item With \code{BENCH_MEMORY_MB}, set the soft \code{RLIMIT_AS} to that many MiB.
\item Rebuild the instance from the configuration and the job (the $F$ instance when
  \code{false} is set). For a rival, print \code{@@STAGE prepare} and arm a CPU timer of
  \code{stages.prepare_timeout_seconds} (default \code{timeout_seconds}) over its preparation;
  the adapter pauses and resumes it around the rival's own call (\py!end_prepare!,
  \py!resume_prepare!).
\item Call \py!adapter.run(inst, job)!. An expired preparation yields \texttt{\{outcome: timeout, killed\_in: prepare, limit: prepare\_cpu, reason: "prepare: CPU limit"\}}; an exception
  yields \py!failure(tool, exc)!.
\item Print \code{@@RESULT JSON}. When encoding raises \py!MemoryError!, retry without span trees
  (\code{result_encode: "MemoryError"}), then with only the fields whose JSON is under 4096
  characters.
\end{enumerate}
\end{algorithm}

\py!failure(tool, exc)!: a translator or rival \py!Unsupported! gives the adapter's unsupported
record; \py!MemoryError! gives \texttt{\{outcome: oom, killed\_in: prepare\}}; any other exception
in ArchytasZX gives \code{error}; in a rival, \code{harness_error} with \code{reason "spec bound: ..."} and \code{harness_cause: spec_bound} for a refused specification, else \code{reason "prepare: ..."}. The \code{SIGTERM} handler prints \code{@@OOMKILL} with the counter's rise,
then \code{@@SPANS} with the snapshot, kills the running Lean checks, and exits with
$128+\mathit{signal}$. \py!mark(stage)! prints \code{@@STAGE stage}.

\section{The results log}\label{sec:bench-log}

\begin{definition}[Run key]\label{def:bench-key}
\py!key(cfg_hash, equation, version, tool, point, repeat, trace)! is
\code{HASH|EQUATION|VERSION|TOOL|n=N|d=D|r=R|t=T}, with \code{n=sym|d=sym} for the symbolic
point and \code{T} 0 or 1. The runner passes \code{EQUATION@SOURCE} as the equation, so an edited
subject never reuses old runs.
\end{definition}

\module{benchmarks.log} keeps one file per worker thread,
\file{results/raw/<host>-<job>-<task>-w<k>.jsonl} (the job from \code{SLURM_ARRAY_JOB_ID}, else
\code{SLURM_JOB_ID}, else \code{local}; the task from \code{SLURM_ARRAY_TASK_ID}, else 0). Each
record is one compact JSON line (separators \code{,} and \code{:}, non-JSON values as text),
appended under a lock, flushed, and \code{fsync}ed; no line is ever edited.
\py!lines(results)! reads every file in sorted order, stops a file at a line without a final
newline, skips malformed lines, and yields objects with a \code{key}. \py!index(results, cfg_hash)! keeps the first record of every key (of that hash when given), except that a unit
record with \code{differs_checked} replaces one without.

\section{Outcomes, stages, and causes}\label{sec:bench-outcomes}

\module{benchmarks.outcomes} gives every run record an \code{outcome}, a \code{stage}, a
\code{cause}, and a \code{cause_message}. \py!classify! reads only a record's fields, so it
applies to new and logged records alike.

\subsection{Outcomes and stages}

\begin{xltabular}{\textwidth}{@{}>{\ttfamily}p{0.22\textwidth} X@{}}
\caption{Run outcomes (\texttt{OUTCOMES}).}\label{tab:bench-outcome-list}\\
\toprule
\normalfont Outcome & Meaning\\
\midrule\endfirsthead
\bottomrule\endfoot
ok & The expected verdict, certified (equal on $S$, unequal on $F$).\\
uncertified & The expected verdict without a certificate.\\
wrong & The opposite verdict, certified.\\
wrong (uncertified) & The opposite verdict, uncertified.\\
unknown & No verdict.\\
unsupported & The tool or translator cannot express the instance.\\
timeout & A CPU or wall limit.\\
oom & A memory kill or allocation failure.\\
error & A failure of the tool.\\
harness\_error & A failure of the harness itself; not scored.\\
invalid\_point & A verdict at a point where the $F$ pair is equal; not scored.\\
\end{xltabular}

The verdict outcomes are the first five; the wrong outcomes are \code{wrong} and \code{wrong (uncertified)}; \code{harness_error} and \code{invalid_point} are not scored. The stages are
\code{generate}, \code{decide}, \code{certificate}, \code{lean_export}, \code{lean_check} (the
engine stages), \code{prepare}, and \code{rival}. Further constants: the exact tools
\code{archytaszx}, \code{pyzx}, \code{quizx}, \code{cirq}, \code{qiskit}; the up-to-scalar tools
\code{qcec}, \code{dizx}, \code{feynman}, \code{autoq_para}, \code{quantomatic}; the scalar kinds
\code{scalar_small_one}, \code{global_phase}, \code{box_boundary}; the float simulators
\code{qiskit}, \code{cirq} with \py!BREAK_FLOOR! $=10^{-6}$.

\subsection{Cause codes}

\begin{xltabular}{\textwidth}{@{}>{\ttfamily}p{0.35\textwidth} X@{}}
\caption{Cause codes (\texttt{CAUSES}).}\label{tab:bench-causes}\\
\toprule
\normalfont Cause & Meaning\\
\midrule\endfirsthead
\toprule
\normalfont Cause & Meaning\\
\midrule\endhead
\bottomrule\endfoot
verdict.<verdict> & The run's verdict.\\
pair\_equal\_at\_point & The $F$ pair is equal at this point.\\
break\_below\_float\_tolerance & A float simulator measured the break below \py!BREAK_FLOOR!.\\
false\_record\_mismatch & The $F$ record's edit is not a variant of the current form.\\
rival.no\_verdict & The rival returned no verdict.\\
kill.stage\_timer & The runner's wall timer killed the worker.\\
kill.sigkill & The worker died of SIGKILL or exit 137.\\
kill.cgroup\_oom\_kill & The scope's oom\_kill counter rose.\\
kill.exit & The worker died with another exit code.\\
limit.stage\_cpu & A stage used up its CPU limit.\\
limit.generate\_cpu & Script generation used up its CPU limit (timeout).\\
limit.prepare\_cpu & A rival's preparation used up its CPU limit (timeout).\\
limit.rival\_cpu, limit.rival\_wall & The rival used up its CPU or wall limit.\\
memory.error & \py!MemoryError! or a failed allocation in the worker.\\
memory.rival & The rival was killed for memory or raised \py!MemoryError!.\\
memory.lean & Lean was killed for memory.\\
rival.exception & An exception in the rival's library.\\
rival.exit & The rival's binary or driver exited nonzero.\\
engine.load\_refused & A load line of the script failed in the engine.\\
limit.load\_size & The engine's parser refused a load line by a size bound (a power or product
  too large, a port count above the bound).\\
engine.exception & An exception in an ArchytasZX stage.\\
lean.rejected & Lean did not accept the exported file.\\
harness.prepare & An exception while preparing a rival's input.\\
harness.spec\_bound & The engine's loader refused a rival's specification script.\\
harness.driver & An exception in the harness's own driver code.\\
harness.generate & An exception while generating the ArchytasZX script.\\
unsupported.<code> & A code of \cref{tab:bench-unsupported}.\\
unsupported.input\_size & A \py!MemoryError! or memory kill while generating the ArchytasZX
  script outside the engine's input spans.\\
unknown.<rung code> & The last rung tried and its stop code (\cref{tab:bench-rungs}).\\
\end{xltabular}

\begin{xltabular}{\textwidth}{@{}>{\ttfamily}p{0.3\textwidth} >{\ttfamily}X@{}}
\caption{Unsupported codes (\texttt{UNSUPPORTED}): code and the message fragment it matches.}\label{tab:bench-unsupported}\\
\toprule
\normalfont Code & \normalfont Fragment\\
\midrule\endfirsthead
\toprule
\normalfont Code & \normalfont Fragment\\
\midrule\endhead
\bottomrule\endfoot
register\_count\_mismatch & gate across registers of different counts\\
phase\_table\_on\_register & a phase table on a register\\
phase\_table\_dims & a phase table needs literal concrete dimensions\\
phase\_table\_bound & above the engine loader's bound of\\
phase\_table\_irrational & a phase table entry is not a rational number of turns\\
controlled\_mixed\_dims & a controlled gate between wires of different dimensions\\
multi\_controlled\_dims & a multi-controlled gate needs literal concrete dimensions\\
oracle\_dims & an oracle needs literal concrete dimensions\\
dense\_gate\_no\_zxw & a dense reference gate has no ZXW translation\\
permutation\_on\_register & a permutation on a register\\
permutation\_dims & a permutation needs literal concrete dimensions\\
no\_general\_form & the circuit has no general form in n\\
needs\_concrete\_d & the equation needs a concrete dimension\\
symbolic\_point & no symbolic n or d\\
qubits\_only & qubits only\\
odd\_prime\_only & odd prime d only\\
fixed\_dimension & fixed dimension\\
channel & measurement or discard (a channel, not a pure map)\\
wire\_dim\_not\_d & a wire of dimension other than d\\
mixed\_wire\_dims & wires of mixed dimensions\\
gate\_unsupported & "gate "\\
scalar\_modulus & a circuit scalar of modulus other than 1\\
circuit\_scalar & a circuit scalar\\
spec\_mixed & a mixed specification\\
spec\_axes & the specification has more than\\
spec\_too\_large & the specification tensor is too large\\
spec\_generator & "generator "\\
spec\_weighted & a weighted spider\\
spec\_no\_wires & the specification has no wires\\
spec\_zero\_state & the specification is the zero state\\
spec\_isometry & the specification is an isometry, not a unitary\\
spec\_not\_unitary & the specification is not unitary\\
no\_reference & no gate-level reference circuit for this equation\\
gate\_set & not expressible in the gate set\\
tensor\_axes & " tensor axes: above numpy's"\\
einsum\_subscripts & "-qudit gate: above numpy.einsum's"\\
too\_large\_to\_simulate & ": too large to simulate"\\
post\_selected & post-selected wires\\
phase\_not\_pi4 & a phase that is not a multiple of pi/4\\
not\_clifford\_t & " outside Clifford+T"\\
not\_dizx\_clifford & " is not a DiZX Clifford gate"\\
width\_mismatch & circuit and reference differ in width\\
autoq\_amplitude & an amplitude outside sqrt2 powers and eighth turns\\
states\_only & only states (every wire initialized, none post-selected)\\
zero\_state & the zero state\\
no\_public\_release & symbolic n needs AutoQ-Para, which has no public release\\
quantomatic\_bang\_boxes & symbolic n through Quantomatic's own !-boxes is not built\\
mixed\_diagram & a mixed diagram\\
hadamard\_arity & a hadamard vertex of arity other than 2\\
published\_once & published proofs are checked once, at their stated size\\
stated\_outside & the published proof's stated size lies outside the domain\\
proof\_file\_missing & " proof file "\\
not\_installed & not installed\\
qubit\_cap & Requested too many qubits from package\\
input\_size & the input format is intractable at this size\\
\end{xltabular}

\py!unsupported_code(message)! strips a leading \code{[dotted.]Unsupported: } and returns the
code whose fragment occurs in the text, the longest fragment winning; \code{other} when none
matches.

{\small
\begin{xltabular}{\textwidth}{@{}>{\ttfamily}p{0.5\textwidth} >{\ttfamily}X@{}}
\caption{Rung stop codes (\texttt{RUNGS}) and their patterns, matched at the start of a note.}\label{tab:bench-rungs}\\
\toprule
\normalfont Code & \normalfont Pattern (regular expression)\\
\midrule\endfirsthead
\toprule
\normalfont Code & \normalfont Pattern\\
\midrule\endhead
\bottomrule\endfoot
normal\_form.differ & normal forms differ\\
normal\_form.agree & normal forms agree\\
oracle.refused & \textbackslash d+ sample\textbackslash(s\textbackslash) refused\\
oracle.no\_mismatch & no oracle mismatch in \textbackslash d+ sample\\
oracle.none\_evaluated & no oracle sample evaluated\\
oracle.mismatch & oracle mismatch at\\
saturation.iteration\_limit & saturation: .*stopped by iteration\_limit\\
saturation.enode\_limit & saturation: .*stopped by enode\_limit\\
saturation.saturated & saturation: .*stopped by saturated\\
symbolic.declined.box\_meets\_wire & symbolic: symbolic contraction declined:
  SymbolicContractionUnsupportedError: bang box \textbackslash d+ has symbolic multiplicity
  \textbackslash w+ and meets the rest\\
symbolic.declined.scalar\_budget & symbolic: symbolic contraction declined: ScalarBudget\\
symbolic.declined.other & symbolic: symbolic contraction declined\\
symbolic.indeterminate.replicated\_factors & symbolic: indeterminate: the replicated fact\\
symbolic.indeterminate & symbolic: indeterminate\\
symbolic.entries\_differ & symbolic: entries differ\\
symbolic.differ & symbolic contractions differ\\
induction.base\_not\_proved & induction: base case \textbackslash S+ not proved\\
induction.step\_not\_discharged & induction: step case over \textbackslash w+ not proved: no discharge\\
induction\_reversed.base\_not\_proved & reversed: base case\\
induction\_reversed.step\_not\_discharged & reversed: step c\\
refutation.recheck\_failed & uncertified: (?:no Lean refutation \textbackslash((?:residue|numeric)
  recheck failed|interface recheck failed)\\
refutation.no\_lean & (?:not reported unequal|uncertified): no Lean refutation\\
\end{xltabular}
}

\py!clip_reason(reason)! keeps a decide reason whole up to $4000+12000$ characters; a longer one
becomes its first 4000 characters, \code{ [... K characters cut ...] }, and its last 12000.

\begin{algorithm}[\texttt{rung\_codes(reason)}]
Split the reason at \code{"; "}. For each note: strip one prefix \code{doubled, WORD: } (and
remember it); take the first pattern of \cref{tab:bench-rungs} that matches, else
\code{RUNG.other} when the note starts with \code{word:}, else skip the note. After an
\code{induction} or \code{induction_reversed} note, only notes of the rungs induction,
induction\_reversed, refutation, and reversed are read. A note holding a \py!clip_reason! cut
mark, and the last note of a reason of 399 or 400 characters without one (records logged when
reasons were cut at 400), becomes \code{RUNG.truncated}. A remembered prefix gives \code{doubled.CODE}.
Consecutive duplicates collapse. \py!stop_codes(stops)! maps a decision's (rung, code) stops to
\code{rung.code}, prefixing \code{doubled.} after the stop \code{("doubled", "doubled")} and
writing another doubled stop as \code{doubled.CODE}.
\end{algorithm}

\subsection{The classifier}

\begin{algorithm}[\texttt{classify(rec, false, free, stale)}]\label{alg:bench-classify}
\begin{enumerate}
\item A record with \code{killed_in} and no \code{stages} is a kill record
  (\cref{alg:bench-killed}); its message is prefixed \code{open span PATH; } when a span was open.
\item An \code{unsupported} outcome: stage \code{generate} (ArchytasZX) or \code{prepare}
  (rivals), cause \code{unsupported.CODE} (the logged code, else \py!unsupported_code!).
\item ArchytasZX (\cref{alg:bench-engine-cause}) or a rival (\cref{alg:bench-rival-cause}).
\item With an $F$ record and a concrete $d$: the run's region is its point plus, for each free
  axis, every value in $0,\dots,64$ ($d\ge1$). When the record is stale, a verdict outcome becomes
  \code{invalid_point} with cause \code{false_record_mismatch}; when the pair is provably equal
  (\py!equal_at! true) on the whole region, a verdict becomes \code{invalid_point} with cause
  \code{pair_equal_at_point}; when a float simulator's \code{max_rel_diff} is below
  $10^{-6}$, a verdict becomes \code{invalid_point} with cause
  \code{break_below_float_tolerance}. Stale or equal-region runs are flagged
  \code{point_invalid}.
\item The message is clipped to 300 characters: the tail when it holds a traceback frame, else
  the head. A scalar-only $F$ record adds the tag \code{global_phase} (the counterexample ratio
  has modulus 1 within $10^{-9}$, or by kind \code{global_phase}) or \code{scalar_factor} to
  \code{tags}. A changed outcome keeps the original as \code{raw_outcome}; \code{wrong} is set
  when present before or when the outcome is wrong.
\end{enumerate}
\end{algorithm}

\py!free_axes(inst, rec)! is empty for rivals and for scripts with a point claim; for ArchytasZX
it holds \code{n} unless $n$ is \code{None}, the form is unrolled, or there is no size parameter,
and \code{d} unless the subject has a fixed $d$ or the script was literal. A break is scalar-only
by the counterexample's \code{proportional} flag when set, else by kind (the scalar kinds and the
boxed \code{beyond_N}).

\begin{algorithm}[Kill records]\label{alg:bench-killed}
In order: a preparation CPU limit gives (\code{timeout}, \code{prepare},
\code{limit.prepare_cpu}); a stage-timer kill (or a timeout with no kill source) gives
\code{timeout} with \code{limit.generate_cpu} in generation outside the engine's \code{input.*}
spans, else \code{kill.stage_timer}; a memory sign (cgroup kill, \code{SIGKILL}, journal, exit
$-9$ or 137, an \code{oom} outcome, or an allocation failure text) in generation with an open
span outside the input spans gives (\code{unsupported}, \code{unsupported.input_size});
a cgroup or journal oom-kill gives (\code{oom}, \code{kill.cgroup_oom_kill}); \code{SIGKILL} or
exit $-9$/137 gives (\code{oom}, \code{kill.sigkill}); an \code{oom} outcome or an allocation
failure text gives (\code{oom}, \code{memory.error}); else (\code{error}, \code{kill.exit}). The
stage is \code{killed_in} (\code{None} for \code{start}).
\end{algorithm}

The allocation-failure texts are \code{out of memory}, \code{Heap exhausted}, \code{memory allocation of}, \code{std::bad_alloc}, \code{MemoryError}, \code{OutOfMemoryError}, \code{failed to map segment from shared object}, \code{cannot allocate memory}, \code{Cannot allocate memory},
\code{Memory allocation still failed}, and \code{Unable to allocate}.

\begin{algorithm}[ArchytasZX causes]\label{alg:bench-engine-cause}
A verdict outcome has stage \code{decide}: \code{unknown} has cause \code{unknown.} plus the last
rung code (\code{none} when there is none) and the reason as message; the others
\code{verdict.<verdict>} with the method as message. Otherwise take the first engine stage whose
status is neither \code{ok} nor \code{not_run} (none: \code{error}, \code{engine.exception}). A
timeout gives \code{limit.generate_cpu} in generation outside input spans, else
\code{limit.stage_cpu}. An oom in generation outside input spans gives (\code{unsupported},
\code{unsupported.input_size}); another oom gives \code{memory.lean} in lean\_check, else
\code{memory.error}. A generation error gives (\code{harness_error}, \code{harness.generate}); a
lean\_check error gives \code{lean.rejected}. A detail matching \code{RuntimeError: ... error: MemoryError} or \code{... Unable to allocate}, or the text \code{error return without exception set} with a peak at least 0.9 of the memory cap (3481 MiB when unrecorded), gives (\code{oom},
\code{memory.error}). A decide detail \code{RuntimeError: ...(load|@NAME =): error:} gives
(\code{oom}, \code{limit.load_size}) when it contains \texttt{error: power (...)\^{}(...) is too large}, \code{error: product is too large}, or \code{error: ... leg count above the bound of },
else (\code{error}, \code{engine.load_refused}); otherwise \code{engine.exception}.
\end{algorithm}

\begin{algorithm}[Rival causes]\label{alg:bench-rival-cause}
Without stages: \code{oom} gives (\code{oom}, \code{prepare}, \code{memory.error});
\code{error} or \code{harness_error} gives \code{memory.error} when the detail's head contains
\code{MemoryError} or is \code{ImportError}/\code{OSError} with an allocation text, else
\code{harness.spec_bound} when so marked, else \code{harness.prepare}. With the \code{check}
stage: a verdict gives stage \code{rival} and \code{verdict.<verdict>}; \code{unknown} gives
\code{rival.no_verdict}; a timeout gives \code{limit.stage_cpu} or \code{limit.rival_<cpu|wall>}
by the check's \code{limit} field (\code{wall} when absent and no detail, else \code{cpu}); an
oom or an allocation error gives \code{memory.rival}, as does the null-return text near the cap;
\code{Requested too many qubits from package} gives \code{unsupported.qubit_cap}. Otherwise the
traceback decides (\py!_frame_cause!): a Rust panic or \code{PanicException} gives
\code{rival.exception}; no frame gives \code{rival.exit}; the deepest frame in
\file{/tools/drivers/} being one of the helpers \code{number}, \code{_build}, \code{_identity},
\code{add}, \code{scalar_parts}, \code{norm_identity} gives \code{harness.driver}; the deepest
frame in the drivers calling none of the rival's names (per tool: \code{pyzx}; \code{quizx};
\code{mqt}, \code{QuantumComputation}, \code{verify(}; \code{cirq}; \code{qiskit},
\code{qasm2}, \code{Operator}; \code{pytket}, \code{FullPeepholeOptimise},
\code{circuit_from_qasm_str}; \code{dizx}, \code{clifford_simp}; else the tool name) gives
\code{harness.driver}; otherwise \code{rival.exception}. A harness cause makes the outcome
\code{harness_error}, else \code{error}.
\end{algorithm}

\py!counted(rec)! is true when the outcome is scored and the run is not flagged
\code{point_invalid}. \py!reclassify_record(rec, insts, falses)! looks up the instance (an $F$
id uses its $S$ instance) and the $F$ record, computes the free axes and staleness
(\py!mismatch!), and calls \py!classify!.

\section{Repeat statistics}\label{sec:bench-stats}

\module{benchmarks.stats} gives the exact distribution-free confidence interval of the median
by order statistics~\cite{david2003order} and the adaptive stopping rule of the cluster
profile, a rule of the kind used for rigorous performance measurement~\cite{georges2007statistically}.

\begin{definition}[Median interval]
For a sample of size $m$, let $j$ be the largest integer in $1,\dots,\lfloor m/2\rfloor$ with
\[
  2\,\Pr[\mathrm{Bin}(m, \tfrac12)\le j-1] \;=\; 2^{1-m}\sum_{i=0}^{j-1}\binom{m}{i} \;\le\; 1-c,
\]
the search stopping at the first $j$ that fails. The interval is $[x_{(j)}, x_{(m+1-j)}]$ of the
sorted sample; it covers the population median with probability at least $c$. No $j$ qualifies
for small $m$ (\py!median_ci_ranks! returns \code{None}).
\end{definition}

\py!min_runs(c)! is the least $m$ with an interval; at $c=0.95$ it is 6 ($2\cdot2^{-6} =
0.03125 \le 0.05$). \py!precise(values, c, p)! is true when the half-width is at most $p$ times
the absolute sample median (when the median is 0, when the half-width is 0). \py!summary(values, confidence=None)! returns \code{median}, \code{q1}, \code{q3} (linear-interpolated 25th and 75th
percentiles), \code{min}, \code{max}, \code{n}, the raw \code{values}, and with $c$ the bounds
\code{ci_lo}, \code{ci_hi} (\code{None} without an interval); an empty sample gives
\texttt{\{n: 0, values: []\}}.

\section{Aggregates}\label{sec:bench-aggregate}

\py!aggregate.main(cfg)! writes the aggregates of one configuration hash to
\file{results/aggregate/}, each table as \file{NAME.csv} and \file{NAME.json} (indent 1). A CSV
has the union of the rows' keys in first-seen order as header, and list or dictionary cells as
JSON text; an empty table is an empty file.

\begin{algorithm}[\texttt{aggregate.main}]
\begin{enumerate}
\item Read the log index of the hash; keep the records whose \code{source} equals today's source
  digest of their instance (\py!current!: edited or removed subjects drop out).
\item \py!upgrade! each ArchytasZX record logged before certified verdicts (no
  \code{certified} field): \code{equal} with a rejected kernel becomes
  \code{equal_uncertified}; \code{unknown} whose reason contains \code{not reported unequal:}
  becomes \code{unequal_uncertified}; the outcome is recomputed by the adapter's
  \py!verdict_outcome! and \code{certified}, \code{wrong} are set.
\item Classify (\py!classified!): a run record that the journal shows oom-killed
  (\cref{sec:bench-reclassify}) is reclassified from its raw outcome with
  \code{journal_oom_kill}; a run record without \code{cause} is reclassified; the $F$ records come
  from \file{false_set.json} through \py!tabled!.
\item Drop the runs with outcome \code{invalid_point} or flag \code{point_invalid}.
\item Write the tables of \cref{tab:bench-aggregates}, then print \code{aggregate: R runs, P points, U units, W wrong verdicts, H harness errors, X runs at points where the F pair is equal (invalid_point, point_invalid) excluded -> DIR} and one line \code{ENGINE WRONG VERDICT[ (uncertified)]: EQ VERSION n=N d=D: VERDICT} per wrong ArchytasZX point.
\end{enumerate}
\end{algorithm}

\begin{xltabular}{\textwidth}{@{}>{\ttfamily}p{0.18\textwidth} X@{}}
\caption{Aggregate tables.}\label{tab:bench-aggregates}\\
\toprule
\normalfont Table & Rows and columns\\
\midrule\endfirsthead
\toprule
\normalfont Table & Rows and columns\\
\midrule\endhead
\bottomrule\endfoot
points & One row per (equation, version, tool, $n$, $d$); untraced runs are left out when tracing
  is enabled. Columns from the first repeat: \code{equation}, \code{version}, \code{tool},
  \code{n}, \code{d}, \code{subject}, \code{category}, \code{truth}, \code{kind},
  \code{phase}, \code{outcome}, \code{verdict}, \code{wrong}, \code{runs}, \code{unstable} (first
  run ok but a repeat not ok), \code{noisy} (adaptive and runs $\ge$ \code{max_runs}),
  \code{node}, \code{cpu_model}, \code{calibration}, \code{detail} (200 characters),
  \code{reason}, \code{stage}, \code{cause}. Per stage $s\in$ \{generate, decide, certificate,
  lean\_export, lean\_check, check\} over the successful runs whose stage is ok:
  \code{s_cpu_<summary key>} (with the median interval in adaptive mode), \code{s_wall_median},
  \code{s_peak_mb} (maximum), \code{s_status}, \code{s_cpu_scaled} (median over calibration);
  \code{s_status} from the first run when no run qualifies. \code{run_cpu_*}: the summary of
  total stage CPU, without values. Size fields \code{decide.derivation_steps},
  \code{certificate.inferences}, \code{certificate.certificate_bytes},
  \code{lean_export.lean_lines}, \code{lean_export.lean_bytes}, \code{lean_check.lean_cpu},
  \code{check.nodes_after}, \code{check.proof_lines}, \code{check.proof_bytes}, and
  \code{axiom_class}. Sorted by equation, tool, version, $n$, $d$ (\code{None} as $-1$).\\
units & Per unit record: \code{equation}, \code{tool}, \code{subject}, \code{category},
  \code{truth}, \code{kind}, \code{K}, \code{successes}, \code{K_reached}
  (\code{successes == K}), \code{region_exhausted}, \code{unsupported}, \code{frontier} (the
  failures), \code{n_max_ok} and \code{d_max_ok} (largest ok concrete point).\\
counts & Points per (\code{category}, \code{version}, \code{tool}, \code{set} S/F, \code{kind},
  \code{n_bin}, \code{d_bin}), tket excluded: \code{points} (scored), \code{proved} (S ok),
  \code{refuted} (F ok), \code{wrongly_refuted}, \code{wrongly_accepted}, and one column per
  other outcome. Bins: $n$ \code{none}, \code{0}, \code{1e<k>} with $k=\lfloor\log_{10}n\rfloor$;
  $d$ \code{none}, \code{1}, \code{2}, \texttt{<2\^{}(k-1)+1>-<2\^{}k>} with $k=\lceil\log_2 d\rceil$.\\
causes & Points per (\code{tool}, \code{version}, \code{set}, \code{outcome}, \code{stage},
  \code{cause}, \code{tags}): \code{points}.\\
equations & Per (\code{category}, \code{version}, \code{tool}, \code{set}), tket excluded:
  \code{equations}, \code{settled} (some point ok), \code{wrong} (some point wrong).\\
breakdown & Per (equation, tool, version, $n$, $d$, category, group, name) over traced ok runs:
  the medians \code{cpu} and \code{calls}, and \code{mean} $=$ cpu/calls. Groups: \code{rung}
  (spans \code{decide.*}, \code{saturation.*}), \code{rule} (\code{rewrite.apply} by label),
  \code{match} (\code{rewrite.match}), \code{inference} (\code{kernel.*}), \code{lean}
  (\code{lean.*}, plus \code{lean_check:<kind>} from \code{decl_ms}), \code{other}
  (\code{stage:name}).\\
b3 & Per repeat-0 traced run with a \code{b3} record holding \code{naive}: \code{equation},
  \code{tool}, \code{version}, \code{n}, \code{d}, \code{category}, \code{ratio},
  \code{naive.*}, \code{smallest.*} (\cref{sec:bench-size}).\\
overhead & Per point with a repeat-0 ok run traced and untraced (positive untraced CPU):
  \code{traced_cpu}, \code{untraced_cpu}, \code{ratio} of the decide (else check) stage.\\
calibration & Per environment file: \code{job}, \code{node}, \code{cpu_model}, \code{governor},
  \code{calibration_median}, \code{calibration_runs}, \code{git_commit}, \code{git_dirty}.\\
wrong & The point rows with \code{wrong}.\\
harness\_errors & Per point whose first run is \code{harness_error}: \code{equation},
  \code{version}, \code{tool}, \code{n}, \code{d}, \code{reason}.\\
\end{xltabular}

\section{Plots}\label{sec:bench-plots}

\py!plots.main(cfg)! reads the aggregates and writes PNG (300 dpi) and PDF files to
\code{output.plots} with \tool{Matplotlib}{https://matplotlib.org}~\cite{hunter2007matplotlib}
(backend Agg; white background; TrueType fonts embedded, type 42; font size 8; light grid
\texttt{\#e5e5e5}; solid lines of width 1.5; no legend frame). Tool styles take the Okabe-Ito
hues~\cite{okabe2008color} \texttt{\#0072B2}, \texttt{\#D55E00}, \texttt{\#009E73}, \texttt{\#CC79A7},
\texttt{\#E69F00}, \texttt{\#56B4E9} (yellow and black left out) in the fixed tool order
\code{archytaszx}, \code{pyzx}, \code{quizx}, \code{qcec}, \code{feynman}, \code{dizx},
\code{qiskit}, \code{cirq}, \code{autoq_para}, \code{tket}, \code{quantomatic},
\code{published_proofs}; tool $k$ has hue $k \bmod 6$ and marker \code{o} for $k<6$, \code{s}
otherwise (an unknown tool takes the last position). Every $n$ axis plots $n+1$ on a log scale.

\begin{xltabular}{\textwidth}{@{}>{\ttfamily}p{0.27\textwidth} X@{}}
\caption{Plot files.}\label{tab:bench-plots}\\
\toprule
\normalfont File & Content\\
\midrule\endfirsthead
\bottomrule\endfoot
time\_<stage>\_in\_<axis> & For stages decide, certificate, lean\_export, lean\_check and axes
  $n$, $d$: four panels (categories), S points of concrete versions, per tool one faint line per
  equation (alpha 0.18) and the median over equations in bold with markers; log-log; written
  only when some panel has data.\\
time\_3d & ArchytasZX decide CPU against $\log_{10}(n+1)$ and $\log_{10}d$ per category, 3D
  scatter.\\
frontier & Per tool, the largest ok $(n+1, d)$ of each S unit; log $n$, log$_2$ $d$.\\
lean\_sizes & Lean check CPU and Lean lines against $n$ and $d$ (ArchytasZX, S, concrete), one
  color per category.\\
b3\_sizes & Smallest diagram total against $n$ and $d$, and the median naive-to-smallest ratio
  per category.\\
breakdown\_<group> & For rung, rule, inference: share of traced ArchytasZX CPU per category by
  the up to six largest keys and \code{other}.\\
b2\_kinds & Per tool, F points per break kind as stacked bars refuted, wrongly accepted, unknown,
  other.\\
proved\_bins & S points proved per $n$ bin and per $d$ bin, per tool.\\
\end{xltabular}

The run prints \code{plots -> DIR}.

\section{Reclassification}\label{sec:bench-reclassify}

\code{python -m benchmarks.reclassify RAW_DIR FALSE_SET OUT_JSONL [--journal] [--annotate-seconds S] [--jobs J]} (defaults $S=10$, $J=4$) reclassifies logged runs without
running anything. It refuses an output inside \code{RAW_DIR} and never writes there.

\begin{algorithm}[Reclassification]
\begin{enumerate}
\item Read every record of \file{RAW_DIR/*.jsonl} in file order (non-empty lines).
\item Annotate the $F$ records of \code{FALSE_SET} with \py!tabled! ($S$ seconds per point,
  $J$ processes).
\item For every run record: restore \code{raw_outcome} as the outcome; drop \code{stage},
  \code{cause}, \code{cause_message}, \code{raw_outcome}, \code{point_invalid}; with
  \code{--journal}, set \code{journal_oom_kill} when \py!journal_oom! holds; apply the adapter's
  \code{rederive} hook when it has one; classify by \py!reclassify_record!. Unit records pass
  unchanged, except that the units of check-once tools are recomputed with $K=1$ and
  \code{successes} 1 when some ok run of the unit exists (a symbolic run counts only when the
  adapter's \code{symbolic_covered} accepts the subject).
\item Write every record as one JSON line to \code{OUT_JSONL}; print per tool
  \code{TOOL: OUTCOME=B->A, ...} for every outcome with a nonzero count before or after, then per
  tool \code{TOOL units: K_reached=B->A, K_not_reached=B->A} over supported units.
\end{enumerate}
\end{algorithm}

\py!journal_oom_times(since=None, scope=None)! runs \code{journalctl --user -o short-unix --no-pager [--since @T] -g "Failed with result 'oom-kill'"} (timeout 60\,s) and returns the sorted
Unix times of lines containing \code{.scope:} (and \code{scope} when given); empty when
\code{journalctl} is missing. \py!journal_oom(rec, times)! holds for a kill record with exit 143,
no kill source (or \code{exit 143}), raw outcome \code{error}, and an oom-kill time in
$[\mathit{finished} - 10, \mathit{finished} + 2]$ seconds.

\section{B3 size counts}\label{sec:bench-size}

\module{benchmarks.size} counts the elements of a diagram and finds the smallest diagram of a
decision. For a run whose decide stage returned \code{equal} on an $S$ instance, the ArchytasZX
adapter stores \texttt{b3 = \{naive, smallest, ratio\}}.

\begin{definition}[Element counts]
\py!counts(diagram)! returns \code{node_G} for each generator $G\in\{Z, X, W, F, T, Ti, B, S,
G\}$ (and any other present), \code{nodes}, \code{node_boxes} and \code{port_boxes} (node-scoped
and port-scoped \bboxes), \code{box_depth} (the deepest parent chain), \code{phases} (phase
entries, a rule counting once), \code{weights} (weight entries, a weight rule counting once),
\code{wires}, \code{self_loops} (wires with both ends on one node), \code{inputs},
\code{outputs}, \code{scalar_size} (SymPy operation count plus 1; the scalar text length when
that fails), and
\[
  \mathit{total} = \mathit{nodes} + \mathit{node\_boxes} + \mathit{port\_boxes} + \mathit{phases}
  + \mathit{weights} + \mathit{wires}.
\]
\end{definition}

\py!diagrams_of(decision, left, right)! yields both sides, both normal forms when present, and
every diagram of every certificate's derivation: its initial diagram, the diagram after each
replayed step (stopping at a failing step), its final diagram, and recursively its children's.
\py!smallest! returns the counts with the least (\code{total}, \code{wires}); diagrams whose
counting raises are skipped. \py!b3! returns the naive (circuit) counts, the smallest counts, and
\code{ratio} $=$ naive total over smallest total (\code{None} when zero or absent).

\section{The Lean runner}\label{sec:bench-lean}

\module{benchmarks.lean} checks an exported Lean file with \code{lake env lean FILE} in
\file{archytaszx/lean} (\cref{ch:lean-export}), at most a fixed number of checks at a time per job,
each killed when its anonymous memory passes a budget.

\subsection{Memory plan}

Constants: \py!OLEAN_MB! $=5120$ (the memory-mapped \file{.olean} pages of one check, shared),
\py!LEAN_WORK_MB! $=8192$ (the least budget), \py!WORKER_RESERVE_MB! $=1024$, \py!MAX_SLOTS!
$=5$, the watch period 0.25\,s.

\py!total_memory_mb()! is \code{MemTotal} in MiB, capped by the Slurm allocation
(\code{SLURM_MEM_PER_NODE}, else \code{SLURM_MEM_PER_CPU} $\times$ \code{SLURM_CPUS_PER_TASK})
and by the smallest memory limit of the process's cgroup and its ancestors (v2
\file{memory.max}, v1 \file{memory.limit_in_bytes}; values $\ge 2^{60}$ mean none).

\begin{definition}[Slots and budget]
With pool $P = T - 5120 - 1024\max(1, \mathit{jobs})$ MiB, \py!plan(T, jobs, fixed)! gives
$\mathit{slots} = \max(1, \min(5, \lfloor P/8192\rfloor))$ (or \code{fixed}) and
$\mathit{budget} = \max(8192, \lfloor P/\mathit{slots}\rfloor)$. \py!slots(jobs)! returns
\code{BENCH_LEAN_SLOTS} when set, else the plan with \code{jobs} (else \code{BENCH_JOBS}, else 1);
\py!budget_mb(jobs)! returns \code{BENCH_LEAN_BUDGET_MB} when set, else the plan's budget.
\end{definition}

A slot is an exclusive non-blocking \code{flock} on
\file{archytaszx-bench-lean-<scope>-<k>.lock} in \file{/dev/shm} (when writable, else the
temporary directory), the scope \code{<SLURM_JOB_ID>-<SLURM_ARRAY_TASK_ID>} under Slurm, else
\code{local}; a check retries every 0.5\,s until a slot is free.

\py!cap_mb(fraction, budget)! is $\min(\lfloor \mathit{fraction}\cdot T\rfloor, \lfloor 5120 +
\mathit{budget}\rfloor)$ MiB. Every attempt runs under a \py!memcap.Cap! of that size (tag
\code{lean}, swap allowed) by the chosen mechanism (\cref{sec:bench-memcap}).
\py!cgroups()! is \py!memcap.systemd_ok()!; \py!scoped(cmd, fraction, budget)! wraps a command
in a systemd scope of \py!cap_mb! when systemd is usable (the check itself uses the cap).

\subsection{The check}

\begin{algorithm}[\texttt{check(path, timeout, memory\_fraction)}]\label{alg:bench-lean-check}
\begin{enumerate}
\item Read the source; record \code{file_bytes}, \code{file_lines}, \code{final_theorem} (the
  last \code{theorem}/\code{lemma} name, qualified by the enclosing \code{namespace}s),
  \code{statement_level}, and \code{unproved_declared} (lines starting \code{axiom unproved_}).
\item When a header line (before the last \code{import} line) is exactly \code{import ArchytasZX}, first make the \emph{light} attempt on the source with that line replaced by
  \code{import ArchytasZX.Light}. When it is not \code{accepted}, or the final theorem's axioms
  were not printed, make the \emph{full} attempt on the source as written. The result is the last
  attempt's, with \code{lean_import} (\code{light} or \code{full}) and \code{lean_attempts}, one
  entry per attempt with \code{import} and the attempt's \code{lean_status}, \code{lean_cpu},
  \code{lean_wall}, \code{lean_peak_mb}, \code{lean_peak_anon_mb}, \code{lean_kill},
  \code{lean_detail} where present.
\end{enumerate}
Each attempt:
\begin{enumerate}
\item Write the probe copy \file{<stem>_probe.lean}: after the imports, \code{set_option profiler true} and \code{set_option profiler.threshold 0}; before every declaration (theorem, lemma, def,
  abbrev, instance, example, with its doc comment, attributes, and modifiers),
  \texttt{\#print "@@DECL KIND NAME"} (\code{_} for an anonymous one); at the end \texttt{\#print "@@DECL end \_"} and \texttt{\#print axioms NAME}.
\item Create the cap; take a slot; start \code{lake env lean PROBE} (wrapped by the cap) in a
  new session; before exec, join the cap's cgroup (\code{cgroup_v2}), raise the soft
  \code{RLIMIT_AS} to the hard limit, and set \code{RLIMIT_CPU} at
  $\lfloor\mathit{timeout}\rfloor$ (hard $+5$). A watcher polls every 0.25\,s the
  \code{RssAnon} and the CPU ticks (\code{utime}, \code{stime}, \code{cutime}, \code{cstime} of
  \file{/proc/PID/stat}) summed over the process tree, keeps their maxima, and kills the group
  past the budget; a wall timer at \code{timeout} kills it (both by \py!Cap.kill_group!).
  Afterwards read the cap's \py!usage()!, close the cap, remove its stats file, release the
  slot, and remove the probe.
\item Status, in order: over budget, \code{oom} with \code{lean_kill = budget}; a cap kill
  (\py!Cap.oom!), \code{oom} with \code{lean_kill = memcap.<mechanism>}; timer kill or
  \code{SIGXCPU}, \code{timeout}; \code{SIGKILL}, \code{oom}; exit 0 with no \code{error:} in the
  output, \code{accepted}; an out-of-memory text (\code{out of memory}, \code{std::bad_alloc},
  \code{Cannot allocate memory}), \code{oom}; else \code{rejected} with the first error line (300
  characters).
\item Record \code{lean_budget_mb}, \code{lean_memcap} (the mechanism), \code{lean_cpu} (the
  larger of the rusage user-plus-system time, \code{lean_cpu_rusage}, and the polled tree CPU,
  \code{lean_cpu_tree}), \code{lean_wall}, \code{lean_peak_mb} (rusage), \code{lean_peak_anon_mb}
  (polled), \code{decl_ms}, and \code{axioms} (parsed from \texttt{\#print axioms} when accepted).
  Under \code{srun} the wrapper's \code{cpu} and \code{maxrss_kb} replace the rusage CPU and
  peak, and its \code{peak_anon_mb} and \code{cpu_tree} raise the polled maxima.
\end{enumerate}
The result finally receives \code{axiom_class} and \code{unproved} (the count of
\code{unproved_*} axioms).
\end{algorithm}

\code{decl_ms} lists, per marker, \texttt{\{kind, name, ms, by\}}: the profiler lines
\code{CATEGORY took X(ms|s)} after the marker summed in milliseconds, also per category (the text
before \code{ of }). The axiom class is \code{opaque} when the source lacks \code{import ArchytasZX} or contains \code{namespace ArchytasZX.Opaque}; \code{unknown} with no axiom list;
\code{clean} when only \code{propext}, \code{Classical.choice}, and \code{Quot.sound} occur;
\code{unproved} when every other axiom's last component starts \code{unproved_}; else
\code{other}; for a non-accepted check, the status. The statement level is \code{opaque}
likewise; \code{source} when the final theorem's statement (from its name to \code{:=}) contains
a string quote or a word starting \code{elab}, \code{parse}, \code{Src}, \code{src},
\code{source}; \code{diagram} when it mentions a word starting \code{DEqual}; else
\code{network}. \py!kill_active()! sends \code{SIGKILL} to every running check's process group.
The Lean toolchain is \tool{Lean 4}{https://lean-lang.org}~\cite{demoura2021lean4} with
Mathlib~\cite{mathlib2020}; \code{lake} is found on \code{PATH}, else in
\file{\textasciitilde/.elan/bin}.

\section{Memory caps}\label{sec:bench-memcap}

\module{benchmarks.memcap} caps the memory of one child process tree by a mechanism chosen at
run time.

\begin{xltabular}{\textwidth}{@{}>{\ttfamily}p{0.14\textwidth} X@{}}
\caption{Memory-cap mechanisms (\texttt{ALL}).}\label{tab:bench-memcap}\\
\toprule
\normalfont Mechanism & How the child is capped; when it is usable\\
\midrule\endfirsthead
\bottomrule\endfoot
systemd & \code{systemd-run --user --scope --quiet --collect -p MemoryMax=<cap>M [-p MemorySwapMax=0]}; usable when that command with \code{MemoryMax=64M} runs \code{true}
  (cached).\\
cgroup\_v2 & A child cgroup \code{azx-<tag>-<pid>-<k>} in a delegated cgroup v2 directory with
  \file{memory.max} at the cap, \file{memory.swap.max} 0 (unless swap is allowed), and
  \file{memory.oom.group} 1; the child writes 0 to its \file{cgroup.procs} before exec. The
  delegated directory (cached) is the first of \code{BENCH_CGROUP_PARENT}, the parent of the
  process's own cgroup, and that cgroup itself in which a 64\,MiB probe cgroup can be made and
  \code{true} run inside it.\\
srun & \code{srun --overlap --quiet -n1 -N1 --mem=<cap>M python -I memcap_step.py STATS CMD}
  (the wrapper below, \code{STATS} a temporary file \code{azx-<tag>-*.json}); usable inside a
  Slurm job with \code{srun} on \code{PATH}.\\
watchdog & A thread polling the tree's \code{RssAnon} every 0.25\,s and killing the tree past
  the cap; always usable.\\
\end{xltabular}

\py!mechanism()! returns \code{BENCH_MEMCAP} when it names a usable mechanism; else
\code{srun} when \code{srun} is usable and neither \code{systemd} nor \code{cgroup_v2} is; else
the first usable of \code{systemd}, \code{cgroup_v2}, \code{watchdog}. \py!Cap(limit_mb, mech=None, tag="child", no_swap=True)! creates the cgroup at construction (falling back to \code{watchdog}
when it cannot); \py!command(cmd)! wraps the command (systemd, srun); \py!preexec()! is the
cgroup join; \py!watch(pid, kill)! starts the watchdog; \py!oom(code, text)! reports a cap kill:
the watchdog fired, the cgroup's \code{oom_kill} counter is positive, or under srun a nonzero
exit whose output matches \code{oom[-_ ]kill}, \code{Exceeded (step|job) memory limit}, or
\code{out of memory} (case-insensitive); \py!close()! stops the watchdog and removes the cgroup
(retried for 2\,s). \py!kill_group(pgid)! sends \code{SIGKILL} to the process group; under
\code{srun} it sends \code{SIGTERM} (which \code{srun} passes to the step) and \code{SIGKILL}
3\,s later. \py!usage()! is, under \code{srun}, the wrapper's stats read from \code{STATS}
(\py!None! when absent or not a JSON object); \py!discard()! removes \code{STATS}.
\py!child_cap(tag)! gives a rival child a cap at \code{BENCH_MEMORY_MB}
(none when unset); \py!run(cmd, cap, timeout)! runs a command to its end under a cap in a
session of its own, killing the group on timeout (by \py!kill_group! under a cap), removing
the stats file and re-raising \py!TimeoutExpired!, and returns (the completed process, whether the cap killed it).
\py!worker_mechanism()! is \code{BENCH_MEMCAP_WORKER}, else \code{rlimit_as};
\py!describe()! reports each mechanism's usability and the chosen one.

\paragraph{The in-step wrapper.} \module{benchmarks.memcap_step} uses the standard library only.
\code{python -I memcap_step.py STATS CMD...} starts \code{CMD} as its child and polls every
0.25\,s the \code{RssAnon} and the CPU ticks (\code{utime}, \code{stime}, \code{cutime},
\code{cstime}) summed over the child's process tree, keeping their maxima. When the child ends it
writes to \code{STATS} the JSON object \texttt{\{code, cpu, maxrss\_kb, peak\_anon\_mb,
cpu\_tree\}}, \code{cpu} and \code{maxrss_kb} from the \code{RUSAGE_CHILDREN} of the reaped
child, and exits with the child's code ($128+N$ for a signal $N$).

\section{The ancillary export}\label{sec:bench-anc}

\code{python -m benchmarks anc OUT_DIR} copies into \code{OUT_DIR}: the directory
\file{benchmarks/} without \file{\_\_pycache\_\_}, \file{*.pyc}, \file{.venvs}, \file{results},
\file{plots}, and \file{slurm}; the subdirectories \file{raw}, \file{aggregate},
\file{environment}, and \file{configs} of the results directory and its \file{false_set.json},
under \file{OUT_DIR/results/}; the plot directory as \file{OUT_DIR/plots}; and the file
the files \file{benchmarks-pipeline.tex} and \file{benchmarks-tools.tex} of \file{docs/manual/chapters} into \file{OUT_DIR/docs/}. It prints \code{anc -> DIR}.

\section{The Slurm profile}\label{sec:bench-slurm}

\code{python -m benchmarks slurm [--set k=v ...]} prepares the cluster profile's job array for
\tool{Slurm}{https://slurm.schedmd.com}~\cite{yoo2003slurm}. It sets the profile to
\code{cluster}, checks out the commit to run, writes \file{benchmarks/slurm/run.sbatch}, creates
\file{benchmarks/slurm/logs/}, and prints \code{wrote PATH at SHA; submit with: sbatch benchmarks/slurm/run.sbatch}. Submission is left to the user, from the repository root.

\begin{algorithm}[Commit selection]
\begin{enumerate}
\item The target is \code{slurm.commit} when set, else the head SHA of the latest successful run
  of workflow \file{ci.yml} on branch \code{dev}, from \code{gh run list --branch dev --status success --workflow ci.yml --limit 1 --json headSha -q .[0].headSha} (timeout 60\,s); none
  exits with \code{slurm: no green dev commit found (gh unavailable?); set slurm.commit=SHA}.
\item \code{git fetch -q origin dev}; resolve the target (\code{slurm: commit X not in this checkout} on failure).
\item The checkout is clean when \code{git status --porcelain --untracked-files=no} over every
  tracked path except \file{benchmarks/results/false_set.json} prints nothing (the pathspec
  \py!CLEAN_PATHS!); a dirty checkout exits with \code{slurm: checkout has local changes}.
\item When \code{HEAD} differs, \code{git checkout -q --detach SHA}, verified.
\end{enumerate}
\end{algorithm}

The batch file refuses any partition but \code{scavenger} (\py!ValueError!) and reads:
\begin{lstlisting}[style=shell]
#!/bin/bash
#SBATCH --job-name=archytaszx-bench
#SBATCH --partition=scavenger
#SBATCH --requeue
#SBATCH --array=0-<array - 1>
#SBATCH --ntasks=1
#SBATCH --cpus-per-task=<cpus_per_task>
#SBATCH --hint=nomultithread
#SBATCH --mem-per-cpu=<mem_per_cpu>
#SBATCH --time=<time>
#SBATCH --output=benchmarks/slurm/logs/%A_%a.out
[#SBATCH --account=<account>]
set -euo pipefail
cd "$SLURM_SUBMIT_DIR"
[ "$(git rev-parse HEAD)" = "<sha>" ] || { echo "checkout is not at <sha>" >&2; exit 1; }
[ -z "$(git status --porcelain --untracked-files=no -- ':(top)' ':(top,exclude)benchmarks/results/false_set.json')" ] || { echo "dirty checkout" >&2; exit 1; }
[source <setup>]
source <venv>/bin/activate
export PYTHONPATH="$SLURM_SUBMIT_DIR"
python -m benchmarks run --profile cluster --shard ${SLURM_ARRAY_TASK_ID}/<array> \
    --set platform.jobs=<cpus_per_task> <the --set options given, shell-quoted>
\end{lstlisting}

Each array task thus runs one shard with one worker per core on the scavenger partition, is
requeued on preemption, and resumes from the shared log. The environment variables the suite
reads are \code{SLURM_ARRAY_JOB_ID}, \code{SLURM_JOB_ID}, \code{SLURM_ARRAY_TASK_ID} (log,
environment, and false-set temporary file names, Lean slot scope), \code{SLURM_MEM_PER_NODE},
\code{SLURM_MEM_PER_CPU}, \code{SLURM_CPUS_PER_TASK} (memory caps), \code{SLURM_SUBMIT_DIR}
(the batch file), every \code{SLURM_*} variable (the environment record), and the suite's own
\code{BENCH_MEMORY_MB}, \code{BENCH_MEMCAP}, \code{BENCH_MEMCAP_WORKER}, \code{BENCH_WORKER_STEP},
\code{BENCH_CGROUP_PARENT}, \code{BENCH_SPAN_FILE}, \code{BENCH_LEAN_SLOTS},
\code{BENCH_LEAN_BUDGET_MB}, \code{BENCH_JOBS}, and \code{ARCHYTASZX_TRACE}.
